\documentclass[12pt,a4paper]{report}
\usepackage{elegant-cours}

\title{Recueil des Théorèmes, Propriétés et Démonstrations}
\date{\today}

\begin{document}
\maketitle
\section*{Extraits : 260109 CH10}


\begin{property}
    Soit $(E,*)$ un magma tel que $*$ admet un élément neutre $e$.

    \begin{enumerate}
        \item $e$ est unique.
        \item Si $*$ est associative, et $a\in E$ est symétrisable, alors le symétrique est unique ; i.e. $\exists ! b \in G, a * b = b * a = e$.
    \end{enumerate}
\end{property}

\begin{proof}
    \begin{enumerate}
        \item Supposons qu'il existe deux éléments neutres $e$ et $e'$. Alors, par définition de l'élément neutre, on a $e*e' = e$ et $e*e' = e'$. Donc, $e = e'$.
        \item Soit $a\in G$ symétrisable, et supposons qu'il existe deux symétriques $b$ et $b'$ de $a$. Alors, par définition du symétrique, on a $a*b = e$ et $a*b' = e$. En multipliant à gauche par $b'$ dans la première égalité, on obtient :
        $$b'*(a*b) = b'*e \implies (b'*a)*b = b'.$$
        En utilisant l'associativité, on a :
        $$e*b = b' \implies b = b'.$$
    \end{enumerate}
\end{proof}

\begin{property}
    Soient $(G,*)$ un groupe, $a,b\in G$.

    Alors $(a*b)^{-1}=b^{-1} * a^{-1}$.
\end{property}

\begin{proof}
    $(a*b)*(b^{-1} * a^{-1}) = a*(b*b^{-1})*a^{-1} = a*e*a^{-1} = a*a^{-1} = e$.
    De même, $(b^{-1} * a^{-1})*(a*b) = e$.
    On en déduit que $(a*b)^{-1}=b^{-1} * a^{-1}$.
\end{proof}

\begin{property}[Caractérisation d'un sous-groupe]
    Soient $(G,*)$ un groupe et $H$ une partie de $G$.

    On peut dire que :

    $$H \text{ est un sous-groupe de }(G,*) \iff \begin{cases}H\ne \emptyset\\ \forall x,y\in H, x*y^{-1} \in H\end{cases}$$
\end{property}

\begin{proof}
    \begin{outline}
        \1 \textbf{Démonstration dans le sens direct $\implies$ :} Trivial (mais à refaire).
        \1 \textbf{Démonstration dans le sens réciproque $\impliedby$ :}
            \2 On a $H\ne \emptyset$.
            \2 Donc $\exists a \in H$.
            \2 On a $a*a^{-1} = \boxed{e_G \in H}$.
                \3 $H$ contient l'élément neutre de $G$.
            \2 On a $\forall x \in H, e_G * x^{-1} = x^{-1} \in H$.
            \2 Soient $x,y\in H$. Alors, $y^{-1} \in H$.
            \2 Donc, $x*y = x*(y^{-1})^{-1} \in H$.
                \3 C'est-à-dire que $H$ est stable par $*$.
            \2 $H$ est donc un sous-groupe de $(G,*)$.
    \end{outline}
\end{proof}

\begin{property}[Sous-groupes triviaux]
    Soit $(G,*)$ un groupe.

    Alors, $\{e_G\}$ et $G$ sont des sous-groupes de $(G,*)$, appelés \textbf{sous-groupes triviaux} de $(G,*)$.
\end{property}

\begin{property}
    Tous les sous-groupes de $(\Z,+)$ sont de la forme $n\Z$ avec $n\in\N$.
\end{property}

\begin{proof}
    Pour démontrer ça, on commence par vérifier que $n\Z$ est bien un sous-groupe de $(\Z,+)$ pour tout $n\in\N$. Ensuite, on montre que si $H$ est un sous-groupe quelconque de $(\Z,+)$, alors il existe $n\in\N$ tel que $H=n\Z$.

    \begin{outline}
        \1 $\forall n \in \N$, $n\Z$ est un sous-groupe de $(\Z,+)$.
            \2 $n\Z \subset \Z$ car $\forall k \in \Z, nk \in \Z$.
            \2 $n\Z \ne \emptyset$ car $0 \in n\Z$.
            \2 Soient $a,b \in n\Z$.
                \3 $\exists k,l \in \Z$ tels que $a= nk$ et $b=n l$.
                \3 Donc, $a - b = n(k - l) \in n\Z$.
            \2 Donc, $n\Z$ est un sous-groupe de $(\Z,+)$.
        \1 Soit $H$ un sous-groupe quelconque de $(\Z,+)$, non réduit au singleton $\{0\}$.
            \2 On pose $H^+ = H \cap \N^*$.
            \2 $H^+ \ne \emptyset$ car $H \ne \{0\}$ et $H$ est stable par l'opposé.
            \2 En effet, $H$ est un sous-groupe, alors $-a \in H$ pour tout $a \in H$.
            \2 Donc $a$ et $-a \in H \implies H^+ \ne \emptyset$.
            \2 On a $H^+ \subset \N^*$, donc $H^+$ admet un minimum $n$.
            \3 Nous allons essayer de montrer que $H = n\Z$.
            \2 On a $n \in H$, donc $nk = \begin{cases} \underbrace{n + \dots + n}_{k \text{ fois}}\pour k \ge 0 \\ \underbrace{-(n + \dots + n)}_{k \text{ fois}}\pour k \le 0\end{cases} \in H$.
            \2 Donc $x\in H$, i.e. $n\Z \subset H$.
            \2 Soit $a \in H$.
                \3 Par le théorème de la division euclidienne, il existe un unique couple $(q,r) \in \Z \times \N$ tels que $a = nq + r$ et $0 \le r < n$.
                \3 Or, $nq \in H$ (car $n\Z \subset H$) et $a \in H$, donc $r = a - nq \in H$ (car $H$ est stable par l'opposé et par l'addition).
                \3 Si $r \ne 0$, alors $r \in H^+$ et $r < n$, ce qui contredit la minimalité de $n$.
                \3 Donc, $r = 0$, et ainsi $a = nq \in n\Z$.
            \2 Donc, $H \subset n\Z$.
            \2 On en déduit que $H = n\Z$.
    \end{outline}
\end{proof}

\begin{property}
    Soit $(G,*)$ un groupe.
    \begin{enumerate}
        \item Si $H_{i_{i\in I}}$ est une famille de sous-groupes de $(G,*)$, alors l'intersection $\bigcap_{i\in I} H_i$ est un sous-groupe de $(G,*)$.
        \item $H \cup K$ est un sous-groupe de $(G,*)$ si et seulement si $H \subset K$ ou $K \subset H$.
    \end{enumerate}
\end{property}

\begin{proof}
    \begin{outline}
        \1 \textbf{Démonstration de 1. :}
        \2 $\forall i \in I, H_i \subset G \implies \bigcap_{i\in I} H_i \subset G$.
        \2 $\forall i \in I, H_i \ne \emptyset \implies \bigcap_{i\in I} H_i \ne \emptyset$ (car $e \in H_i$).
        \2 Soient $a,b \in \bigcap_{i\in I} H_i$.
        \2 Alors $\forall i \in I, a,b \in H_i \implies \forall i \in I, a*b^{-1} \in H_i$ (car $H_i$ est un sous-groupe).
        \2 Donc, $a*b^{-1} \in \bigcap_{i\in I} H_i$.
        \2 Donc, $\bigcap_{i\in I} H_i$ est un sous-groupe de $(G,*)$.
        \1 \textbf{Démonstration de 2. :}
        \2 Dans le sens indirect ($\impliedby$), la démonstration est triviale.
        \2 Dans le sens direct ($\implies$), on procède par absurde. On admet que $H \not\subset K$ et $K \not\subset H$.
        \2 Donc $\exists h \in H, h \notin K$ et $\exists k \in K, k \notin H$.
        \2 On a $h,k\in H\cup K \implies h*k \in H \cup K$ (car $H$ et $K$ sont des sous-groupes).
        \2 Si $h*k \in H$, alors $\exists h_1 \in H \tq h*k = h_1$.
        \3 Donc $k=h^{-1}*h_1 \in H$ (car $H$ est stable par $*$ et par l'inverse). (On compose à gauche !)
        \3 Ce qui est absurde.
        \2 Si $h*k \in K$, alors $\exists k_1 \in K \tq h*k = k_1$.
        \3 Donc $h = k_1 * k^{-1} \in K$ (car $K$ est stable par $*$ et par l'inverse). (On compose à droite !)
        \3 Ce qui est absurde.
        \2 On en déduit que $H \subset K$ ou $K \subset H$.
    \end{outline}
\end{proof}

\begin{property}[Groupe des applications à valeurs dans un groupe]
    Soit $(G,*)$ un groupe et $D$ un ensemble non-vide.

    On note $G^D$ l'ensemble des applications de $D$ vers $G$.

    Alors $(G^D, \star)$ est un groupe, où $\star$ est définie par :

    $$\forall f,g \in G^D, \quad f \star g : \begin{cases}D \to G\\x \mapsto f(x) * g(x)\end{cases}$$
\end{property}

\begin{proof}
    \begin{outline}
        \1 \textbf{Démonstration que $\star$ est une LCI dans $G^D$ :}
        \2 On a $\forall f,g \in G^D, (f \star g)(x)=f(x) * g(x) \in G$ (car $*$ est une LCI dans $G$).
        \2 Donc $f \star g \in G^D$.
        \2 D'où $\star$ est une LCI dans $G^D$.
        \1 \textbf{Démonstration que $\star$ est associative :}
        \2 Soient $f,g,h \in G^D,x\in D$.
        \2 On a $\parens{(f \star g) \star h}(x) = (f \star g)(x) * h(x) = (f(x) * g(x)) * h(x)$.
        \2 De même, on a $\parens{f \star (g \star h)}(x) = f(x) * (g \star h)(x) = f(x) * (g(x) * h(x))$.
        \2 Or, $*$ est associative dans $G$, donc $(f(x) * g(x)) * h(x) = f(x) * (g(x) * h(x))$.
        \2 Donc, $\parens{(f \star g) \star h}(x) = \parens{f \star (g \star h)}(x)$.
        \2 Donc, $(f \star g) \star h = f \star (g \star h)$.
        \2 Donc, $\star$ est associative.
        \1 \textbf{Démonstration de l'existence d'un élément neutre :}
        \2 On pose $e : \begin{cases}D \to G\\x \mapsto e_G\end{cases}$, où $e_G$ est l'élément neutre de $G$.
        \2 Soit $f \in G^D, x \in D$.
        \2 On a $(f \star e)(x) = f(x) * e_G = f(x)$, car $f(x) \in G$.
        \2 De même, on a $(e \star f)(x) = e_G * f(x) = f(x)$, car $f(x) \in G$.
        \2 Donc, $f \star e = f$ et $e \star f = f$.
        \2 Donc, $e$ est l'élément neutre de $\star$ dans $G^D$.
        \1 \textbf{Démonstration que tout élément de $G^D$ est symétrisable pour $\star$ :}
        \2 Soit $f \in G^D$.
        \2 On pose $f' : \begin{cases}D \to G\\x \mapsto f(x)^{-1}\end{cases} \in G^D$, où $f(x)^{-1}$ est le symétrique de $f(x)$ pour $*$ dans $G$.
        \2 On a $\forall x \in D, (f \star f')(x) = f(x) * f(x)^{-1} = e_G = e(x)$.
        \2 De même, on a $\forall x \in D, (f' \star f)(x) = f(x)^{-1} * f(x) = e_G = e(x)$.
        \2 Donc, $f \star f' = e$ et $f' \star f = e$.
        \2 Donc, $f$ est symétrisable pour $\star$.
        \2 Donc, tout élément de $G^D$ est symétrisable pour $\star$.
        \1 On en déduit que $(G^D, \star)$ est un groupe.
    \end{outline}
\end{proof}

\begin{property}
    $\forall x,y \in E, x\mathcal{R}y \iff \cl{x} = \cl{y}$.

    i.e. deux éléments sont en relation si et seulement si leurs classes d'équivalence sont égales.
\end{property}

\begin{proof}
    \begin{itemize}
        \item \textbf{Démonstration dans le sens réciproque $\impliedby$ :} On a $x\in\cl{x} = \cl{y}$, donc $x\in\cl{y}$, i.e. $x\mathcal{R}y$.
        \item \textbf{Démonstration dans le sens direct $\implies$ :} Soit $z \in \cl{x}$. Alors, $z \mathcal{R} x$ et $x \mathcal{R} y$, donc $z \mathcal{R} y$ (par transitivité). Donc, $z \in \cl{y}$. Ainsi, $\cl{x} \subset \cl{y}$. De même, on montre que $\cl{y} \subset \cl{x}$. Donc, $\cl{x} = \cl{y}$.
    \end{itemize}
\end{proof}

\begin{property}
    Si $\mathcal{R}$ est une relation d'équivalence sur un ensemble $E$, alors les classes d'équivalence forment \textbf{une partition} de $E$.
    
    i.e. : Elles sont non-vides, disjointes deux à deux, et leur réunion est égale à $E$.
\end{property}

\begin{proof}
    \begin{itemize}
        \item $\forall x \in E, x\in \cl x \implies \cl{x} \ne \emptyset$.
        \item Soient $x,y \in E$ tels que $\cl{x} \ne \cl{y}$.\\Soit $z \in \cl{x} \cap \cl{y}$. Alors, $z \mathcal{R} x$ et $z \mathcal{R} y$, donc $x \mathcal{R} y$ (par symétrie et transitivité). Donc, $\cl{x} = \cl{y}$, ce qui est absurde. Donc, $\cl{x} \cap \cl{y} = \emptyset$.
        \item Montrons que $\bigcup_{x \in E} \cl{x} = E$.\\On a $\forall x \in E, \cl{x} \subset E$, donc $\bigcup_{x \in E} \cl{x} \subset E$.\\Soit $y \in E$. On a $y \in \cl{y}$, donc $y \in \bigcup_{x \in E} \cl{x}$. Donc, $E \subset \bigcup_{x \in E} \cl{x}$.\\On en déduit que $\bigcup_{x \in E} \cl{x} = E$.
    \end{itemize}
\end{proof}

\begin{proof}
    Montrons que $+$ est bien définie.

    + est bel et bien définie si et seulement si elle ne dépend pas du choix du représentant de la classe d'équivalence.

    Soient $\overline{x} = \overline{x'}$ et $\overline{y} = \overline{y'}$ dans $\Z/n\Z$.

    Donc $x \equiv x' [n]$ et $y \equiv y' [n]$.

    i.e. $\exists k_1,k_2 \in \Z$ tels que $x' = x + k_1 n$ et $y' = y + k_2 n$.

    Donc $x+y = x' + y' - (k_1 + k_2) n$.
    Donc $x + y \equiv x' + y' [n]$.
    Donc, $\overline{x+y} = \overline{x' + y'}$.
\end{proof}

\begin{property}
    Soit $(G,\cdot)$ un groupe fini d'élément neutre $e$. Alors, $\forall a \in G, o(a) = \min(\{k\in\N^* \tq a^k = e\})$.
\end{property}

\begin{proof}
    \begin{outline}
        \1 On pose $A=\{k\in\N^* \tq a^k = e\}$. Montrons que $A\ne\emptyset$ et qu'il admet un minimum.

        \2 On a $\angles{a} = \{a^n \tq n \in \Z\} \subset G$ est fini.

        \2 Donc l'application $\psi : \begin{cases}\overbrace{\Z}^{\text{infini}} \to \overbrace{\angles{a}}^{\text{fini}}\\ n \mapsto a^n\end{cases}$ n'est pas injective.

        \3 \textit{Si l'application était injective, alors on aurait une infinité d'éléments dans $\langle a \rangle$, ce qui est absurde, car $\langle a \rangle$ est fini.}

        \2 Donc $\exists k_1,k_2 \in \Z \tq a^{k_1} = a^{k_2}$ et $k_1 \ne k_2$.

        \2 Donc $a^{\abs{k_1 - k_2}} = e$.

        \2 Donc $\abs{k_1 - k_2} \in A$, donc $A \ne \emptyset$.

        \2 En plus, $A \subset \N^*$, donc $A$ admet un minimum. On pose $S=\min(A)$.

        \1 Montrons que $\angles{a} = \{e, a, a^2, \ldots, a^{S-1}\}$.

        \2 On pose $H = \{e, a, a^2, \ldots, a^{S-1}\}$.

        \2 On a $H \subset \angles{a}$.

        \2 Soit $x\in\angles{a}$, donc $\exists n \in \Z \tq x = a ^n$.
        \2 Par le théorème de la division euclidienne, il existe $(q,r) \in \Z \times \N$ tels que $n = qS + r$ et $0 \le r < S$.
        \2 Donc, $x = a^n = a^{qS + r} = (a^S)^q * a^r = e^q * a^r = a^r$, sachant que $a^S = e$ (car $S \in A$).
        \2 Donc, $x \in H$.
    \end{outline}
\end{proof}

\begin{theorem}[Théorème de Lagrange]
    Soit $(G,\cdot)$ un groupe fini et $(H,\cdot)$ un sous-groupe de $(G,\cdot)$.

    Alors, $\card{H}$ divise $\card{G}$ (on note $\card{H} \mid \card{G}$).
\end{theorem}

\begin{proof}
    Soit $H$ un sous-groupe de $(G,\cdot)$.

    On définit une relation binaire $\mathcal{R}$ sur $G$ par : $\forall x,y \in G, x\mathcal{R}y \iff x^{-1} \cdot y \in H$.

    \begin{outline}
        \1 Montrons que $\mathcal{R}$ est une relation d'équivalence.
            \2 \textbf{Réflexivité :} Soit $x \in G$. On a $x^{-1} \cdot x = e_G \in H$, donc $x \mathcal{R} x$.
            \2 \textbf{Symétrie :} Soient $x,y \in G$ tels que $x \mathcal{R} y$. Donc, $x^{-1} \cdot y \in H$. Donc, $(x^{-1} \cdot y)^{-1} = y^{-1} \cdot x \in H$ (car $H$ est stable par l'inverse). Donc, $y \mathcal{R} x$.
            \2 \textbf{Transitivité :} Soient $x,y,z \in G$ tels que $x \mathcal{R} y$ et $y \mathcal{R} z$. Donc, $x^{-1} \cdot y \in H$ et $y^{-1} \cdot z \in H$. Donc, $(x^{-1} \cdot y) \cdot (y^{-1} \cdot z) = x^{-1} \cdot z \in H$ (car $H$ est stable par $\cdot$). Donc, $x \mathcal{R} z$.
        \1 On en déduit que $\mathcal{R}$ est une relation d'équivalence sur $G$.
        \1 Soit $\cl{x}$ la classe d'équivalence de $x\in G$ pour la relation $\mathcal{R}$.
            \2 On a $\cl{x} = \{y\in G\tq x^{-1}\cdot y\in H\} = \{x\cdot h\tq h\in H\}$.
        \1 On montre que toutes les classes d'équivalence ont le même cardinal que $H$.
            \2 Soit $x\in G$. On définit l'application $\varphi : \begin{cases}H \to \cl{x}\\ h\mapsto x\cdot h\end{cases}$.
            \2 Montrons que $\varphi$ est une bijection.
                \3 \textbf{Injectivité :} Soient $h_1,h_2 \in H$ tels que $\varphi(h_1) = \varphi(h_2)$. Donc, $x\cdot h_1 = x\cdot h_2$. Donc, $h_1 = h_2$ (en composant à gauche par $x^{-1}$). Donc, $\varphi$ est injective.
                \3 \textbf{Surjectivité :} Soit $y \in \cl{x}$. Donc, $\exists h \in H \tq y = x\cdot h$. Donc, $\varphi(h) = y$. Donc, $\varphi$ est surjective.
            \2 On en déduit que $\varphi$ est une bijection.
            \2 Donc, $\card{\cl{x}} = \card{H}$.
        \1 Soit $\{\cl{x_i}\}_{i\in I}$ l'ensemble des classes d'équivalence de $G$ pour la relation $\mathcal{R}$.
            \2 On a $\bigcup_{i\in I} \cl{x_i} = G$ et les $\cl{x_i}$ sont disjointes deux à deux.
            \2 Donc, $\card{G} = \sum_{i\in I} \card{\cl{x_i}}$.
            \2 Donc, $\card{G} = \sum_{i\in I} \card{H} = \card{H} \cdot \card{I}$.
        \1 On en déduit que $\card{H} \mid \card{G}$.
    \end{outline}
\end{proof}

\begin{property}
    Soit $(A,+,\times)$ un anneau.

    \begin{itemize}
        \item $a\in A$ est dit \textbf{inversible} lorsque $a$ est inversible pour la loi $\times$, i.e. $\exists b \in A \tq a \times b = b \times a = 1_A$. Dans ce cas, $b$ est unique et on le note $a^{-1}$.
        \item On note $A^* = \{a \in A \tq a \text{ est inversible}\}$ l'ensemble des éléments inversibles de l'anneau $(A,+,\times)$.
    \end{itemize}
\end{property}

\begin{property}
    Soit $(A,+,\times)$ un anneau. Alors, $(A^*=\Uset_A,\times)$ est un groupe, appelé le groupe des inversibles de l'anneau $A$.
\end{property}

\begin{proof}[Deuxième exemple]
    
\end{proof}

\begin{property}
    Soient $(A,+,\times)$ un anneau commutatif et $a,b\in A$. Alors, $\forall n \in \N^*, (a + b)^n = \sum_{k=0}^n \binom nk a^k \times b^{n-k}$ (formule du binôme de Newton), avec la convention $a^0 = 1_A$ et $0_A^0 = 1_A$.

    \textit{Preuve par récurrence sur $n$.}
\end{property}

\begin{property}[Idéal engendré par un élément]
    Soit $(A,+,\times)$ un anneau commutatif et $x \in A$. Alors, l'ensemble $xA = \{a \times x \tq a \in A\}$ est le plus petit idéal de $A$ contenant $x$.

    \textbf{Cet idéal est appelé l'idéal engendré par $x$.}
\end{property}

\begin{proof}
    \begin{outline}
        \1 Montrons que $xA$ est un idéal de $A$.
            \2 On a clairement $xA \subset A$.
            \2 Montrons que $xA$ est un sous-groupe de $(A,+)$.
                \3 $0_A = 0_A \times x \in xA$, donc $xA$ est non-vide.
                \3 Soient $u,v \in xA$. Donc, $\exists a,b \in A$ tels que $u = a \times x$ et $v = b \times x$.
                \3 Donc, $u - v = (a - b) \times x \in xA$ (car $(A,+)$ est un groupe).
            \2 Donc, $xA$ est un sous-groupe de $(A,+)$.
            \2 Montrons que $\forall a \in A, \forall y \in xA, a \times y \in xA$.
                \3 Soit $a \in A$ et $y \in xA$. Donc, $\exists b \in A$ tel que $y = b \times x$.
                \3 Donc, $a \times y = a \times (b \times x) = (a \times b) \times x \in xA$ (car $(A,\times)$ est associative).
            \2 Donc, $\forall a \in A, \forall y \in xA, a \times y \in xA$.
        \1 On en déduit que $xA$ est un idéal de $A$.
        \1 Montrons que $xA$ est le plus petit idéal de $A$ contenant $x$. Nous allons, pour cela, montrer que pour tout idéal $I$ de $A$ contenant $x$, on a $xA \subset I$.
            \2 Soit $I$ un idéal de $A$ tel que $x \in I$.
            \2 Soit $y \in xA$. Donc, $\exists a \in A$ tel que $y = a \times x$.
            \2 Or, comme $x \in I$ et que $I$ est stable par la multiplication par un élément de $A$, on a $y = a \times x \in I$. Cela veut dire que $xA \subset I$.
        \1 On en déduit que $xA$ est le plus petit idéal de $A$ contenant $x$.
    \end{outline}
\end{proof}

\begin{proof}
    Soit $x \in A\setminus\{0_A\}$. Montrons que $x$ est inversible dans $A$.

    On sait que $xA$ est un idéal de $A$.

    Donc $xA = \{0_A\}$ ou $xA = A$.

    Or, $x \in xA$ et $x \ne 0_A$, donc $xA \ne \{0_A\}$.

    Donc, $xA = A$.

    On a donc $1_A \in xA$, sachant que $1_A\in A$.

    Donc $\exists y \in A$ tel que $1_A = y \times x$.

    Donc $x$ est inversible dans $A$.

    On en déduit que $(A,+,\times)$ est un corps commutatif.
\end{proof}

\begin{property}
    Soient $(A,+,\times)$ un anneau commutatif et $I,J$ deux idéaux de $A$. Alors,$I\cap J$ et $I + J = \{x + y \tq (x,y) \in I \times J\}$ sont des idéaux de $A$.
\end{property}

\begin{proof}
    \begin{outline}
        \1 Montrons que $I \cap J$ et $I+J$ sont des sous-groupes de $(A,+)$.
            \2 $I\cap J$ est clairement un sous-groupe de $(A,+)$ (propriétés d'intersection).
            \2 Montrons que $I + J$ est un sous-groupe de $(A,+)$.
                \3 $0_A = 0_A + 0_A \in I + J$, donc $I + J$ est non-vide.
                \3 Soient $u,v \in I + J$. Donc, $\exists (a,b), (c,d) \in I \times J$ tels que $u = a + b$ et $v = c + d$.
                \3 Donc, $u - v = (a - c) + (b - d) \in I + J$ (car $I$ et $J$ sont des sous-groupes de $(A,+)$).
            \2 Donc, $I + J$ est un sous-groupe de $(A,+)$.
        \1 Montrons que $\forall a \in A, \forall x \in I \cap J, a \times x \in I \cap J$ et $\forall a \in A, \forall y \in I + J, a \times y \in I + J$.
            \2 Soit $a \in A$ et $x \in I \cap J$. Donc, $x \in I$ et $x \in J$.
            \2 Donc, $a \times x \in I$ (car $I$ est un idéal de $A$) et $a \times x \in J$ (car $J$ est un idéal de $A$).
            \2 Donc, $a \times x \in I \cap J$.
            \2 Soit $a \in A$ et $y \in I + J$. Donc, $\exists (b,c) \in I \times J$ tel que $y = b + c$.
            \2 Donc, $a \times y = a \times (b + c) = (a \times b) + (a \times c) \in I + J$ (car $I$ et $J$ sont des idéaux de $A$).
    \end{outline}
\end{proof}

\begin{property}
    Soient $f:(G,*)\to(G',T)$ un morphisme de groupes et $e, e'$ les éléments neutres de $G$ et $G'$.
    \begin{enumerate}
        \item $f(e) = e'$.
        \item $\parens{f(x)}^{-1} = f(x^{-1})$, pour tout $x \in G$.
        \item Si $H$ est un sous-groupe de $(G,*)$, alors $f(H)=\{f(h) \tq h \in H\}$ est un sous-groupe de $(G',T)$.
        \item Si $K$ est un sous-groupe de $(G',T)$, alors $f^{-1}(K) = \{g \in G \tq f(g) \in K\}$ est un sous-groupe de $(G,*)$.
    \end{enumerate}
\end{property}

\begin{proof}
    \begin{outline}
        \1 Preuve du premier point
            \2 On a $f(e) = f(e * e) = f(e) T f(e)$.
            \2 En ``simplifiant'' $f(e)$ des deux côtés, on obtient $f(e) = e'$.
        \1 Preuve du quatrième point
            \2 On a $f^{-1}(K) = \{x\in G \tq f(x) \in K\} \subset G$.
            \2 Donc $x^{-1} \in f^{-1}(K) \iff f(x) \in K$.
            \2 Montrons que $f^{-1}(K)$ est non-vide.
                \3 On a $e' \in K$ (car $K$ est un sous-groupe de $(G',T)$).
                \3 Donc, $f(e) = e' \in K$, donc $e \in f^{-1}(K)$.
            \2 Soient $x,y \in f^{-1}(K)$. Montrons que $x * y^{-1} \in f^{-1}(K)$.
                \3 On a $f(x*y^{-1}) = f(x) T f(y^{-1}) = f(x) T \parens{f(y)}^{-1}$ (d'après le deuxième point).
                \3 Or, $f(x) \in K$ et $f(y) \in K$ (car $x,y \in f^{-1}(K)$).
                \3 Donc, $f(x) T \parens{f(y)}^{-1} \in K$ (car $K$ est un sous-groupe de $(G',T)$).
                \3 Donc, $f(x*y^{-1}) \in K$, donc $x * y^{-1} \in f^{-1}(K)$.
            \2 Donc, $f^{-1}(K)$ est un sous-groupe de $(G,*)$.
    \end{outline}
\end{proof}

\begin{property}
    Soient $f:(G,*)\to(G',T)$ un morphisme de groupes et $e, e'$ les éléments neutres de les lois $*$ et $T$ dans $G$ et $G'$ (respectivement).
    
    $f$ est injectif si et seulement si $\Ker(f) = \{e\}$.

    $f$ est surjectif si et seulement si $\Imf(f) = G'$.
\end{property}

\begin{proof}
    \begin{outline}[enumerate]
        \1 Démonstration de la propriété liée à l'injectivité.
            \2 $\implies$ Démonstration dans le sens direct.
                \3 Supposons que $f$ est injectif.
                \3 Soit $x \in \Ker(f)$. Donc, $f(x) = e'$.
                \3 Or, $f(e) = e'$ (d'après la première propriété sur les morphismes de groupes).
                \3 Donc, par injectivité de $f$, on a $x = e$.
                \3 Donc, $\Ker(f) \subset \{e\}$.
                \3 Comme $e \in \Ker(f)$, on a $\{e\} \subset \Ker(f)$.
                \3 Donc, $\Ker(f) = \{e\}$.
            \2 $\impliedby$ Démonstration dans le sens réciproque.
                \3 Supposons que $\Ker(f) = \{e\}$.
                \3 Soient $x,y \in G$ tels que $f(x) = f(y)$.
                \3 Alors, $f(x) T (f(y))^{-1} = e'$.
                \3 Donc, $f(x * y^{-1}) = e'$.
                \3 Donc, $x * y^{-1} \in \Ker(f)$.
                \3 Par hypothèse, $\Ker(f) = \{e\}$, donc $x * y^{-1} = e$.
                \3 Donc, $x = y$.
                \3 Ainsi, $f$ est injectif.
        \1 Démonstration de la propriété liée à la surjectivité.
            \2 Si $f$ est surjectif, alors (équivalence $\iff$) par définition de la surjectivité, on a $f(G) = G'$, ce qui équivaut à dire que $\Imf(f) = G'$. L'équivalence est donc démontrée.
    \end{outline}
\end{proof}

\begin{property}
    Soit $(G,*)$ un groupe. Alors, $(\Aut(G),\circ)$ est un groupe (où $\circ$ est la composition d'applications).
\end{property}

\begin{proof}
    On remarque que $\Aut(G) \subset S_G$, où $S_G$ est l'ensemble des applications bijectives de $G$ dans $G$.

    Montrons alors que $\Aut(G)$ est un sous-groupe de $(S_G,\circ)$.

    \begin{itemize}
        \item On sait déjà que $\Aut(G) \subset S_G$.
        \item L'application identité $\Id_G$ est un automorphisme de $G$, donc $\Id_G \in \Aut(G)$, donc $\Aut(G)$ est non-vide.
        \item Soient $f,g \in \Aut(G)$.
        \item Comme $f$ et $g$ sont bijectives, $f \circ g$ est bijective.
        \item $\forall x,y \in G, (f \circ g)(x * y) = f(g(x * y)) = f(g(x) * g(y)) = f(g(x)) * f(g(y)) = (f \circ g)(x) * (f \circ g)(y)$.
        \item Donc $f \circ g \in \Aut(G)$ (1).
        \item Soit $f \in \Aut(G)$. Montrons que $f^{-1} \in \Aut(G)$.
        \item Comme $f$ est bijective, $f^{-1}$ est bijective.
        \item Soient $x,y \in G$. On a $f^{-1}(x * y) = f^{-1}(f(f^{-1}(x)) * f(f^{-1}(y))) = f^{-1}(f(f^{-1}(x) * f^{-1}(y))) = f^{-1}(x) * f^{-1}(y)$.
        \item Donc, $f^{-1} \in \Aut(G)$ (2).
        \item D'après (1) et (2), on en déduit que $\Aut(G)$ est un sous-groupe de $(S_G,\circ)$.
        \item Donc, $(\Aut(G),\circ)$ est un groupe.
    \end{itemize}
\end{proof}

\begin{property}
    Soit $f:A\to A'$ un morphisme d'anneaux. Alors :
    \begin{enumerate}
        \item L'image de l'image réciproque par $f$ d'un sous-anneau est un sous-anneau.
        \item Si $A$ et $A'$ sont des anneaux commutatifs, que $I$ est un idéal de $(A,+,\times)$, et que $f$ est surjectif, alors $f(I)$ est un idéal de $(A',\oplus,\odot)$.
        \item Si $J$ est un idéal de $(A',\oplus,\odot)$, alors $f^{-1}(J)$ est un idéal de $(A,+,\times)$.
    \end{enumerate}
\end{property}

\begin{proof}
    \begin{outline}
        \1 Preuve du premier point
            \2 Elle est analogue à la preuve du troisième point de la propriété sur les morphismes de groupes.
        \1 Preuve du deuxième point
            \2 Supposons que $f$ est surjectif. Soit $I$ un idéal de $A$.
            \2 On a $f(I)\subset f(A) \subset A'$.
            \2 On a $I$ est un sous-groupe de $(A,+)$, cela implique que $f(I)$ est un sous-groupe de $(A',\oplus)$, car $f$ est un morphisme d'anneaux, donc en particulier un morphisme de groupes.
            \2 Soient $a\in A', y\in f(I)$. Montrons que $ay \in f(I)$.
                \3 Comme $y \in f(I)$, alors $\exists x \in I$ tel que $y = f(x)$.
                \3 Comme $f$ est surjectif, $\exists b \in A$ tel que $f(b) = a$.
                \3 Donc, $a y = f(b) \odot f(x) = f(b \times x) \in f(I)$ (car $f$ est un morphisme d'anneaux), sachant que $b \times x \in I$ (car $I$ est un idéal de $A$).
            \2 Donc, $f(I)$ est un idéal de $(A',\oplus,\odot)$.
        \1 Preuve du troisième point
            \2 Soit $J$ un idéal de $A'$.
            \2 On a $f^{-1}(J)$ est un sous-groupe de $(A,+)$ (propriété sur les morphismes de groupes).
            \2 Soient $a \in A, y \in f^{-1}(J)$. Montrons que $a \times y \in f^{-1}(J)$.
                \3 Comme $y \in f^{-1}(J)$, alors $f(y) \in J$.
                \3 Donc, $f(a \times y) = f(a) \odot f(y) \in J$ (car $f$ est un morphisme d'anneaux), sachant que $f(y) \in J$ et que $J$ est un idéal de $A'$.
                \3 Donc, $a \times y \in f^{-1}(J)$.
            \2 Donc, $f^{-1}(J)$ est un idéal de $(A,+,\times)$.
    \end{outline}
\end{proof}

\begin{proof}
    \begin{itemize}
        \item On pose $E=\{n\in\N^*\tq n\cdot1_A = 0_A\}$.

        \item Dans le cas où $\Ker(f) \neq \{0\}$, on a $\exists n \in \Z \tq n \neq 0$ et $n\cdot 1_A = 0_A$.

        \item Or $-n \in \Ker(f)$ aussi, parce que $(-n)\cdot 1_A = -(n\cdot 1_A) = -0_A = 0_A$, donc $n \in E$.

        \item Donc, $E \neq \emptyset$.

        \item On pose $c=\min(E)$.
        
        \item Montrons que $\Ker(f) = c\Z$. $\Ker(f) = \{n\in\Z \tq n\cdot 1_A = 0_A\}$.

        \item On sait que $c \in \Ker(f)$, donc $c\Z \subset \Ker(f)$.

        \item Soit $n \in \Ker(f)$. Donc, $n\cdot 1_A = 0_A$.

        \item Par division euclidienne, on a $n = qc + r$ avec $q \in \Z$ et $r \in \{0, \ldots, c-1\}$.

        \item Si $r>0$, alors $n-qc\in\N^*$.
        \item Or $n\in\Ker(f)$ et $c\in\Ker(f)$ alors $n - qc \in \Ker(f)$ (car $\Ker(f)$ est un sous-groupe de $(\Z,+)$).
        \item Donc, $r = n - qc \in E$.
        \item Donc $r \ge c$ (car $c = \min(E)$).
        \item Contradiction.
        \item Donc, $r=0$, et par suite $n = qc \in c\Z$.
        \item Donc, $\Ker(f) \subset c\Z$.
        \item On en déduit que $\Ker(f) = c\Z$.
    \end{itemize}
    
\end{proof}

\begin{property}
    Soit $A$ un anneau \textbf{intègre}.
    
    $\Carac(A) = 0$ ou un nombre premier.
\end{property}

\begin{proof}
    On suppose que $\Carac(A) \ne 0$, et $\Carac(A)$ n'est pas premier.

    Donc $\Carac(A) = c = \min(\{n\in\N^*\tq n\cdot1_A = 0_A\})$ est composé.

    Donc $\exists a,b \in \N^*$ tels que $c = a b$, avec $1 \le a < c$ et $1 \le b < c$. (On peut considérer l'inégalité stricte dans le cas où $c\ne 1$).

    On a $c\cdot 1_A = 0_A$.

    Ce qui implique que $(a b)\cdot 1_A = (a\cdot 1_A) \times (b\cdot 1_A) = 0_A$.

    Ce qui implique que $f(ab)=0_A$.
    
    Donc $f(a)\cdot f(b) = 0_A$.

    Comme $A$ est intègre, on a $f(a) = 0_A$ ou $f(b) = 0_A$.

    Cela signifie que soit $a \in \Ker(f)$, soit $b \in \Ker(f)$.

    Donc que $a\cdot 1_A = 0_A$ ou $b \cdot 1_A = 0_A$. Contradiction.

    On conclut que $c=0$ ou $c$ est premier.
\end{proof}


\vspace{1em}


\section*{Extraits : 260128 CH11}


\begin{property}
    \label{property:divisibilite}
    \begin{enumerate}
        \item Pour tout $x\in A$, on a $x\mid x$ (réflexivité).
        \item Pour tout $x,y,z \in A$, si $x\mid y$ et $y\mid z$, alors $x\mid z$ (transitivité).
        \item Pour tout $x,y\in A$, $x\mid y \iff yA \subset xA$. Il faut faire attention à ne pas inverser les deux membres ! Cette relation nous permet de passer de la relation de divisibilité à une inclusion (inégalité).
        \item Pour tout $x,y \in A, (x\mid y \et y \mid x) \iff xA = yA$. Cela équivaut à dire que $\exists \eps \in \mathbb{U}_A \tq x = \eps y$. On dit que $x$ et $y$ sont \textbf{associés} dans ce cas.
        \item La relation de divisibilité n'est pas antisymétrique. Par exemple, dans $\Z$, $2\mid -2$ et $-2\mid 2$, mais $2 \neq -2$.
        \item La relation de divisibilité n'est pas symétrique, par exemple dans $\Z$, $2\mid 4$, mais $4 \nmid 2$.
        \item La relation de divisibilité n'est donc pas une relation d'équivalence.
    \end{enumerate}
\end{property}

\begin{proof}
    \begin{enumerate}
        \item Trivial car $x = x \cdot 1_A$.
        \item $(x \mid y)$ et $(y \mid z)$ impliquent l'existence de $a,b \in A$ tels que $y = xa$ et $z = yb$. Donc $z = (xa)b = x(ab)$, donc $x \mid z$.
        \item ($\impliedby$) : On a $y=y\cdot 1_A \in yA \subset xA$, donc $y \in xA \iff \exists a \in A \tq y = ax$, donc $x\mid y$.\\($\implies$) : Supposons que $x\mid y$, alors $\exists a \in A \tq y = ax$. Soit $t \in yA$, i.e. $\exists \alpha \in A \tq t = y\alpha$.\\Donc $t = (ax)\alpha = x\underbrace{(a\alpha)}_{\in A} \in xA$. Ainsi, $yA \subset xA$.
        \item D'après la propriété précédente, $x\mid y$ et $y \mid x$ équivaut à $yA \subset xA$ et $xA \subset yA$, donc $xA = yA$.\\Maintenant, nous allons montrer que $(x \mid y \et y \mid x) \iff \exists \eps \in \mathbb{U}_n \tq x = \eps y$.\\Dans le sens réciproque, on a $x=\eps y \implies y \mid x$, or $\eps \in \mathbb{U}_A \implies \exists \eps' \in \mathbb{U}_A \tq \eps'x=y$. D'où $x\mid y$.\\Dans le sens direct, on suppose que $x\mid y$ et $y \mid x$. Donc il existe $t,z \in A \tq y=tx \et x=zy$.\\Donc $y=t(zy) = (tz)y$. Donc $y(1_A - tz) = 0_A$. Comme $A$ est intègre, cela veut dire que $y=0_A$ ou $1_A - tz = 0_A \iff tz = 1_A$. Dans le deuxième cas, cela veut dire que $tz=1_A$, donc $z=\eps \in \mathbb{U}_A$, et on a bien $x = zy = \eps y$. Dans le premier cas, on a $y=0_A$, donc $x=zy= z0_A = 0_A$. On peut alors choisir $\eps = 1_A \in \mathbb{U}_A$ et on a bien $x=\eps y$.\\Ainsi, dans tous les cas, on a $\exists \eps \in \mathbb{U}_A \tq x = \eps y$.
    \end{enumerate}
\end{proof}

\begin{theorem}[PGCD dans un anneau principal]
    Soient $A$ un anneau principal et $x_1,\dots,x_n \in A$. Alors il existe un $d \in A$ tel que :

    $$x_1 A + x_2 A + \dots + x_n A = dA$$

    avec $\forall i \in [\![1,n]\!], d \mid x_i$.

    De plus, si $d'$ est un autre élément de $A$ vérifiant ces propriétés, alors $d' \mid d$.

    Dans cas, $d$ est appelé le \textbf{plus grand commun diviseur} (PGCD) des $x_i$.

    On le note $\pgcd(x_1,\dots,x_n)$.
\end{theorem}

\begin{proof}
    \begin{itemize}
        \item On a $\forall i \in [\![1,n]\!], x_iA$ est un idéal principal.
        \item Donc $\sum_{i=1}^{n}x_iA = x_1A + \dots + x_nA$ est aussi un idéal de $A$, qui est principal, puisque l'anneau $A$ est principal.
        \item Donc $\exists d \in A \tq \sum_{i=1}^{n}x_iA = dA$.
        \item On a $\forall i \in [\![1,n]\!], x_i = x_i\times 1_A + \sum_{j=1\et j\neq i}^{n}x_j \times 0_A \in \sum_{i=1}^{n}x_iA = dA$.
        \item Donc $\forall i \in [\![1,n]\!], d \mid x_i$.
        \item Si $d' \in A \tq \forall i \in [\![1,n]\!], d' \mid x_i$.
        \item Alors $\forall i \in [\![1,n]\!], x_iA \subset d'A$.
        \item Comme $dA$ est stable par l'addition (sous-groupe), alors $dA = \sum_{i=1}^{n}x_iA \subset d'A \implies d'|d$.
    \end{itemize}
\end{proof}

\begin{theorem}[Égalité de Bézout]
    Soient $A$ un anneau principal et $x_1,\dots,x_n \in A$. Alors, si $d$ est un $\pgcd(x_1,\dots,x_n)$, alors :
    \[
        \exists u_1,\dots,u_n \in A \tq d = u_1 x_1 + u_2 x_2 + \dots + u_n x_n.
    \]
\end{theorem}

\begin{proof}
    On a $\sum_{i=1}^{n}x_iA = dA$. Donc, $d\in dA = \sum_{i=1}^{n}x_iA$. 
\end{proof}

\begin{property}[Réciproque de l'égalité de Bézout]
    Soient $x_1,\dots,x_n \in A$ (anneau principal), et $d$ un élément de $A$ tel que $\exists u_1,\dots,u_n \in A \tq d = u_1 x_1 + u_2 x_2 + \dots + u_n x_n$.

    $d$ n'est pas forcément le PGCD des $x_i$. Il faut que $d$ soit un diviseur commun des $x_i$ pour que $d$ soit le PGCD des $x_i$.
    
\end{property}

\begin{proof}
    Montrons que $\sum_{i=1}^{n}x_iA = dA$.

    Montrons d'abord l'inclusion $\sum_{i=1}^{n}x_iA \subset dA$.

    Soit $y\in dA$. DOnc $y=da$ pour un certain $a\in A$.

    Donc $y = (\sum_{i=1}^{n}x_iu_i)a=\sum_{i=1}^{n}x_i\underbrace{(u_i a)}_{\in A}$

    Donc $y \in \sum_{i=1}^{n}x_iA$.

    i.e. : $dA \subset \sum_{i=1}^{n}x_i A$.

    On conclut que $\sum_{i=1}^{n}x_iA = dA$.

    D'où $d=\pgcd(x_1,\dots,x_n)$.
\end{proof}

\begin{theorem}[Théorème de Gauss]
    Soient $A$ un anneau principal et $x,y,z \in A$ tels que $x \mid yz$ et $\pgcd(x,z)=1_A$. Alors $x \mid y$.
    
\end{theorem}

\begin{proof}
    On a $x\mid yz$, donc $\exists a \in A \tq yz = xa$.

    On sait aussi que $\pgcd(x,z)=1_A$, cela équivaut à dire que $\exists u_1,u_2 \in A \tq 1_A = u_1 x + u_2 z$.

    Donc $y= xyu_1+yzu_2$.

    Donc $y= xyu_1 + (xa)zu_2 = x\underbrace{(yu_1 + au_2 z)}_{\in A}$.
    
    Donc $x \mid y$.
\end{proof}

\begin{property}
    Soient $x,y,z \in A$ (anneau principal).

    $$(x\wedge y = x \wedge z= 1_A) \iff x\wedge (yz) = 1_A$$
\end{property}

\begin{proof}
    \begin{enumerate}
        \item ($\impliedby$) : Facile.
        \item ($\implies$) : On a $\exists u_1,u_2,v_1,v_2 \in A \tq \begin{cases}xu_1 + y u_2 = 1_A\\xv_1 + z v_2 = 1_A\end{cases}$
        \item Cela implique que $x(\underbrace{xu_1v_1+zu_1v_2+yu_2v_1}_{\in A}) + yz(u_2v_2) = 1_A$.
        \item Donc $x \wedge (yz) = 1_A$.
    \end{enumerate}
\end{proof}

\begin{proof}[Démonstration de la conséquence \ref{consequence:pgcd_produit}]
    En utilisant la propriété précédente, on peut
    
\end{proof}

\begin{property}
    Soient $A$ un anneau principal et $x,y,z$ dans $A$.

    Si $x\mid y$ et $y\mid z$ avec $\pgcd(x,y)=1_A$, alors $xy \mid z$.
\end{property}

\begin{proof}
    On sait que $\exists a,b \in A \tqp z = xa$ et $z=yb$.

    Donc $\exists u,v \in A \tqp xu + yv = 1_A$.

    \begin{align*}
        \implies z &= z(xu + yv)\\
        &= xybu + xyav = xy(bu + av)\\
        &implies xy \mid z.
    \end{align*}
\end{proof}

\begin{proof}
    On procède par récurrence sur $n$.
\end{proof}

\begin{property}
    Soient $t,x,y \in \Z$.

    On a $(tx)\wedge (ty) = |t| (x \wedge y)$.
\end{property}

\begin{proof}
    On pose $d=(tx)\wedge (ty)$.

    On a alors $(tx)\Z + (ty)\Z = d\Z$ (définition du PGCD dans $\Z$).

    Ce qui implique que $t(x\Z + y\Z) = d\Z$ (opérations sur les ensembles).

    Donc $t(x \wedge y)\Z = d\Z$.

    Donc $|t(x \wedge y)| = d$ (car ces deux entiers sont associés, voir propriété \ref{property:divisibilite}).

    Donc $d = |t| (x \wedge y)$.
\end{proof}

\begin{property}
    Soient $A$ un anneau principal et $x,y \in A$.

    $$x\wedge y = d \iff \exists x_1,y_1 \in A \tq \begin{cases}
        x = dx_1\\
        y= dy_1\\
        x_1 \wedge y_1 = 1_A
    \end{cases}$$
\end{property}

\begin{proof}
    Dans le sens direct $(\implies)$, on suppose que $d\mid x$ et $d\mid y$.

    Donc $\exists x_1,y_1 \in A \tq x=dx_1$ et $y=dy_1$.

    Or $x\wedge y = d = (dx_1) \wedge (dy_1)$

    Donc $d = d(x_1 \wedge y_1)$.

    Alors $d(1_A - (x_1 \wedge y_1)) = 0_A$.

    Si $d = 0_A$, alors $x=y=0_A$, donc on peut choisir $x_1=y_1=1_A$ et on a bien $x_1 \wedge y_1 = 1_A$.

    Si $d \neq 0_A$, comme $A$ est intègre, on a $1_A - (x_1 \wedge y_1) = 0_A$, donc $x_1 \wedge y_1 = 1_A$.

    Dans le sens réciproque $(\impliedby)$, on a $x=dx_1$ et $y=dy_1$.

    On sait que $x_1 \wedge y_1 = 1_A$.

    Donc $\exists u,v \in A \tq 1_A = u x_1 + v y_1$.

    Donc $d = d(u x_1 + v y_1) = u (dx_1) + v (dy_1) = ux + vy$.

    Donc $x \wedge y = d$.
\end{proof}

\begin{proof}
    On a $x\mid yz \iff \ovl{\ovl y \cdot \ovl z} = \ovl 0$ dans $\Z/|x|\Z$.

    \begin{align*}
        \ona x\wedge y = 1 &\iff \ovl y \in \Uset_{\Z/|x|\Z}\\
        \text{Hypothèse de départ} &\implies \ovl y \cdot \ovl z = \ovl 0 \text{ dans } \Z/|x|\Z\\
        \implies \ovl y^{-1} \cdot \ovl y \cdot \ovl z &= \ovl y^{-1} \cdot \ovl 0\\
        \ovl z &= \ovl 0 \text{ dans } \Z/|x|\Z\\
        \implies x \mid z.
    \end{align*}
\end{proof}

\begin{property}[Divisibilité par 3]
    Soit $x=a_n\dots a_0$ l'écriture en base décimale de $x \in \N$.

    $$3 \mid x \iff 3 \mid \sum_{k=0}^{n}a_k$$
\end{property}

\begin{proof}
    On a $x=\sum_{k=0}^{n}a_k\cdot 10^k$.

    \begin{align*}
        \ona 3\mid x &\iff \ovl x = \ovl 0 \text{ dans } \Z/3\Z\\
        &\iff \sum_{k=0}^{n}\ovl{a_k \cdot 10^k} = \ovl 0 = \ovl{\sum_{k=0}^{n}a_k} \text{ dans } \Z/3\Z\\
        &\iff \sum_{k=0}^{n}\ovl{a_k} = \ovl 0 \text{ dans } \Z/3\Z\\
        &\iff 3 \mid \sum_{k=0}^{n}a_k
    \end{align*}
\end{proof}

\begin{theorem}[Division euclidienne]
    Soient $(x,y)\in\Z\times\Z^*$.

    Alors $\exists!(q,r)\in\Z\times\N \tqp : \begin{cases}
    x=qy+r\\
    0 \leq r < |y|
    \end{cases}$

    On dit qu'on a effectué la division euclidienne (qu'on notera D.E.) de $x$ par $y$.

    \begin{itemize}
        \item $q$ est appelé le \textbf{quotient} de la D.E. de $x$ par $y$.
        \item $r$ est appelé le \textbf{reste} de la D.E. de $x$ par $y$.
    \end{itemize}
    
\end{theorem}

\begin{proof}
    \begin{align*}
        \begin{cases}
            x = qy + r\\
            0 \leq r < |y|
        \end{cases} &\iff \begin{cases}
        r = x-qy\\
        0 \le x - qy < |y|
        \end{cases}\\
        &\iff \begin{cases}
        r = x - qy\\
        0 \le \frac{x}{|y|} - q \frac{y}{|y|} < 1
        \end{cases}
    \end{align*}

    Si $y > 0$, alors on a $\begin{cases}
    r = x-qy\\
    0 \le \frac{x}{y} - q < 1
    \end{cases}$

    On prend $q=E(\frac{x}{y}) \in \Z$, et $r = x-qy \in \N^*$.

    Si $y < 0$, alors on a $\begin{cases}
    r = x-qy\\
    -q \le -\frac{x}y < 1 - q
    \end{cases}$

    Donc on prend $q = -E(-\frac{x}{y}) \in \Z$, et $r = x-qy \in \N^*$.

    D'où l'existence et l'unicité de $(q,r)$.
\end{proof}

\begin{property}
    Soient $(x,y) \in \Z \times\Z^*$.

    On a $\exists!(q,r)\in\Z\times\N^* \tqp \begin{cases}
    x=qy+r\\
    0 < r \leq |y|
    \end{cases}$

    Le PGCD de $x$ et $y$ est égal au PGCD de $y$ et $r$.

    $$x\wedge y = y \wedge r$$
\end{property}

\begin{proof}
    On pose $d=x\wedge y$, et $d' = y \wedge r$.

    \textbf{Méthode 1 (plus facile) en utilisant les combinaisons linéaires :} (non abordée).

    \textbf{Méthode 2 en utilisant la définition du PGCD :}

    \begin{align*}
        d\Z &= x\Z + y\Z\\
        &= (qy+r)\Z + y\Z\\
        &\subset r\Z + y\Z=d'\Z\\
    \end{align*}

    De même, $d'\Z = y\Z + r\Z \subset x\Z + y\Z = d\Z$.

    Donc $|d| = |d'| \implies d = d'$ car ils sont positifs.

    On en déduit que $x \wedge y = y \wedge r$.

\end{proof}

\begin{property}[Équation $ax+by=c$]
    Soient $A$ un anneau principal et $a,b,c \in A \tq (a,b)\neq (0_A,0_A)$.

    On pose $d = a \wedge b$.

    Les solutions de l'équation $ax + by = c$ sont données par :
    $$\begin{cases}
        \si d \nmid c, \text{ il n'y a pas de solution}\\
        \si d \mid c, \exists \text{ une solution }(x_0,y_0) \et S= \{(x,y) = (x_0 + \frac{b}{d}k, y_0 - \frac{a}{d}k) \text{ avec } k \in A\}
    \end{cases}$$
    
\end{property}

\begin{proof}
    \begin{align*}
        a \wedge b = d &\iff \exists a_1,b_1 \in A \tqp a = da_1 \et b = db_1 \et a_1 \wedge b_1 = 1_A\\
    \end{align*}

    Donc $(E) : d(a_1 x + b_1 y) = c$.

    Dans le cas où $d\nmid c$, il n'y a pas de solution.

    Dans le cas contraire, on a $c = d c_1$ avec $c_1 \in A$.

    On va donc trouver l'équation $(E) : a_1 x + b_1 y = c_1$, sachant que l'anneau $A$ est intègre (donc on peut factoriser par $d$, puis omettre le cas où $d=0$, parce que $d$ est un élément non nul).

    Comme $a_1 \wedge b_1 = 1_A$, il existe un certain couple $(u,v)$ tels que $a_1 u + b_1 v = 1_A$.

    En multipliant par $c_1$, on a $a_1 (c_1 u) + b_1 (c_1 v) = c_1$.

    Donc $(x_0,y_0) = (c_1 u, c_1 v)$ est une solution particulière de $(E)$.

    Soit $(x,y)$ une autre solution de $(E)$.

    On obtient le système suivant :

    $$\begin{cases}
    a_1 x + b_1 y = c_1\\
    a_1 x_0 + b_1 y_0 = c_1
    \end{cases}\implies
    a_1(x - x_0) = b_1 (y - y_0)
    $$

    \begin{center}
    \textit{(Attention : à partir de ce stade, on ne travaille que par implications...)}
    \end{center}

    Cela implique que $b_1 \mid a_1 (x - x_0)$.

    Cela implique que $b_1 \mid (x - x_0)$ car $a_1 \wedge b_1 = 1_A$ (théorème de Gauss).

    Donc $\exists k \in A \tq x - x_0 = b_1 k$.

    Donc $\exists k \in A \tq x = x_0 + b_1 k$.

    En remplaçant dans l'équation $a_1 (x - x_0) = b_1 (y - y_0)$, on obtient :

    \begin{align*}
        a_1 (b_1 k) &= b_1 (y - y_0)\\
        \implies a_1 k &= y - y_0 \scq b_1 \ne 0 \et A \text{ est un anneau intègre }\\
        \implies y &= y_0 + a_1 k
    \end{align*}

    On trouve ensuite que toutes les solutions de l'équation $(E)$ sont de la forme :

    $$(x,y) = (x_0 + b_1 k, y_0 - a_1 k) \text{ avec } k \in A$$

    \textbf{Maintenant, on doit procéder dans le sens réciproque (il ne faut pas oublier qu'on a procédé par implications).}

    On a $a(x_0 + b_1 k) + b(y_0 - a_1 k) = d(a_1 u c_1 + a_1 b_1 k + b_1 v c_1 - b_1 a_1 k) = d(a_1 u c_1 + b_1 v c_1) = d c_1 = c$.

    On conclut que toutes les solutions de l'équation $(E)$ sont de la forme (dans le cas où $d \mid c$) :

    $$S = \{(x,y) = (x_0 + \frac{b}{d}k, y_0 - \frac{a}{d}k) \tq k \in A\}$$

    Où $(x_0,y_0)$ est une solution particulière de l'équation $ax + by = c$.

    Évidemment, la notation $\frac{a}{d}$ et $\frac{b}{d}$ est valide car $d \mid a$ et $d \mid b$, mais elle n'est pas rigoureuse. Pour être rigoureux, il faudrait écrire $a_1$ et $b_1$ à la place de $\frac{a}{d}$ et $\frac{b}{d}$, en mentionnant le fait qu'ils existent etc...

\end{proof}

\begin{theorem}[Lemme]
    $\forall a \in \N^*\setminus\{1\}, \exists p \in \mathcal{P} \tq p \mid a$.
\end{theorem}

\begin{proof}
On pose $E = \{k\in\N^* \setminus \{1\} \tq \forall p \in \mathcal{P}, k \mid a\}$.

On a $q \in E \implies E \neq \emptyset$.

Donc $E \subset \N^*$.

Alors $E$ admet un minimum $p$.

Supposons que $p$ ne soit pas premier. On aura alors :

$\exists 1 < q < p \tq q \mid p$. Or $p \mid a$, alors :

$q \mid a \et q \in \N^* \setminus \{1\} \implies q \in E$.

Ce qui implique que $q\ge \min(E) = p$. Contradiction.

DOnc $p$ est premier et divise $a$.

\end{proof}

\begin{property}[Infinité des nombres premiers]
    Il existe une infinité de nombres premiers. L'ensemble $\mathcal{P}$ est infini.
\end{property}

\begin{proof}
    Supposons que $\mathcal{P}$ soit fini, et notons-le $\mathcal{P} = \{p_1,\dots,p_r\}$, avec chaque $p_i$ deux à deux distincts.

    On pose $a = 1 + \prod_{k=1}^{r} p_k > p_i \forall i \in \{1,\dots,r\}$.

    Donc $\begin{cases}
    a \notin \mathcal{P}\\
    a > 1
    \end{cases}\et \exists j \in \{1,\dots,r\} \tq p_j \mid a$.

    Or $\forall j \in \{1,\dots,r\}, p_j \wedge a = 1$, ce qui implique que $p_j 1$ car $p_j \mid a$. C'est une contradiction.

    Donc $\mathcal{P}$ est infini.
\end{proof}

\begin{proof}
    \textbf{Preuve du deuxième point :}

    On a $\legbinom kp = \frac{p!}{k!(p-k)!}$.

    Donc $\legbinom kp = \frac pk \times \legbinom{k-1}{p-1}$.

    Donc $k\legbinom kp = p \times \legbinom{k-1}{p-1}$.

    Donc $p \mid k \legbinom kp$.

    Or, comme $1 \le k \le p-1$, on a $p \wedge k = 1$.

    Donc $d = p \wedge k = 1$.
\end{proof}

\begin{theorem}[Théorème de Fermat]
    $$\forall n \in \Z,\forall p \in \Pset, n^p \equiv n [p]$$
\end{theorem}

\begin{proof}[Preuve classique par récurrence]
    
\end{proof}

\begin{proof}[Preuve utilisant des notions de structures algébriques]
    \begin{center}
        On rappelle que si $(G,\cdot)$ est un groupe fini de cardinal $n \in \N^*$, alors $\forall a \in G, a^n = e$.
    \end{center}

    On sait que $(\Uset_{\Z/p\Z}, \times)$ est un groupe fini de cardinal $p-1$.

    On a $\forall n \in \Z$, \textbf{si $p\mid n$}, alors $\ovl n = \ovl 0$ et $\ovl n^p = \ovl 0 = \ovl n$. Cela implique que $n^p \equiv n [p]$.

    Le théorème est donc vrai dans ce cas.

    \textbf{Si $p \nmid n$}, c'est-à-dire que $n\wedge p = 1$, alors $\ovl n \in \Uset_{\Z/p\Z}$.

    Donc $\ovl n^{p-1} = \ovl 1$ (car le cardinal de $\Uset_{\Z/p\Z}$ est $p-1$).

    Cela implique que $\ovl n^p = \ovl n$.

    D'où $n^p \equiv n [p]$.

    Le théorème est donc vrai dans ce cas aussi.
\end{proof}

\begin{theorem}
    Soit $x \in \N^*\setminus\{1\}$. Il existe $p_1,\dots,p_r \in \Pset$ tels que $x=\prod_{k=1}^{r}p_k$.
\end{theorem}

\begin{proof}
    Preuve par récurrence (forte) sur $x\in\N^*\setminus\{1\}$.

    \textbf{Initialisation :} Pour $x=2$, on a $2 \in \Pset$, donc la propriété est vraie.

    \textbf{Hérédité :} Supposons que la propriété soit vraie pour tout $k \in \N^*\setminus\{1\}$ tel que $2 \le k < x$, avec $x \ge 3$.

    Si $x+1 \in \Pset$, la propriété est vraie.

    Sinon, $\exists q \in \Pset \tq q \mid x+1$, ce qui implique que $x+1 = q \times a$ pour un certain $a \in \N^*\setminus\{1\} \tq a < x+1$.

    Cela implique que $a\le x$, donc en utilisant l'hypothèse de récurrence, on a $a = \prod_{k=1}^{r} p_k$ avec $p_1,\dots,p_r \in \Pset$.

    Donc $x+1 = q \times \prod_{k=1}^{r} p_k$, en posant $p_{r+1} = q$.
\end{proof}

\begin{proof}[Existence et unicité]
    \textbf{Démonstration de l'existence :}

    On pose $A = \{k \in \N \tq p^k \mid x\}$.

    \textbf{Démonstration de l'unicité :}

    Supposons que $x = p^{\alpha_1}y_1 = p^{\alpha_2}y_2$ avec $p \wedge y_1 = 1$ et $p \wedge y_2 = 1$.

    Supposons que $\alpha_1 \ne \alpha_2$.

    On posera $lambda = \alpha_1 - \alpha_2 > 0$ (sans perte de généralité).

    On aura donc $p^{\lambda} y_1 = y_2$.

    Ce qui implique que $p\mid y_2$, ce qui est une contradiction, car cela implique que $p = 1$.

    On en déduit que $\alpha_1 = \alpha_2$ et $y_1 = y_2$.
    
\end{proof}

\begin{property}
    Si $x=\prod_{i=1}^r  p_i^{\alpha_i}$ avec les $p_i$ des nombres premiers deux à deux distincts, alors $v_{p_i}(x) = \alpha_i$ pour tout $i \in \{1,\dots,r\}$, et $v_p(x) = 0$ pour tout $p \in \Pset\setminus\{p_1,\dots,p_r\}$.

    Et on a $\forall p \in \Pset\setminus\{p_1,\dots,p_r\}, v_p(x) = 0$ (les nombres premiers qui ne divisent pas $x$ ont une valuation nulle).
\end{property}

\begin{proof}
    On a $\forall i \in \{1,\dots,r\}, x = p_i^{\alpha_i}\cdot \underbrace{\prod_{\substack{k=1 \\ k\ne i}}^{r}p_k^{\alpha_k}}_{=y_i}$.

    Les $p_k$ (premiers) sont deux à deux distincts, donc $p_i \wedge \prod_{k=1~;~k\ne i}^{r}p_k^{\alpha_k} = 1$.

    Cela implique que $p_i \wedge y_i = 1$ pour tout $i \in \{1,\dots,r\}$.

    On en déduit que $v_{p_i}(x) = \alpha_i$ pour tout $i \in \{1,\dots,r\}$.

    Soit $p \in \Pset\setminus\{p_1,\dots,p_r\}$.

    On a $x=p^0 \cdot x$, et $p\notin\{p_1,\dots,p_r\} \implies \forall i, p\wedge p_i = 1 \implies p \wedge x = 1$.

    Donc $v_p(x) = 0$ pour tout $p \in \Pset\setminus\{p_1,\dots,p_r\}$.
\end{proof}

\begin{theorem}[Décomposition primaire]

    Soit $x\in \N^*\setminus\{1\}$.

    Il existe $p_1<\dots<p_r \in \Pset$ deux à deux distincts, et $\alpha_1,\dots,\alpha_r \in \N^*$ tels que $x = \prod_{k=1}^{r} p_k^{\alpha_k}$.

    Cette décomposition est unique.

    Si on n'impose pas l'ordre $p_1 < \dots < p_r$, alors la décomposition est unique à l'ordre/la permutation près.

    i.e. si $x=\prod_{i=1}^s q_i^{\beta_i}$ avec $q_1 < \dots < q_s \in \Pset$ et $\beta_j \in \N^* \forall j$, alors $r=s, \forall 1 \le i \le r : \begin{cases}
    p_i = q_i\\
    \alpha_i = \beta_i
    \end{cases}$
\end{theorem}

\begin{proof}
    L'existence a déjà été démontrée.

    \textbf{Démonstration de l'unicité :}

    Supposons que $x=\prod_{k=1}^r p_k^{\alpha_k} = \prod_{i=1}^s q_i^{\beta_i}$ avec $\begin{cases}p_1<\dots<p_r \in \Pset\\q_1<\dots<q_s \in \Pset\end{cases}$ et $\alpha_k,\beta_i \in \N^*$.

    Soit $k=1,\dots,r$. On a $p_k \mid \prod_{i=1}^s q_i^{\beta_i}$.

    On sait que $p_k$ est premier, donc $\exists i \in \{1,\dots,s\} \tq p_k \mid q_i^{\beta_i} \implies p_k \mid q_i$.

    Et on sait aussi que $q_i$ est premier, donc $p_k = q_i$, cela implique que $p_k = 1$ ou $p_k = q_i$.

    Donc $\{p_1,\dots,p_r\}\subset \{q_1,\dots,q_s\}$.

    De même, on trouve que $\{q_1,\dots,q_s\} \subset \{p_1,\dots,p_r\}$.

    Donc $\{p_1,\dots,p_r\} = \{q_1,\dots,q_s\}$, ce qui implique que $r=s$ et $p_i = q_i$ pour tout $i \in \{1,\dots,r\}$.

    On a $x = \prod_{k=1}^r p_k^{\alpha_k} = \prod_{i=1}^r p_i^{\beta_i}$.

    D'après la propriété précédente, on a $v_{p_i}(x) = \alpha_i$ et $v_{p_i}(x) = \beta_i$ pour tout $i \in \{1,\dots,r\}$.

    Donc $\alpha_i = \beta_i$ pour tout $i \in \{1,\dots,r\}$.

    On en déduit que la décomposition primaire de $x$ est unique.
\end{proof}

\begin{property}
    Soient $x,y \in \N^*\setminus\{1\}$.

    $$x\mid y \iff \forall p \in \Pset, v_p(x) \le v_p(y)$$
\end{property}

\begin{proof}
    \textbf{($\implies$) :} Soit $p \in \Pset$. On a donc $x = p^\alpha \cdot z_1$ et $p\wedge z_1 = 1$, avec $\alpha = v_p(x)$.

    Comme $x \mid y$, on a $y = x \cdot k$ pour un certain $k \in \N$.

    De même, il existe un unique $\beta \in \N$, et un unique $z_2 \in \N^*$ tels que $k=p^\beta z_2$ ($\beta = v_p(k)$). On a $p \wedge z_2 = 1$.

    Donc $y = p^{\alpha + \beta} z_1 z_2$.

    Comme $p \wedge z_1 = 1$ et $p \wedge z_2 = 1$, on a $p \wedge z_1 z_2 = 1$.

    Donc $v_p(y) = \alpha + \beta \ge \alpha = v_p(x)$.

    \textbf{($\impliedby$) :} On a $x = \prod_{i = 1}^r p_i^{\alpha_i}$ et $y = \prod_{i=1}^r p_i^{\beta_i}$, avec $p_1,\dots,p_r \in \Pset$ et $\forall 1 \le i \le r, \alpha_i, \beta_i \ge 0$.

    On a $\forall 1 \le i \le r, \alpha_i \le \beta_i$ (parce que $v_{p_i}(x) \le v_{p_i}(y)$).

    Donc $y = \prod_{i=1}^r p_i^{\beta_i} = \prod_{i=1}^r p_i^{\alpha_i} \cdot \prod_{i=1}^r p_i^{\beta_i - \alpha_i} = x \cdot \underbrace{\prod_{i=1}^r p_i^{\beta_i - \alpha_i}}_{\in \N^*}$.

    Donc $x \mid y$.
\end{proof}

\begin{theorem}[Calcul du PGCD par décomposition primaire]
    Soient $x,y \in \N^*\setminus\{1\}$.

    On a $x = \prod_{i=1}^r p_i^{\alpha_i}$ et $y = \prod_{i=1}^r p_i^{\beta_i}$, avec $p_1,\dots,p_r \in \Pset$ et $\alpha_i, \beta_i \in \N$.

    Alors $x \wedge y = \prod_{i=1}^r p_i^{\min(\alpha_i,\beta_i)}$.
\end{theorem}

\begin{proof}
    On pose $d = \prod_{i=1}^r p_i^{m_i}$.

    On a $\forall 1 \le i \le r, m_i \le \alpha_i \et m_i \le \beta_i$.

    Alors $d \mid x$ et $d\mid y$ donc $d\mid x\wedge y$.

    On pose $d' = x \wedge y$, alors $d' \mid x$ et $d' \mid y$.

    Donc $\forall p \in \Pset, \begin{cases}
    v_p(d') \le v_p(x)\\
    v_p(d') \le v_p(y)
    \end{cases}$

    i.e. $\forall 1 \le i \le r, v_{p_i}(d') \le \alpha_i$ et $v_{p_i}(d') \le \beta_i$.

    Donc $\forall 1 \le i \le r, v_{p_i}(d') \le \min(\alpha_i,\beta_i) = m_i$.

    D'après la propriété précédente, on a $d' \mid d$.

    On a $d \mid d'$ et $d' \mid d$, donc $d = d'$, car ils sont tous les deux positifs.

    On en déduit que $x \wedge y = \prod_{i=1}^r p_i^{\min(\alpha_i,\beta_i)}$.
\end{proof}

\begin{proof}
    $\forall 1 \le i \le n, x_i A$ est un idéal de $A$.
    
    Donc $\cap_{i=1}^n x_i A$ est un idéal de $A$, qui est un anneau principal.

    $\implies \cap_{i=1}^n x_i A$ est un idéal principal.

    i.e. $\exists m \in A \tq \cap_{i=1}^n x_i A = mA$.

    Donc $\forall 1 \le i \le n, mA \subset x_i A \implies x_i \mid m$.

    Si $\exists M \tq \forall 1 \le i \le n, x_i \mid M$, alors $M \in x_i A$ pour tout $i$. 

    Donc $MA \subset \cap_{i=1}^n x_i A = mA$.

    Donc $m \mid M$.
\end{proof}

\begin{property}
    \begin{enumerate}
        \item $\forall x,y \in \Z, \forall \lambda \in \Z, (\lambda x) \vee (\lambda y) = \abs{\lambda} (x \vee y)$.
        \item Si $x,y \in \Z \tqp x \wedge y = 1$, alors $x \vee y = \abs{x}\cdot\abs{y}$.
        \item $\forall x,y\in\Z, (x \wedge y) \cdot (x \vee y) = \abs{x}\cdot\abs{y}$.
    \end{enumerate}
\end{property}

\begin{proof}
    \textbf{Preuve du premier point :}

    \begin{align*}
    (\lambda x)\Z \cap (\lambda y)\Z &= (\lambda x \vee \lambda y)\Z\\
    \lambda (x\Z \cap y\Z) &= \lambda(x \vee y)\Z\\
    (\lambda x \vee \lambda y)\Z &= \lambda (x \vee y)\Z\\
    \implies (\lambda x) \vee (\lambda y) &= \abs{\lambda} (x \vee y)
    \end{align*}

    \textbf{Preuve du deuxième point :}

    On pose $m= x \vee y$.

    Donc $x\mid m$ et $y \mid m$.

    Comme $x \wedge y = 1$, on a $xy \mid m$ (propriété de l'application).

    Et on a $m \mid xy$ (car $m \mid x$ et $m \mid y$).

    Donc $m = \abs{xy}$.

    \textbf{Preuve du troisième point :}

    On pose $d=x \wedge y$, donc $\exists x_1,y_1 \in \Z \tqp x= dx_1, y=dy_1 \et x_1 \wedge y_1 = 1$.

    Donc, d'après le deuxième point, on a $x_1 \vee y_1 = \abs{x_1}\cdot\abs{y_1}$.

    Donc $d^2(x_1 \vee y_1) = d^2 \abs{x_1}\cdot\abs{y_1}$.

    Donc $\abs{d}(dx_1 \vee dy_1) = \abs{dx_1}\cdot\abs{dy_1}$.

    Donc $(x \wedge y)(x \vee y) = \abs{x}\cdot\abs{y}$.
\end{proof}

\begin{theorem}
    Soient $x = \prod_{i=1}^{r}p_i^{\alpha_i}$ et $y = \prod_{i=1}^{r} p_i^{\beta_i}$, avec $p_1,\dots,p_r \in \Pset$ et $\alpha_i,\beta_i \in \N$.

    Alors $x \vee y = \prod_{i=1}^{r} p_i^{\max(\alpha_i,\beta_i)}$.
\end{theorem}

\begin{proof}
    On a $xy = \prod_{i=1}^{r}p_i^{\alpha_i + \beta_i}$.

    Or $x \vee y = \frac{xy}{x\wedge y}$.

    Donc $x \vee y = \prod_{i=1}^{r} p_i^{\alpha_i + \beta_i - \min(\alpha_i,\beta_i)} = \prod_{i=1}^{r} p_i^{\max(\alpha_i,\beta_i)}$.

    Sachant que $\alpha_i + \beta_i - \min(\alpha_i,\beta_i) = \max(\alpha_i,\beta_i)$.
\end{proof}


\vspace{1em}


\section*{Extraits : 260206 CH12}


\begin{theorem}
    $(A^\N, +, \times)$ est un anneau commutatif. On l'appelle l'anneau des polynômes à coefficients dans $A$ et on le note $A[X]$.
\end{theorem}

\begin{proof}
    Comme $(A,+)$ est un groupe abélien (sachant que $(A,+,\times)$ est un anneau), il en est de même pour $(A^\N, +)$. \textit{(On rappelle que le groupe des applications de $\N$ dans un groupe abélien est lui-même un groupe abélien)}
    
    La loi $\times$ est bien définie ; c'est une LCI dans $A^\N$ (stabilité par l'addition et la multiplication). De plus, la multiplication $\times$ est associative et commutative, et elle distribue par rapport à l'addition $+$. L'élément neutre de l'anneau $A^\N$ pour la loi $+$ est la suite nulle $(0)_n$ et l'élément neutre pour la loi $\times$ est la suite $(1, 0, 0, \ldots)$. 

    \begin{outline}
        \1 Pour démontrer que $\times$ est une LCI, on utilise les propriétés de l'anneau $A$ et la définition du produit $(c_n)_n$.
        \1 Pour démontrer que $\times$ est associative, on considère trois suites $(x_n)_n$, $(y_n)_n$, et $(z_n)_n$ dans $A^\N$ et on montre que $((x_n)_n \times (y_n)_n) \times (z_n)_n = (x_n)_n \times ((y_n)_n \times (z_n)_n)$.
            \2 Pour cela, on va montrer que tous les termes de ces deux suites sont égaux.
            \2 On pose $(c_n)_n\times (z_n)_n = (\alpha_n)_n$ et on trouve que $\alpha_n = \sum_{k+l+j=n}(x_k \times y_l)\times z_j$.
            \2 De même, on pose $(x_n)_n\times ((y_n)_n\times(z_n)_n) = (\beta_n)_n$ et on trouve que $\beta_n = \sum_{k+l+j=n}x_k\times (y_l \times z_j)$.
            \2 En utilisant l'associativité de la multiplication dans $A$, on trouve que $\alpha_n = \beta_n$ pour tout $n$, \textbf{ce qui prouve l'associativité de $\times$}.
            \2 Il faut noter qu'on a utilisé la distributivité de la multiplication par rapport à l'addition dans $A$ pour réarranger les termes de ces sommes.
        \1 Pour démontrer que $\times$ est commutative, on considère deux suites $(x_n)_n$ et $(y_n)_n$ dans $A^\N$ et on montre que $(x_n)_n \times (y_n)_n = (y_n)_n \times (x_n)_n$.
            \2 On pose une suite $(c_n)_n = \sum_{i+j=n} x_i \times y_j$ et on pose $(d_n)_n = \sum_{i+j=n} y_i \times x_j$.
            \2 En utilisant la commutativité de la multiplication dans $A$, on trouve que $c_n = d_n$ pour tout $n$, \textbf{ce qui prouve la commutativité de $\times$}.
        \1 Pour déterminer l'élément neutre de la multiplication $\times$, on doit trouver une suite $e = (e_n)_n$ dans $A^\N$ telle que pour toute suite $(x_n)_n$ dans $A^\N$, on ait $e \times (x_n)_n = (x_n)_n$ et $(x_n)_n \times e = (x_n)_n$.
            \2 On pose $e = (e_n)_n \in A^\N$ définie par $e_0 = 1_A$, et $\forall n \ge 1, e_n = 0_A$.
            \2 On utilisera le \textit{Symbole de Kronecker} pour montrer que $e$ est l'élément neutre pour la multiplication $\times$.
            \3 Le symbole de Kronecker $\delta_{x,y}$ est défini par $\delta_{x,y} = 1$ si $x=y$ et $\delta_{x,y} = 0$ sinon.
            \2 Soit $(x_n)_n \in A^\N$. On pose $e\times (x_n)_n = (c_n)_n$.
            \2 En utilisant la définition du produit, on trouve que $c_n = \sum_{k=0}^n \delta_{k,0} \times x_{n-k}$.
            \2 En utilisant la définition du symbole de Kronecker, on trouve que $c_n = x_n$ pour tout $n$, et sachant que la multiplication est commutative dans $A$, \textbf{ce qui prouve que $e$ est l'élément neutre pour la multiplication $\times$}.
        \1 Pour montrer que $\times$ est distributive par rapport à $+$, on considère trois suites $(x_n)_n$, $(y_n)_n$, et $(z_n)_n$ dans $A^\N$ et on montre que $(x_n)_n \times ((y_n)_n + (z_n)_n) = (x_n)_n \times (y_n)_n + (x_n)_n \times (z_n)_n$.
            \2 On a $(x_n)_n \times ((y_n)_n + (z_n)_n) = (x_n)_n \times (y_n+z_n)_n$.
            \2 En utilisant la définition du produit, on trouve que $(x_n)_n \times ((y_n)_n + (z_n)_n)\parens{\sum_{i+j} x_i \times (y_j + z_j)}_n$.
            \2 En séparant les termes de cette somme, on trouve que $(x_n)_n \times ((y_n)_n + (z_n)_n) = \parens{\sum_{i+j} x_i \times y_j}_n + \parens{\sum_{i+j} x_i \times z_j}_n$.
            \2 En utilisant la définition du produit, on trouve que $(x_n)_n \times ((y_n)_n + (z_n)_n) = (x_n)_n \times (y_n)_n + (x_n)_n \times (z_n)_n$, \textbf{ce qui prouve la distributivité de $\times$ par rapport à $+$}.
        \1 On conclut que $(A^\N, +, \times)$ est un anneau commutatif, car il satisfait toutes les propriétés nécessaires : 
            \2 $(A^\N, +)$ est un groupe abélien.
            \2 $\times$ est une LCI dans $A^\N$.
            \2 $\times$ est associative et commutative.
            \2 $\times$ distribue par rapport à $+$.
            \2 Il existe un élément neutre pour la multiplication $\times$.
    \end{outline}
    
    Par conséquent, $(A^\N, +, \times)$ satisfait la définition d'un anneau commutatif.
\end{proof}

\begin{property}[Conservation de l'intégrité]
    Si $A$ est un anneau intègre, alors $A[X]$ est aussi un anneau intègre.
\end{property}

\begin{proof}
    \begin{outline}
        \1 Supposons que $\exists (a_n)_n, (b_n)_n \in A^\N$ tels que $(a_n)_n \times (b_n)_n = 0_{A^\N}$ et $\begin{cases} (a_n)_n \neq 0_{A^\N}\\ (b_n)_n \neq 0_{A^\N} \end{cases}$.

    \1 On pose $i_0 = \min\parens{\{n\in\N \tq a_n \ne 0\}}$ et $\forall i < i_0, a_i = 0_A$. On pose aussi $j_0 = \min\parens{\{n\in\N \tq b_n \ne 0\}}$ et $\forall j < j_0, b_j = 0_A$. $i_0$ et $j_0$ représentent les indices du premier terme non nul de $(a_n)_n$ et $(b_n)_n$ respectivement.

    \1 On pose $(a_n)_n \times (b_n)_n = (c_n)_n$. Donc $\forall n \in \N, c_n = 0_A = \sum_{i+j=n} a_i \times b_j$.

    \2 On remarque qu'en particulier, pour $n = i_0 + j_0$, on a $c_{i_0 + j_0} = 0_A = \sum_{i+j=i_0+j_0} a_i \times b_j$.

    \2 Cela implique que $\sum_{i+j=i_0+j_0} a_i \times b_j = 0_A$.
    
    \3 Dans le cas où $i < i_0$ ou $j < j_0$, on a $a_i = 0_A$ ou $b_j = 0_A$, donc les termes correspondants de la somme sont nuls.

    \3 Dans le cas contraire, c'est à dire lorsque $i \ge i_0$ et $j \ge j_0$, on a $c_{i_0 + j_0} = 0_A + a_{i_0} \times b_{j_0} = 0_A$, car les autres termes de la somme sont nuls.

    \2 Comme $A$ est un anneau intègre, on a $a_{i_0} \times b_{j_0} = 0_A \implies a_{i_0} = 0_A$ ou $b_{j_0} = 0_A$, ce qui contredit la définition de $i_0$ et $j_0$.

    \1 Par conséquent, notre supposition est fausse, et il n'existe pas de tels éléments $(a_n)_n$ et $(b_n)_n$ dans $A^\N$.

    \1 Ainsi, $A^\N$ est un anneau intègre.
    \end{outline}
\end{proof}

\begin{property}
    $\forall k \in \N, X^k = \underbrace{X\times\dots\times X}_{k\text{ fois}} = (\delta_{n,k})_n$.
\end{property}

\begin{proof}
    On va démontrer cette propriété par récurrence sur $k$.

    \textbf{Initialisation :} Pour $k=0$, on a $X^0 = 1_{A^\N} = (\delta_{n,0})_n$, ce qui est vrai.

    \textbf{Hérédité :} Supposons que la propriété est vraie pour un certain $k \in \N$, c'est-à-dire que $X^k = (\delta_{n,k})_n$. Nous devons montrer qu'elle est également vraie pour $k+1$.

    On a $X^{k+1} = X^k \times X = \underbrace{(\delta_{n,k})_n}_{=(a_n)_n} \times \underbrace{(\delta_{n,1})_n}_{=(b_n)_n}$.

    On a $X^{k+1} = \parens{\sum_{i=0}^{n}a_i\times b_{n-i}}_n$.

    On obtient donc $X^{k+1} = \parens{\sum_{i=0}^{n}\delta_{i,k}\times \delta_{n-i,1}}_n$.

    Pour chaque $n \in \N$, on a :
    \[\sum_{i=0}^{n}\delta_{i,k}\times \delta_{n-i,1} = \begin{cases}
    1_A & \text{si } n = k+1\\
    0_A & \text{sinon}
    \end{cases}\]

    On en déduit que $X^{k+1} = (\delta_{n-1,k})_n$.

    Et on a $\delta_{n-1,k} = 1 \iff n-1 = k \iff n = k+1 \iff \delta_{n,k+1} = 1$.

    On en conclut que $X^{k+1} = (\delta_{n,k+1})_n$.

    Ainsi, par le principe de récurrence, la propriété est vraie pour tout $k \in \N$.
\end{proof}

\begin{property}
    $A[X]$ est un sous-anneau de $(A^\N, +, \times)$.
\end{property}

\begin{proof}
    \begin{outline}
        \1 Soient $(a_n)_n$ et $(b_n)_n$ deux éléments de $A[X]$. Par définition, il existe des entiers naturels $n_a$ et $n_b$ tels que $\forall n \ge n_a, a_n = 0_A$ et $\forall n \ge n_b, b_n = 0_A$.

        \1 Les éléments neutre $1_{A^\N}$ et $0_{A^\N}$ appartiennent à $A[X]$ car ils sont nuls à partir du rang 1.

        \1 On pose $n_0 = \max(n_a, n_b)$. Alors, pour tout $n \ge n_0$, on a :
            \2 $a_n = 0_A$ (car $n \ge n_a$)
            \2 $b_n = 0_A$ (car $n \ge n_b$)
            \2 Donc $a-b \in A[X]$.
        \1 Ainsi, $A[X]$ est stable par soustraction, et par conséquent, il est stable par addition.

        \1 Considérons maintenant le produit $(d_n)_n = (a_n)_n \times (b_n)_n$.\\Pour tout $n \ge n_0$, on a :
            \2 $d_n = \sum_{i=0}^{n} a_i \times b_{n-i}$
            \2 Si $n \ge n_0$, alors pour chaque terme de la somme, soit $i \ge n_a$ ou $n-i \ge n_b$, ce qui implique que soit $a_i = 0_A$ soit $b_{n-i} = 0_A$. Par conséquent, chaque terme de la somme est nul, et donc $d_n = 0_A$.

        \1 Ainsi, $(d_n)_n$ appartient à $A[X]$, ce qui montre que $A[X]$ est stable par multiplication.

        \1 Par conséquent, $A[X]$ est un sous-anneau de $(A^\N, +, \times)$.
    \end{outline}
\end{proof}

\begin{property}
    \begin{enumerate}
        \item $\forall a \in A^\N, 1_A \cdot a = a$.
        \item $\forall \lambda_1,\lambda_2 \in A, \forall a \in A^\N, (\lambda_1 + \lambda_2) \cdot a = \lambda_1 \cdot a + \lambda_2 \cdot a$.
        \item $\forall \lambda \in A, \forall a,b \in A^\N, \lambda \cdot (a + b) = \lambda \cdot a + \lambda \cdot b$.
        \item $\forall \lambda_1,\lambda_2 \in A, \forall a \in A^\N, (\lambda_1 \times \lambda_2) \cdot a = \lambda_1 \cdot (\lambda_2 \cdot a)$.
    \end{enumerate}

    On sait que $\forall a \in A[X], \exists n_0 \in \N, \forall n \ge n_0, a_n = 0_A$.

    i.e. $a = (a_0, a_1, \ldots, a_{n_0}, 0_A, 0_A, \ldots)$.

    \begin{itemize}
        \item On peut encore réécrire $a$ sous la forme $a = (a_0, 0_A, 0_A, \ldots) + (0_A, a_1, 0_A, \ldots) + \ldots + (0_A, 0_A, \ldots, a_{n_0}, 0_A, \ldots)$.

        \item On peut aussi réécrire $a$ sous la forme $a = a_0 \cdot (1_A, 0_A, 0_A, \ldots) + a_1 \cdot (0_A, 1_A, 0_A, \ldots) + \ldots + a_{n_0} \cdot (0_A, 0_A, \ldots, 1_A, 0_A, \ldots)$.
    \end{itemize}

    En utilisant la définition de $X$, on trouve que $a = a_0 \cdot 1_{A^\N} + a_1 \cdot X + \ldots + a_{n_0} \cdot X^{n_0}$.

    Donc $a = \sum_{k=0}^{n_0}a_k \cdot X^k$.
\end{property}

\begin{property}[Préservation de l'intégrité]
    Si $A$ est un anneau intègre, alors $A[X]$ est aussi un anneau intègre.
\end{property}

\begin{proof}
    $A$ intègre $\implies A^\N$ intègre, et $A[X]$ est un sous-anneau de $A^\N$. Donc $A[X]$ est un anneau intègre.
\end{proof}

\begin{property}
    Soient $a,b \in A[X]$. On a $a = b \iff \forall n, a_n = b_n$.

    Donc, en particulier, $\sum a_k\cdot X^k = \sum_{b_k}\cdot X^k \iff \forall k, a_k = b_k$.
\end{property}

\begin{property}
    Soient $A$ un anneau \textit{intègre} (\faIcon{exclamation-circle}) et $P,Q \in A[X]$.

    \begin{enumerate}
        \item $\deg(P\cdot Q) = \deg(P) + \deg(Q)$ et $\cfd(P\cdot Q) = \cfd(P) \times \cfd(Q)$.
        \item $\deg(P+Q) \le \max(\deg(P), \deg(Q))$\\avec égalité si et seulement si $\begin{cases} \deg(P) \ne \deg(Q)\\ \text{ou } \deg(P) = \deg(Q) \text{ et } \cfd(P) + \cfd(Q) \ne 0_A\end{cases}$.
    \end{enumerate}
\end{property}

\begin{proof}
    \begin{enumerate}
        \item Si $P=0$ ou $Q=0$, alors $P\cdot Q = 0$, et $\deg(P\cdot Q) = -\infty = \deg(P) + \deg(Q)$.\\Dans le cas contraire, on pose $P = \sum_{k=0}^{n}a_k X^k$ et $Q = \sum_{k=0}^{m}b_k X^k$ avec $a_n \ne 0_A$ et $b_m \ne 0_A$.\\En utilisant la définition du produit, on trouve que $P\cdot Q = \sum_{k=0}^{n+m}c_k X^k$ avec $c_{k} = \sum_{i+j=k} a_i \times b_j$.\\En particulier, on a $c_{n+m} = a_n \times b_m \ne 0_A$ (car $A$ est intègre), ce qui implique que $\deg(P\cdot Q) = n+m = \deg(P) + \deg(Q)$ et $\cfd(P\cdot Q) = c_{n+m} = a_n \times b_m = \cfd(P) \times \cfd(Q)$.
        \item Si $P = 0 \ou Q= 0$, alors la démonstration est claire.\\Dans le cas contraire, on pose : $P=\sum_{k=0}^{n}a_kX^k, Q = \sum_{k=0}^{m}b_kX^k$. On suppose par exemple (et sans perte de généralité) que $n \le m$, donc $P = \sum_{k=0}^{m}a_kX^k$ avec $\forall k > n, a_k = 0$.\\Donc $(P+Q) = \sum_{k=0}^{m}(a_k+b_k)X^k$.\\Cela implique que $\deg(P+Q) \le m = \max(n,m) = \max(\deg(P), \deg(Q))$.\\Dans le cas où $n < m$, on a $\deg(P+Q) = m = \max(\deg(P), \deg(Q))$.\\Dans le cas où $n = m$, on a $\deg(P+Q) = n = m$ si et seulement si $a_n + b_n \ne 0_A$, c'est à dire $\cfd(P) + \cfd(Q) \ne 0_A$.\\Dans le cas contraire, on a $\deg(P+Q) < n = m$.
    \end{enumerate}
\end{proof}

\begin{theorem}[Formule de Taylor]
    Soient $P \in \Kset[X] \tq \deg(P) = n$ et $a \in \Kset$. Alors :

    $$P = \sum_{k=0}^{n} \frac{P^{(k)}(a)}{k!} (X-a)^k$$
\end{theorem}

\begin{proof}
    On pose $\forall i \in [\![0,n]\!], Q_i = X^i$.

    \begin{align*}
        \ona \forall i \in [\![0,n]\!], Q_i &= (X-a+a)^i\\
        &= \sum_{k=0}^{i} \legbinom{i}{k} (X-a)^k a^{i-k}
    \end{align*}

    \begin{align*}
        \ona \forall k \in [\![0,i]\!], Q_i^{(k)}(a) &= \frac{i!}{(i-k)!} a^{i-k}\\
        \implies Q_i^{(k)} &= k! \cdot \legbinom{i}{k} a^{i-k} (X-a)^k\\
        \implies \forall k \in [\![0,i]\!], \frac{Q_i^{(k)}(a)}{k!} &= \legbinom{i}{k} a^{i-k}\\
    \end{align*}

    Donc $Q_i = \sum_{k=0}^{i} \frac{Q_i^{(k)}(a)}{k!} (X-a)^k$.

    On écrit $P = \sum_{i=0}^{n} \lambda_i X^i = \sum_{i=0}^{n} \lambda_i Q_i$.

    Donc, pour tout $k\le n$, $P^{(k)} = \sum_{i=k}^n \lambda_i Q_i^{(k)}$.

    Donc, pour tout $k \le n$, $\boxed{P^{(k)}(a) = \sum_{i=k}^n \lambda_i Q_i^{(k)}(a)}$.
    
    Comme $P = \sum_{i=0}^{n}\lambda_i Q_i$, on trouve que $P = \sum_{i=0}^{n} \lambda_i \sum_{k=0}^{i} \frac{Q_i^{(k)}(a)}{k!} (X-a)^k$.

    En échangeant les sommes, on trouve que $P = \sum_{k=0}^{n} \frac{1}{k!} \underbrace{\parens{\sum_{i=k}^{n} \lambda_i Q_i^{(k)}(a)}}_{=P^{(k)}(a)} (X-a)^k$.

    Le terme entre parenthèses est égal à $P^{(k)}(a)$, d'après la formule encadrée ci-dessus.

    Donc $P = \sum_{k=0}^{n} \frac{P^{(k)}(a)}{k!} (X-a)^k$.

\end{proof}

\begin{property}
    Les éléments inversibles de $\Kset[X]$ sont les polynômes de degré 0.

    $$P \in \Uset_{\Kset[X]} \iff \deg(P) = 0$$

    $$\Uset_{\Kset[X]} = \Kset^* = \Kset_0[X]\setminus\{0\}$$
\end{property}

\begin{proof}
    \begin{outline}
        \1 Démonstration dans le sens $\implies$ :

        \2 On a $P \in \Uset_{\Kset[X]} \iff \exists Q \in \Kset[X], P\cdot Q = 1$.

        \2 Donc $\deg(P\cdot Q) = \deg(P) + \deg(Q) = 0$, sachant que $P \ne 0$ et $Q \ne 0$ (car $1 \ne 0$).

        \2 Et comme $\deg(P), \deg(Q) \in \N$, on trouve que $\deg(P) = 0$ et $\deg(Q) = 0$.

        \1 Démonstration dans le sens $\impliedby$ :

        \2 Si $\deg(P) = 0$, alors $P$ est de la forme $P = a_0$ avec $a_0 \in \Kset$ et $a_0 \ne 0$ (car $P \ne 0$).

        \2 Comme $\Kset$ est un corps, $a_0$ est inversible dans $\Kset$, donc il existe $a_0^{-1} \in \Kset$ tel que $a_0 \times a_0^{-1} = 1$.

        \2 On pose $Q = a_0^{-1}$. Alors $P\cdot Q = a_0 \cdot a_0^{-1} = 1$, ce qui implique que $P$ est inversible dans $\Kset[X]$.

        \1 Par conséquent, les éléments inversibles de $\Kset[X]$ sont exactement les polynômes de degré 0.
    \end{outline}

\end{proof}

\begin{property}
    Soient $A,B,C \in \Kset[X]$.
    \begin{outline}[enumerate]
        \1 Si $A \mid B$ et $B \mid C$, alors $A \mid C$.
        \1 Si $A \mid B$ et $B \mid A$, alors :
        \2 $A$ et $B$ sont associés.
        \2 Il existe $\lambda \in \underbrace{\Uset_{\Kset[X]}}_{=\Kset^*}$ tel que $A = \lambda \cdot B$.
        \1 $\Kset[X]$ est un anneau intègre, donc toutes les propriétés de divisibilité dans un anneau intègre sont vérifiées dans $\Kset[X]$.
    \end{outline}
\end{property}

\begin{theorem}
    Soient $A,B \in \Kset[X]$ avec $B \ne 0$. 
    Alors il existe un unique couple $(Q,R) \in \Kset[X]^2$ tels que $A = B\cdot Q + R$ et $\deg(R) < \deg(B)$.
\end{theorem}

\begin{proof}
    \begin{outline}
        \1 Si $\deg(B) < \deg(A)$,
        
        \2 Alors $B=0\cdot A + B$ est une division euclidienne de $B$ par $A$.

        \2 Cela implique que $(Q,R) = (0,B)$ et $\deg(R) = \deg(B) < \deg(A)$.

        \1 Si $\deg(B) \ge \deg(A)$,

        \2 Cela implique que $B \ne 0$.

        \2 On pose $B = \sum_{k=0}^{m}b_k X^k$ avec $b_m \ne 0$, et $A = \sum_{k=0}^{n}a_k X^k$ avec $a_n \ne 0$.

        \2 Le degré de $B$ est $m$ et le degré de $A$ est $n$, donc $m \ge n$.

        \2 Nous allons procéder par récurrence forte sur $m$.

        \3 \textbf{Initialisation :} Si $m < n$, la propriété est vérifiée comme montré dans le premier cas.

        \3 \textbf{Hérédité :} Supposons que la propriété est vérifiée pour tout $k \le m - 1$. Nous allons montrer qu'elle est également vérifiée pour $m$.

        \3 On pose $B_1 = B - b_m a_n^{-1} X^{m-n} \cdot A$. Alors $B_1$ est un polynôme de degré inférieur à $m$ (en prenant le degré de $B$ et le degré de $b_m a_n^{-1} X^{m-n} \cdot A$).

        \3 Le coefficient dominant de $B_1$ est nul, car $b_m - b_m a_n^{-1} a_n = 0$.

        \3 Le degré de $B_1$ est donc strictement inférieur à $m$. On peut donc appliquer l'hypothèse de récurrence à $B_1$. On peut donc écrire $B_1 = A \cdot Q_1 + R$ avec $\deg(R) < \deg(A)$.

        Donc $B = A \cdot Q + R$ avec $Q = Q_1 + b_m a_n^{-1} X^{m-n}$.

        \3 On a donc trouvé un couple $(Q,R) \in \Kset[X]^2$ tel que $B = A\cdot Q + R$ et $\deg(R) < \deg(B)$.

        \2 Montrons maintenant l'unicité de ce couple.

        \3 Supposons qu'il existe un autre couple $(Q',R') \in \Kset[X]^2$ tel que $B = A\cdot Q' + R'$ et $\deg(R') < \deg(B)$.

        \3 En soustrayant les deux égalités, on trouve que $A\cdot \underbrace{(Q-Q')}_{=Q} = \underbrace{R' - R}_{=R}$.

        \3 Comme $\deg(AQ) = \deg(A) + \deg(Q) = \deg(R) \le \max(\deg(R), \deg(R'))$, on trouve que $\deg(Q) < 0$, ce qui implique que $Q = 0$, et par conséquent que $R = 0$.

        \3 Donc $Q = Q'$ et $R = R'$, ce qui montre l'unicité du couple $(Q,R)$.
    \end{outline}
\end{proof}

\begin{property}
    L'anneau $\Kset[X]$ est principal.
\end{property}

\begin{proof}
    Soit $I$ un idéal de $\Kset[X]$.

    \begin{outline}
        \1 Si $I = \{0\}$, alors $I = \langle 0 \rangle = 0\cdot\Kset[X]$.

        \1 Si $I$ n'est pas réduit à $\{0\}$, alors il existe un polynôme non nul $P \in I$ de degré minimal $d$.

        \2 On pose $E = \{\deg(P) \tq P \in I\setminus\{0\}\}$

        \2 On a $E \ne \emptyset$ (car $I \ne \{0\}$) et $E \subset \N$. $E$ admet un minimum $d$ (car $\N$ est bien ordonné).

        \2 On peut donc écrire $\exists P_0 \in I \setminus \{0\} \tq \deg(P_0) = d$.

        \2 Montrons que $I = \langle P_0 \rangle = P_0 \cdot \Kset[X]$.

        \2 On a $\forall Q \in \Kset[X]$, $P_0 \cdot Q \in I$ car $P_0 \in I$ et $I$ est un idéal.

        \2 Donc $P_0 \cdot \Kset[X] \subset I$.

        \2 Soit $B \in I$. Par la division euclidienne, il existe un unique couple $(Q,R) \in \Kset[X]^2$ tel que $B = P_0 \cdot Q + R$ et $\deg(R) < \deg(P_0) = d$.

        \2 On a $R = B - P_0 \cdot Q \in I$ (car $B \in I$, $P_0 \cdot Q \in I$, et $I$ est un idéal, donc stable par soustraction).

        \3 Si $R \ne 0$, cela implique que $R \in I\setminus\{0\}$, et donc que $\deg(R) \in E$, ce qui implique que $\deg(R) \ge \deg(P_0) = d$ (car $d$ est le minimum de $E$). C'est une contradiction, parce que $R = B- P_0 \cdot Q$ et $\deg(R) < \deg(P_0) = d$.

        \3 Donc $R = 0$, et par conséquent, $B = P_0 \cdot Q \in P_0 \cdot \Kset[X]$.

        \2 Ainsi, $I \subset P_0 \cdot \Kset[X]$.

        \2 Par conséquent, $I = P_0 \cdot \Kset[X] = \langle P_0 \rangle$.

        \1 Donc, dans tous les cas, $I$ est principal.
    \end{outline}
\end{proof}

\begin{proof}
    Soit $I$ un idéal non-nul.

    Nous allons montrer que $I$ est engendré par un polynôme unitaire de degré minimal.
    
    Donc $\exists P_1 \in \Kset[X] \setminus \{0\} \tq I = \langle P_1 \rangle = P_1 \cdot \Kset[X]$.

    On a donc $I = (\cfd(P_1))^{-1} \cdot P_1 \cdot \Kset[X]$, sachant que $\cfd(P_1) \in \Kset^*$ (car $P_1 \ne 0$ et $\Kset$ est un corps), et que donc les deux polynômes $P_1$ et $(\cfd(P_1))^{-1} \cdot P_1$ sont associés.

    On pose $P_0 = (\cfd(P_1))^{-1} \cdot P_1$. Alors $P_0$ est unitaire, et $I = P_0 \cdot \Kset[X]$.

    Nous allons maintenant montrer que $P_0$ est le \textit{seul} polynôme unitaire de degré minimal qui engendre $I$.

    Si $\exists Q_0 \text{ unitaire } \tq I = P_0 \cdot \Kset[X] = Q_0 \cdot \Kset[X]$, alors $P_0$ et $Q_0$ sont associés.

    Donc $\exists \lambda \in \Kset^* \tq P_0 = \lambda \cdot Q_0$.

    On sait que $\cfd(P_0) = 1$ et $\cfd(Q_0) = 1$, donc $\lambda = 1$, et par conséquent, $P_0 = Q_0$.

    Donc $P_0$ est unique.
\end{proof}

\begin{property}
    Soit $P \in \Kset[X]$ un polynôme à coefficients dans $\Kset$. Alors $a \in \Kset$ est une racine de $P$ si et seulement si $(X-a) \mid P$.
\end{property}

\begin{proof}
    \begin{outline}
        \1 La démonstration dans le sens indirect $\impliedby$ est claire, car si $(X-a) \mid P$, alors $P = (X-a) \cdot Q$ pour un certain $Q \in \Kset[X]$, et donc $P(a) = (a-a) \cdot Q(a) = 0$.
        \1 Dans le sens direct $\implies$, si $P(a) = 0$, alors par la division euclidienne, il existe un unique couple $(Q,R) \in \Kset[X]^2$ tel que $P = (X-a) \cdot Q + R$ et $\deg(R) < \deg(X-a) = 1$. Donc $R$ est une constante. En évaluant en $a$, on obtient $0 = P(a) = (a-a) \cdot Q(a) + R = R$, donc $R = 0$ et $P = (X-a) \cdot Q$, ce qui montre que $(X-a) \mid P$.
    \end{outline}
\end{proof}

\begin{proof}
    On procède par récurrence sur $k \in \N^*$.
    
    \textbf{Initialisation :} Pour $k = 1$, le résultat est trivialement vrai, car $(X-a_1) \mid P$ si $a_1$ est une racine de $P$.

    \textbf{Hérédité :} Supposons que le résultat soit vrai pour un certain $k \ge 1$, c'est-à-dire que si $a_1, a_2, \ldots, a_k$ sont des racines deux à deux distinctes de $P$, alors $(X-a_1)(X-a_2)\cdots(X-a_k) \mid P$.
    
    Soit $a_{k+1}$ une racine distincte de $P$. Alors, par la propriété précédente, $(X-a_{k+1}) \mid P$.

    On sait qu'il existe un $Q \in \Kset[X]$ tel que $P = (X-a_{k}) \cdot Q$.

    Et on a $P(a_{k+1}) = 0$, donc $Q(a_{k+1}) = 0$, ce qui implique que $a_{k+1}$ est une racine de $Q$, et donc que $(X-a_{k+1}) \mid Q$.
    
    Par hypothèse de récurrence, $(X-a_1)(X-a_2)\cdots(X-a_k) \mid P$. 

    On obtient donc $(X-a_1)(X-a_2)\cdots(X-a_k)(X-a_{k+1}) \mid P$.

    \textbf{Conclusion :} Par le principe de récurrence, le résultat est vrai pour tout $k \in \N^*$.
\end{proof}

\begin{property}
    Soient $P \in \Kset[X]$, $\alpha \in \Kset$ et $s\in \N^*$. Alors $\alpha$ est une racine de multiplicité $s$ de $P$ si et seulement si $\forall k \in [\![0,s-1]\!], P^{(k)}(\alpha) = 0$, et $P^{(s)}(\alpha) \ne 0$.

    Concrètement, une racine de multiplicité $s$ d'un polynôme est une racine qui annule les $s$ premières dérivées du polynôme, mais pas la $s$-ième dérivée du polynôme.
\end{property}

\begin{proof}
    La démonstration utilisera la formule de Taylor.

    On pose $n = \deg(P)$.

    \begin{align*}
        P &= \sum_{k=0}^{n} \frac{P^{(k)}(\alpha)}{k!} (X-\alpha)^k\\
        &= (X-\alpha)^s \cdot \underbracket{\sum_{k=s}^{n} \frac{P^{(k)}(\alpha)}{k!} (X-\alpha)^{k-s}}_{=Q} + \underbracket{\sum_{k=0}^{s-1} \frac{P^{(k)}(\alpha)}{k!} (X-\alpha)^k}_{=R}
    \end{align*}

    On remarque que $Q$ et $R$ sont respectivement le quotient et le reste de la division euclidienne de $P$ par $(X-\alpha)^s$.

    On suppose maintenant que $\alpha$ est une racine de multiplicité $s$ de $P$ dans $\Kset[X]$.

    \begin{align*}
        \alpha \text{ racine de multiplicité } s &\iff \big(((X-\alpha)^s \mid P) \et ((X-\alpha)^{s+1} \nmid P)\big)\\
        &\iff (R=0 \et P^{(s)}(\alpha) \ne 0)\\
        &\iff (R=0 \et Q(\alpha) \ne 0)\\
        &\iff \begin{cases}
        \forall k \in [\![0,s-1]\!], P^{(k)}(\alpha) = 0\\
        P^{(s)}(\alpha) \ne 0
        \end{cases}
    \end{align*}

    On en déduit donc que $\alpha$ est une racine de multiplicité $s$ de $P$ si et seulement si $\forall k \in [\![0,s-1]\!], P^{(k)}(\alpha) = 0$, et $P^{(s)}(\alpha) \ne 0$.
\end{proof}

\begin{property}[Relations racines/coeffs. d'un polynôme scindé]
    Soit $P = \sum_{k=0}^{n}a_kX^k \in \Kset[X]$ un polynôme scindé dans $\Kset[X]$, et $x_1, x_2, \ldots, x_n$ les racines de $P$ (pas forcément distinctes) dans $\Kset$.

    On pose $\forall k \in [\![1,n]\!], \sigma_k = \sum_{1 \le i_1 < \cdots < i_k \le n} x_{i_1} \cdots x_{i_k}$, c'est à dire que $\sigma_k$ est la somme de tous les produits de $k$ racines distinctes de $P$.

    Alors $\forall k \in [\![0,n]\!], \sigma_k = (-1)^k \cdot \frac{a_{n-k}}{a_n}$, où $a_n$ est le coefficient dominant de $P$ (Formule de Viète).
\end{property}

\begin{proof}
    La démonstration se fait par récurrence sur $n = \deg(P) \in \N^*$. Elle est assez complexe et ne sera pas détaillée dans ce cours.
\end{proof}

\begin{theorem}
    Soient $P_1,\dots,P_n \in \Kset[X]$.
    \begin{outline}[enumerate]
        \1 \textbf{PGCD :}
        \2 $\exists ! D \in \Kset[X]$ avec $\begin{cases}D \text{ unitaire }\\\ou\\D= 0\end{cases}$ tel que :\\ $\sum_{i=1}^{n}P_i \Kset[X] = D\cdot\Kset[X]$ et $\forall i \in [\![1,n]\!], D \mid P_i$.
        \2 Si $\exists \Delta \in \Kset[X]$ tel que $\forall i \in [\![1,n]\!], \Delta \mid P_i$, alors $\Delta \mid D$.
        
        $D$ est appelé le \textbf{plus grand commun diviseur} (PGCD) de $P_1,\dots,P_n$, et est noté $\pgcd(P_1,\dots,P_n)$.

        \1 \textbf{PPCM :}
        \2 $\exists ! M \in \Kset[X]$ avec $\begin{cases}M \text{ unitaire }\\\ou\\M= 0\end{cases}$ tel que :\\ $\bigcap_{i=1}^{n} P_i \cdot \Kset[X] = M\cdot\Kset[X]$ et $\forall i \in [\![1,n]\!], P_i \mid M$.
        \2 Si $\exists A \in \Kset[X]$ tel que $\forall i \in [\![1,n]\!], P_i \mid A$, alors $M \mid A$.

        $M$ est appelé le \textbf{plus petit commun multiple} (PPCM) de $P_1,\dots,P_n$, et est noté $\ppcm(P_1,\dots,P_n)$.
    \end{outline}
\end{theorem}

\begin{proof}
    La preuve est analogue à celle du théorème sur le PGCD et le PPCM dans un anneau (chapitre 11).
\end{proof}

\begin{property}
    Soient $P_1, \dots , P_n, Q \in \Kset[X]$.
    \begin{outline}[enumerate]
        \1 $\pgcd(P_1Q, P_2Q, \dots, P_nQ) = \pgcd(P_1, P_2, \dots, P_n) \cdot \frac{Q}{\cfd(Q)}$.
        \1 $\ppcm(P_1Q, P_2Q, \dots, P_nQ) = \ppcm(P_1, P_2, \dots, P_n) \cdot \frac{Q}{\cfd(Q)}$.
    \end{outline}
\end{property}

\begin{property}
    Soient $A,B,C \in \Kset[X] \setminus \{0\}$.

    $$\begin{cases}
    A \wedge B = 1\\
    C \mid B
    \end{cases}
    \implies A \wedge C = 1$$
\end{property}

\begin{theorem}[Égalité de Bézout]
    Soient $P_1,\dots,P_n \in \Kset[X]$.

    $$D = \pgcd(P_1,\dots,P_n) \implies \exists Q_1, \dots, Q_n \in \Kset[X] \tq D = \sum_{i=1}^{n}P_iQ_i$$
\end{theorem}

\begin{property}
    Soient $P_1,\dots,P_n \in \Kset[X] \et D \in \Kset[X]$.

    $$\si \exists U_1,\dots,U_n \in \Kset[X] \tq D = P_1U_1 + P_2U_2 + \cdots + P_nU_n$$

    avec $\forall 1 \le i \le n, D \mid P_i$ et $D$ est unitaire (ou nul).

    Alors $D = \pgcd(P_1,\dots,P_n)$.
\end{property}

\begin{theorem}[Théorème de Gauss]
    Soient $A,B,C \in \Kset[X]\setminus\{0\}$.

    $$\begin{cases}
    A \wedge B = 1\\
    A \mid C\\
    B \mid C
    \end{cases}
    \implies AB \mid C$$
\end{theorem}

\begin{property}
    Si $P_1, \dots, P_n$ sont premiers entre eux deux à deux, alors :

    $$\ppcm(P_1,\dots,P_n) = \frac{1}{\prod_{i=1}^{n} \cfd(P_i)} \prod_{i=1}^{n} P_i$$
\end{property}

\begin{theorem}[Lemme d'Euclide]
    Si $A = B\cdot Q + R$ avec $A,B,Q,R \in \Kset[X]$ et $\deg(R) < \deg(B)$, alors $A\wedge B = B\wedge R$.
\end{theorem}

\begin{theorem}[Théorème de D'Alembert-Gauss]
    Tout polynôme non-constant $P \in \C[X]$ de degré $n \ge 1$ admet au moins une racine dans $\C$.
\end{theorem}

\begin{proof}
    À faire comme DL (ou voir CNC 2007).
\end{proof}

\begin{property}[Conjugaison complexe d'un polynôme]
    Soit $P \in \C[X]$.

    $$P \in \R[X] \iff \forall z \in \C, \ovl{P(z)} = P(\ovl z)$$
\end{property}

\begin{proof}
    Dans le sens direct $(\implies)$, la démonstration est triviale. Si $P = \sum_{k=0}^{n}a_kX^k \in \R[X]$, alors $\ovl{P(z)} = \sum_{k=0}^{n}\ovl{a_k}\ovl{z^k} = \sum_{k=0}^{n}a_k\ovl{z^k} = P(\ovl z)$.

    Dans le sens réciproque $(\impliedby)$, supposons que $\forall z \in \C, \ovl{P(z)} = P(\ovl z)$. On pose $P = \sum_{k=0}^{n}a_kX^k$. Alors $\ovl{P(z)} = \sum_{k=0}^{n}\ovl{a_k}\ovl{z^k}$ et $P(\ovl z) = \sum_{k=0}^{n}a_k\ovl{z^k}$.

    Donc $\sum_{k=0}^{n}(\ovl{a_k} - a_k)\ovl{z^k} = 0$ pour tout $z \in \C$.

    Donc $Q$ admet une infinité de racines, et par conséquent $Q$ est le polynôme nul, c'est à dire que $\forall k \in [\![0,n]\!], \ovl{a_k} - a_k = 0$, ce qui implique que $\forall k \in [\![0,n]\!], a_k \in \R$, et donc $P \in \R[X]$.
\end{proof}

\begin{proof}
    Nous allons utiliser les propriétés liées aux polynômes dérivés et à la conjugaison complexe d'un polynôme.

    On sait que $z$ est une racine de multiplicité $s$ de $P$ dans $\C[X]$.

    Par définition, cela équivaut à dire que $\forall k \in [\![0,s-1]\!], P^{(k)}(z) = 0$ et $P^{(s)}(z) \ne 0$.

    Cela implique que $\forall k \in [\![0,s-1]\!], \ovl{P^{(k)}(z)} = 0$ et $\ovl{P^{(s)}(z)} \ne 0$.

    Donc $\forall k \in [\![0,s-1]\!], P^{(k)}(\ovl z) = 0$ et $P^{(s)}(\ovl z) \ne 0$, ce qui équivaut à dire que $\ovl z$ est une racine de multiplicité $s$ de $P$ dans $\C[X]$.
\end{proof}

\begin{proof}
    Soit $P \in \R[X]$.
    \begin{outline}
        \1 Si $P$ est de degré $1$, alors $P$ est irréductible dans $\R[X]$.
        \1 Si $P$ est de degré $2$...
        \2 Si $\Delta \ge 0$, alors $P$ admet au moins une racine dans $\R$, et par conséquent $P$ est réductible dans $\R[X]$.
        \2 Dans le cas où $\Delta < 0$, on procède par absurde en supposant que $P$ est réductible.
        \3 Supposons que $P$ est réductible dans $\R[X]$. Alors il existe $P_1 \in \R[X]$ un polynôme de degré $1$ tel que $P_1 \mid P$, or $P_1$ admet une racine réelle, et par conséquent $P$ admet une racine réelle, ce qui est absurde, car $\Delta < 0$.
        \3 On en déduit que $P$ est irréductible dans $\R[X]$.
        \1 Si $P$ est de degré $\ge 3$, alors :
        \2 Si $P$ admet un au moins une racine réelle $a$, alors $X-a \mid P$, et par conséquent $P$ est réductible dans $\R[X]$.
        \2 Dans le cas contraire, on sait que $P \in \R[X] \subset \C[X]$, donc d'après le théorème de D'Alembert-Gauss, $P$ admet au moins une racine complexe $z$, cette racine $z$ n'est pas réelle (elle appartient à $\C\setminus\R$). $\ovl z$ est aussi une racine de $P$, donc $X-z \mid P$ et $X-\ovl z \mid P$, et par conséquent $(X-z)(X-\ovl z) \mid P$, et comme $(X-z)(X-\ovl z) = X^2 - 2\Re(z)X + |z|^2 \in \R[X]$, cela implique que $P$ est réductible dans $\R[X]$.
    \end{outline}
\end{proof}

\begin{property}
    Soit $A \in \Kset[X] \tq \deg(A) \ge 1$. Alors $\exists P \in \Kset[X]$ irréductible tel que $P \mid A$.
\end{property}

\begin{proof}
    On pose $E = \{\deg(D) \tq D\mid A \et \deg(D) \ge 1\}$.

    On a $A \mid A$ donc $\deg(A) \in E$ et par conséquent $E$ est non-vide. Donc $E$ admet un minimum $m \in \N^*$, et il existe $P \in \Kset[X]$ tel que $\deg(P) = m$ et $P \mid A$.

    Supposons que $P$ est réductible dans $\Kset[X]$. Alors il existe $Q \in \Kset[X]$ tel que $1 \le \deg(Q) < \deg(P)$ et $Q \mid P$. Or $P \mid A$, alors $Q \mid A$, donc $\deg Q \in E$ et $\deg(Q) < m$, ce qui est absurde, car $m$ est le minimum de $E$.

    On en déduit que $P$ est irréductible dans $\Kset[X]$ et que $P \mid A$.
\end{proof}

\begin{theorem}[Décomposition primaire d'un polynôme]
    Soit $A \in \Kset[X]$. Alors il existe une décomposition primaire de $A$, c'est-à-dire des polynômes irréductibles $P_1, \dots, P_s \in \Kset[X]$ tels que :
    $$A = \cfd(A) P_1 \dots P_s$$

    Cette décomposition est unique à une permutation près.

    i.e. si $A = \cfd(A) \prod_{i=1}^r Q_i$ avec $Q_i$ irréductible, alors $r=s$ et $\{P_1,\dots,P_s\} = \{Q_1,\dots,Q_r\}$.
\end{theorem}

\begin{proof}
    L'existence est démontrée par récurrence forte sur $\deg(A)$.

    \textbf{Démonstration de l'unicité :}

    Supposons que $A = \cfd(A) P_1 \dots P_s$ et $A = \cfd(A) Q_1 \dots Q_r$ avec $P_i, Q_j$ irréductibles.

    On a donc $\forall i \in [\![1,s]\!], P_i \mid A$

    Donc $\exists j \in [\![1,r]\!]$ tel que $P_i \mid Q_j$, car $P_1, \dots, P_s$ sont irréductibles.

    Comme $Q_j$ est irréductible, cela implique que :

    $$\begin{cases}
        P_i \text{ constant non-nul }\\
        \exists \lambda \in \Kset^* \tq P_i = \lambda Q_j
    \end{cases}$$

    Or $\cfd(P_i) = 1$ et $\cfd(Q_j) = 1$, donc $\lambda = 1$, et par conséquent $P_i = Q_j$.
\end{proof}

\begin{theorem}[PGCD et PPCM par décomposition primaire]
    Soient $A,B$ deux polynômes de $\Kset[X]$ de degré $\ge 1$.

    On écrit :

    $$\begin{cases}
    A = \cfd(A) \prod_{i=1}^{r} P_i^{\alpha_i}\\
    B = \cfd(B) \prod_{i=1}^{r} P_i^{\beta_i}
    \end{cases}$$

    avec $P_1,\dots,P_r \in \Kset[X]$ irréductibles deux à deux distincts
    
    et $\alpha_1,\dots,\alpha_r, \beta_1,\dots,\beta_r \in \N$.

    Le PGCD et le PPCM de $A$ et $B$ sont alors donnés par les formules suivantes :

    $$A \wedge B = \prod_{i=1}^{r} P_i^{\min(\alpha_i,\beta_i)}$$
    
    $$A \vee B = \prod_{i=1}^{r} P_i^{\max(\alpha_i,\beta_i)}$$
    
\end{theorem}

\begin{proof}[Démonstration de la caractérisation (2)]
    \begin{outline}
        \1 La démonstration de l'implication $(\implies)$ est assez simple.
        \2 On suppose que $P \wedge Q = 1$. 
        \2 Par conséquent, il existe $A,B \in \Kset[X]$ tels que $1 = AP + BQ$.
        \2 Soit $z \in \C$ une racine de $P$ et de $Q$. Alors $1 = AP(z) + BQ(z) = 0$, ce qui est absurde.
        \2 On en déduit que $P$ et $Q$ n'ont aucune racine commune dans $\C$.
        \1 La démonstration de l'implication $(\impliedby)$ est plus complexe. % ←←← ba dum tss, hahahah complexe
        \2 On suppose que les deux polynômes $P$ et $Q$ n'ont aucune racine commune dans $\C$.
        \2 On pose $D = P \wedge Q$.
        \2 Donc $\exists P_1, Q_1 \in \Kset[X]$ tels que $P = D P_1$ et $Q = D Q_1$.
        \2 Nous allons raisonner par absurde en supposant que $\deg(D) \ge 1$.
        \2 Comme $\deg(D) \ge 1$, alors $D$ admet au moins une racine $z \in \C$.
        \2 Donc $P(z) = D(z) P_1(z) = 0$ et $Q(z) = D(z) Q_1(z) = 0$, ce qui est absurde, car $P$ et $Q$ n'ont aucune racine commune dans $\C$.
        \2 On en déduit que $\deg(D) = 0$, et comme $D$ est unitaire, alors $D = 1$, et par conséquent $P \wedge Q = 1$.
    \end{outline}
\end{proof}

\begin{theorem}
    Soient $f$ une fonction continue sur le segment $[a,b]$ et $a \le x_0 < x_1 < \cdots < x_n \le b$ des éléments de $[a,b]$.

    On souhaite ``approximer'' en quelque sorte la fonction $f$ par un polynôme $P \in \R_n[X]$ tel que $P(x_i) = f(x_i)$ pour tout $i \in [\![0,n]\!]$.

    Alors $P \in \R_n[X]$ est unique et est donné par la formule suivante :

    $$P = \sum_{k=0}^{n} f(x_k) l_k$$

    où $(l_k)_{0 \le k \le n}$ sont les polynômes élémentaires de Lagrange associés à $(x_k)_{0 \le k \le n}$.

    $P$ est appelé le \textbf{polynôme d'interpolation de Lagrange} (ou l'interpolé) de $f$ aux points $x_0,\dots,x_n$.
\end{theorem}

\begin{proof}
    \begin{outline}
        \1 Démontrons que $P$ est un polynôme.
            \2 On sait que $\forall k \in [\![0,n]\!], l_k \in \R_n[X]$, donc $P = \sum_{k=0}^{n}f(x_k) l_k \in \R_n[X]$.
            \2 Et on a bien $P(x_i) = \sum_{k=0}^{n} f(x_k) l_k(x_i) = f(x_i)$ pour tout $i \in [\![0,n]\!]$.
            \2 On a donc construit un polynôme $P$ tel que $P(x_i) = f(x_i)$ pour tout $i \in [\![0,n]\!]$.
        \1 Démontrons que $P$ est unique.
            \2 Supposons qu'il existe un autre polynôme $Q \in \R_n[X]$ tel que $Q(x_i) = f(x_i)$ pour tout $i \in [\![0,n]\!]$.
            \2 Alors $P(x_i) = Q(x_i)$ pour tout $i \in [\![0,n]\!]$.
            \2 Donc $(P-Q)(x_i) = 0$ pour tout $i \in [\![0,n]\!]$.
            \2 Par conséquent, le polynôme $P-Q$ admet au moins $n+1$ racines distinctes, ce qui est absurde, car $\deg(P-Q) \le n$.
            \2 On en déduit que $P-Q$ est le polynôme nul, c'est à dire que $P=Q$, et par conséquent $P$ est unique.
    \end{outline}
\end{proof}

\begin{theorem}[Erreur d'interpolation de Lagrange]
    Soient $f$ une fonction de classe $\cclas{n+1}$ à image dans $\Kset$ sur le segment $[a,b]$ et $a \le x_0 < x_1 < \cdots < x_n \le b$ des éléments de $[a,b]$.

    L'erreur d'interpolation de Lagrange est donnée par la formule suivante :

    $$\forall x \in [a,b], \exists \eps_x \in [a,b], e(x) = f(x) - P(x) = \frac{f^{(n+1)}(\eps_x)}{(n+1)!} \prod_{k=0}^{n} (x-x_k)$$
\end{theorem}

\begin{proof}
    \begin{outline}
    \1 Soit $x \in [a,b]$.

    \1 Si $x = x_k \tq 0 \le k \le n$, alors $e(x_k) = f(x_k) - P(x_k) = 0$, et la formule est vérifiée.

    \2 De même, $\prod_{i=0}^n (x_k - x_i) = 0$ pour tout $k \in [\![0,n]\!]$, donc la formule est vérifiée.

    \1 Dans le cas où $x \ne x_k$ pour tout $k \in [\![0,n]\!]$, on pose la fonction $g : [a,b] \to \Kset$ définie par $g : t \mapsto f(t) - P(t) - e(x) \prod_{k=0}^{n} (t-x_k)$.

    \2 Cette fonction est de classe $\cclas{n+1}$ (car toutes les fonctions qui la composent sont aussi de classe $\cclas{n+1}$).

    \2 $\forall k \in [\![0,n]\!], g(x_k) = f(x_k) - P(x_k) - e(x) \prod_{i=0}^{n} (x_k - x_i) = 0$. La fonction admet donc au moins $n+1$ racines distinctes.

    \2 Comme $x$ a été fixé au début de la preuve, on a $g(x) = f(x) - P(x) - e(x) \prod_{k=0}^{n} (x-x_k) = 0$, donc $g$ admet au moins $n+2$ racines distinctes.

    \2 On peut donc procéder en utilisant le théorème de Rolle (plus précisément, Rolle itéré) pour montrer que $g^{(n+1)}$ admet au moins une racine dans $[a,b]$, sachant que Rolle itéré nécessite $n+2$ racines distinctes pour garantir l'existence d'une racine de $g^{(n+1)}$, car $g$ est de classe $\cclas{n+1}$.

    \2 Donc $\exists \eps_x \in [a,b]$ tel que $g^{(n+1)}(\eps_x) = 0$.

    \2 Or $g^{(n+1)}(t) = f^{(n+1)}(t) - P^{(n+1)}(t) - e(x) \prod_{k=0}^{n} (t-x_k)$. Comme $P$ est de degré $\le n$, alors $P^{(n+1)}$ est le polynôme nul, et par conséquent $g^{(n+1)}(t) = f^{(n+1)}(t) - e(x) \prod_{k=0}^{n} (t-x_k)$.

    \2 On en déduit que $0 = g^{(n+1)}(\eps_x) = f^{(n+1)}(\eps_x) - e(x) \prod_{k=0}^{n} (\eps_x-x_k)$, et par conséquent $e(x) = \frac{f^{(n+1)}(\eps_x)}{(n+1)!} \prod_{k=0}^{n} (x-x_k)$.
    \end{outline}
\end{proof}


\vspace{1em}


\section*{Extraits : 260220 CH13}


\begin{property}
    Soit $A$ un anneau intègre commutatif. On définit une relation d'équivalence sur $A \times A^*$ par :

    $$(N, D) \sim (N', D') \iff ND' = N'D$$

    Cette relation est bien une relation d'équivalence.
\end{property}

\begin{proof}
    \begin{outline}
        \1 Réflexivité : $ND = ND$ pour tout $(N, D) \in A \times A^*$.
        \1 Symétrie : Si $ND' = N'D$, alors $N'D = ND'$, donc $(N', D') \sim (N, D)$ (la relation $=$ est symétrique).
        \1 Transitivité :
        \2 Si $(N, D) \sim (N', D')$ et $(N', D') \sim (N'', D'')$, alors $ND' = N'D$ et $N'D'' = N''D'$.
        \2 En multipliant la première égalité par $D''$ et la seconde par $D$, on obtient $ND'D'' = N'DD''$ et $N'DD'' = N''DD'$.
        \2 En combinant ces égalités, on trouve $ND'D'' = N''DD'$, ce qui implique que $ND'' = N''D$ (car $A$ est intègre).
        \2 Ainsi, $(N, D) \sim (N'', D'')$.
    \end{outline}
\end{proof}

\begin{property}
    $$\begin{cases}
    \cl{N,D} = \cl{N', D'}\\
    \cl{A,B} = \cl{A', B'}
    \end{cases}
    \implies
    \begin{cases}
    \cl{NB + AD, BD} = \cl{N'B' + A'D', B'D'}\\
    \cl{AN,BD} = \cl{A'N', B'D'}
    \end{cases}$$
\end{property}

\begin{proof}
    \begin{outline}
        \1 Montrons que $\cl{NB + AD, BD} = \cl{N'B' + A'D', B'D'}$.
        \2 En utilisant les égalités $\cl{N,D} = \cl{N', D'}$ et $\cl{A,B} = \cl{A', B'}$, on a $ND' = N'D$ et $AB' = A'B$.
        \2 En multipliant la première égalité par $B$ et la seconde par $D$, on obtient $NDB = N'DB$ et $ABD = A'BD$.
        \2 En combinant ces égalités, on trouve $NB + AD = N'B' + A'D'$ et $BD = B'D'$.
        \1 Montrons que $\cl{AN,BD} = \cl{A'N', B'D'}$.
        \1 En utilisant les mêmes égalités, on a $AN = A'N'$ et $BD = B'D'$, ce qui implique que $\cl{AN,BD} = \cl{A'N', B'D'}$.
    \end{outline}
\end{proof}

\begin{property}
    Avec les opérations $+$ et $\times$ définies ci-dessus, $\Kset$ est un corps commutatif.
\end{property}

\begin{proof}
    \begin{outline}
        \1 On remarque que $\frac ND = \frac{ND'}{DD'}$ et $\frac{N'}{D'} = \frac{N'D}{DD'}$.
        \1 On peut donc supposer que les ``fractions'' sont écrites avec un dénominateur commun.
        \1 Les lois $+$ et $\times$ sont déjà définies d'après la remarque précédente.
        \1 Il ne reste plus qu'à vérifier que $\Kset$ est un corps commutatif.
        \2 L'addition est associative et commutative, avec $0_\Kset$ comme élément neutre. On utilise les propriétés de l'addition dans $A$ (car $A$ est un anneau commutatif) pour vérifier ces propriétés dans $\Kset$.
        \2 La multiplication est associative et commutative, avec $1_\Kset$ comme élément neutre. On utilise les propriétés de la multiplication dans $A$ pour vérifier ces propriétés dans $\Kset$.
        \2 La multiplication est distributive par rapport à l'addition. On utilise les propriétés de l'addition et de la multiplication dans $A$ pour vérifier cette propriété dans $\Kset$.
        \2 Tout élément nul de $\Kset$ est de la forme $\frac{0_A}{D}$, ce qui correspond à $0_\Kset$. Tout élément non nul de $\Kset$ est de la forme $\frac{N}{D}$ avec $N \neq 0_A$, et son inverse est $\frac{D}{N}$, qui est bien défini car $A$ est intègre (et on démontre facilement que $\frac{N}{D} \cdot \frac{D}{N} = 1_\Kset$).
    \end{outline}
    On conclut que $\Kset$ est un corps commutatif.
\end{proof}

\begin{property}[Injection canonique de $A$ dans $\Kset$]
    
    $$f : \substack{A \to \Kset\\ N \mapsto \frac{N}{1_A}}$$

    $f$ est un morphisme d'anneaux injectif.
\end{property}

\begin{proof}
    On peut procéder par deux méthodes, soit en montrer que le noyau $\ker(f) = \{0_A\}$ (par équivalences successives sans avoir à passer par une double inclusion), soit en montrant que $f$ est injective directement (en étudiant l'égalité $f(N) = f(N')$ avec $N, N' \in A$).
\end{proof}

\begin{theorem}
    $\Kset$ est le plus petit corps contenant $A$.

    $\Kset$ est appelé \textbf{le corps de fractions} de $A$.
\end{theorem}

\begin{proof}
    Soit $L$ un corps contenant $A$. On souhaite montrer que $\Kset \subseteq L$.

    \begin{outline}
        \1 Soit $\frac ND \in \Kset$ avec $N \in A$ et $D \in A^*$. Comme $L$ est un corps contenant $A$, on a $N \in L$ et $D \in L^*$ (car $D$ est non nul dans $A$, il est aussi non nul dans $L$).
        \1 Par conséquent, $\frac ND = N \cdot D^{-1} \in L$.
        \1 Ainsi, tout élément de $\Kset$ est aussi un élément de $L$, ce qui implique que $\Kset \subseteq L$.
    \end{outline}
\end{proof}

\begin{property}
    Soient $(N,D),(N_1,D_1) \in \Kset[X]\times\Kset[X]^*$.

    Si $\frac ND = \frac{N_1}{D_1}$, alors : $N \cdot D_1 = N_1 \cdot D$.

    On a donc une égalité de deux polynômes, et on peut comparer leurs degrés. En particulier, on a $\deg(N) + \deg(D_1) = \deg(N_1) + \deg(D)$, ce qui implique que :

    $$\deg(N) - \deg(D) = \deg(N_1) - \deg(D_1)$$

    Cela veut dire que $\deg(N) - \deg(D)$ est une quantité qui ne dépend pas du représentant choisi pour cette fraction rationnelle.
\end{property}

\begin{property}
    Soient $F, G \in \Kset(X)$. Alors :

    $$\deg(F\cdot G) = \deg(F) + \deg(G)$$
    $$\deg(F + G) \leq \max(\deg(F), \deg(G)) \text{ avec égalité si } \deg(F) \neq \deg(G)$$
\end{property}

\begin{proof}
    \begin{enumerate}
        \item La première égalité est vérifiée trivialement, en effet :
        
        \begin{align*}
            \deg(F\cdot G) &= \deg\left(\frac{A}{B} \cdot \frac{C}{D}\right)\\
            &= \deg\left(\frac{AC}{BD}\right)\\
            &= \deg(AC) - \deg(BD)\\
            &= (\deg(A) + \deg(C)) - (\deg(B) + \deg(D))\\
            &= (\deg(A) - \deg(B)) + (\deg(C) - \deg(D))\\
            &= \deg(F) + \deg(G)
        \end{align*}
        \item Pour démontrer la seconde inégalité, on pose $F = \frac AD$ et $G = \frac BD$ (on suppose que les fractions sont écrites avec un dénominateur commun, ce qui est toujours possible). Alors : $F + G = \frac{A + B}{D}$, et $\deg(F + G) = \deg(A + B) - \deg(D)$. \\Or, $\deg(A + B) \leq \max(\deg(A), \deg(B))$ (propriété issue des polynômes !), ce qui implique que $\deg(F + G) \leq \max(\deg(A), \deg(B)) - \deg(D) = \max(\underbrace{\deg(A) - \deg(D)}_{\text{degré de fraction}}, \underbrace{\deg(B) - \deg(D)}_{\text{degré de fraction}}) = \max(\deg(F), \deg(G))$.\\Ensuite, on démontre facilement l'égalité dans le cas où $\deg(F) \neq \deg(G)$.
    \end{enumerate}
\end{proof}

\begin{property}
    Soit $F = \frac ND$.

    La fraction $\frac{N'D - ND'}{D^2}$ ne dépend pas du représentant choisi pour $F$.
\end{property}

\begin{proof}
    Soit $F = \frac ND = \frac AB \iff NB = AD$.

    Nous voulons montrer que $\frac{N'D - ND'}{D^2} = \frac{A'B - AB'}{B^2} \iff D^2\parens{A'B-AB'} = B^2\parens{N'D - ND'}$.

    \begin{outline}
        \1 Dans le sens direct $\implies$ :
        \2 On a N'B + N B' = A'D + A D' (en dérivant l'égalité $NB = AD$).
        \2 Et on a aussi $D^2(A'B - AB') = A'D\cdot DB - AD^2 B$ et $B^2(N'D - ND') = N'B\cdot BD - NB^2 D$.
        \2 En combinant ces égalités, on trouve que $D^2(A'B - AB') = B^2(N'D - ND')$.
        \1 Dans le sens inverse $\impliedby$, la démonstration est plus simple, on peut directement utiliser les égalités $D^2(A'B - AB') = B^2(N'D - ND')$ et $NB = AD$ pour trouver que $N'B + N B' = A'D + A D'$.
    \end{outline}
\end{proof}

\begin{property}
    Soient $F_1,F_2 \in \Kset(X)$. Alors :

    \begin{enumerate}
        \item $\forall \lambda \in\Kset, (F_1 + \lambda F_2)' = F_1' + \lambda F_2'$
        \item $(F_1 \cdot F_2)' = F_1' \cdot F_2 + F_1 \cdot F_2'$
    \end{enumerate}
\end{property}

\begin{proof}
    \begin{enumerate}
        \item Pour démontrer la première propriété, on peut utiliser le fait que $\Kset(X)$ est un corps (et techniquement, un sous-espace vectoriel), et considérer $\lambda \in \Kset$ comme un polynôme de degré $0$.
        \item La démonstration est triviale, on procède par calcul direct.
    \end{enumerate}
\end{proof}

\begin{theorem}[Forme irréductible d'une fraction rationnelle]
    Soit $F \in \Kset(X)$. Alors, il existe un unique couple $(A,B) \in \Kset[X]^2$ tel que $A \wedge B = 1$, $B$ est unitaire et $F = \frac AB$.
\end{theorem}

\begin{proof}
    Soit $F = \frac ND \in \Kset(X)$.

    \begin{outline}
        \1 \textbf{Existence} : On peut trouver un couple $(A, B) \in \Kset[X]^2$ tel que $A \wedge B = 1$, $B$ est unitaire et $\frac AB = \frac ND$.
        \2 Montrons que $F$ peut s'écrire sous la forme $\frac AB$ avec $A, B \in \Kset[X]$ et $B$ unitaire.
        \3 On sait que $D$ admet un coefficient dominant non nul, disons $d$, car $D \neq 0_\Kset[X]$.
        \3 On pose $N_1 = d^{-1} \cdot N$ et $D_1 = d^{-1} \cdot D$. Alors, $N_1$ et $D_1$ sont des polynômes de $\Kset[X]$ avec $D_1$ unitaire
        \2 On a alors : $\frac{N_1}{D_1} = \frac ND$.
        \3 On peut se permettre d'écrire $F = \frac{N_1}{D_1}$ car ils agissent comme d'autres représentants de cette classe d'équivalence.
        \1 Montrons qu'on peut écrire $F$ sous la forme $\frac AB$ avec $A, B \in \Kset[X]$, $A \wedge B = 1$, $B$ unitaire et $\cfd(B) = 1$.
        \2 Soit $C = N_1 \wedge D_1$.
        \2 Alors il existe un couple $(A,B) \in \Kset[X]^2$ tel que $N_1 = CA$ et $D_1 = CB$.
        \2 Donc $F = \frac AB \avec A \wedge B = 1 \et \cfd(B) = 1$, sachant que $D_1$ et $C$ sont unitaires.
        \1 \textbf{Unicité} : Supposons que $F$ puisse s'écrire sous la forme $\frac AB$ et $\frac{A'}{B'}$ avec $A, B, A', B' \in \Kset[X]$, $A \wedge B = 1$, $A' \wedge B' = 1$, $B$ unitaire et $B'$ unitaire. Alors, on a :
        \2 $A'B = AB'$ (en utilisant l'égalité $\frac AB = \frac{A'}{B'}$).
        \2 Comme $B \mid A' B$, alors en utilisant le théorème de Gauss, on trouve que $B \mid B'$. De même, $B' \mid A B'$, donc $B' \mid B$ (sachant que $A$ et $B$ sont premiers entre eux).
        \2 Ainsi, $B$ et $B'$ sont deux polynômes qui divisent mutuellement l'un l'autre, ils sont donc associés à un élément unitaire près. Mais comme ils sont déjà unitaires, alors $B = B'$.
    \end{outline}
\end{proof}

\begin{property}[Partie entière et partie fractionnaire]
    Soit $F \in \Kset(X)$. Alors, il existe un unique couple $(E, G) \in \Kset[X] \times \Kset(X)$ tel que :

    $$\begin{cases}
    F = E + G\\
    \deg(G) < 0
    \end{cases}$$

    $E$ est appelé la partie entière de $F$, et $G$ est appelé la partie fractionnaire de $F$.
\end{property}

\begin{proof}
    On pose $F = \frac AB \in \Kset(X)$ avec $A, B \in \Kset[X]$ et $B \neq 0_\Kset[X]$.

    On effectue la division euclidienne de $A$ par $B$, ce qui nous donne un quotient $E \in \Kset[X]$ et un reste $R \in \Kset[X]$ tels que $A = EB + R$ avec $\deg(R) < \deg(B)$.

    Donc $F = \frac{EB + R}{B} = E + \frac RB$.

    L'existence du couple $(E, G)$ est ainsi démontrée en posant $G = \frac RB$. L'unicité est démontrée implicitement, car le quotient et le reste de la division euclidienne sont uniques.
\end{proof}

\begin{property}
    Soient $A$ et $B_1,\dots,B_n \in \Kset[X]$ tels que $B_i$ et $B_j$ sont deux à deux premiers entre eux.

    On pose $F = \frac{A}{B_1 \cdots B_n}$ avec $\deg(F) < 0$.

    Alors...

    $$\exists! (A_1, \dots, A_n) \in \Kset[X]^n \tqp F = \sum_{i=1}^{n}\frac{A_i}{B_i} \et \underbrace{\forall i \in [\![1,n]\!], \deg(A_i) < \deg(B_i)}_{\text{condition importante}}$$
\end{property}

\begin{proof}
    La démonstration de cette propriété se fait par récurrence sur $n \in \N^*$. Le cas de base $n=1$ est trivial, mais le cas $n=2$ est plus utile est plus intéressant à démontrer. Il faut y employer le théorème de Bézout et décomposer le numérateur $A$ en $A\cdot 1$ puis en $A\cdot (B_1 u + B_2 v)$, ce qui permet de trouver les éléments $A_1$ et $A_2$ de la décomposition.

    La preuve a été faite en tant qu'exercice hors-classe. Il ne faut pas oublier de démontrer l'existence (via division euclidienne et en considérant la somme des quotients comme une partie entière de la fraction) et l'unicité (en considérant un autre couple vérifiant les mêmes propriétés, et/ou par absurde) de la décomposition pendant l'initialisation de la récurrence, dans le cas où $n=2$. L'hérédité devient ensuite triviale, en recyclant tout simplement la démonstration du cas de base $n=2$.

    \textit{L'unicité ne peut pas être démontrée par la division euclidienne parce qu'elle dépend des constantes $U$ et $V$ de Bézout, qui ne sont pas uniques.}
\end{proof}

\begin{property}
    Soit $F=\frac{A}{B^n} \avec n \in \N^* \et \deg(F) < 0$.
    $$\text{Alors } \exists ! (A_1, \dots, A_n) \in \Kset[X]^n \tq F = \sum_{k=1}^{n} \frac{A_k}{B^k} \et \forall k \in [\![1,n]\!], \deg(A_k) < \deg(B)$$
\end{property}

\begin{proof}
    La preuve se fait par récurrence sur $n \in \N^*$. Le cas de base $n=1$ est trivial. L'hérédité se fait en utilisant la division euclidienne pour trouver un quotient $Q$ et un reste $R$ tels que $A = QB + A_{n+1}$ avec $\deg(R) < \deg(B)$, ce qui permet d'écrire $F = \frac{Q}{B^n} + \frac{A_{n+1}}{B^{n+1}}$, et de conclure en utilisant l'hypothèse de récurrence sur la fraction $\frac{Q}{B^{n-1}}$.
\end{proof}

\begin{theorem}[D.E.S. d'une fraction rationnelle]
    Soit $F = \frac AB \in \Kset(X)$ et $A\wedge B = 1$ et $B$ unitaire.

    On écrit $B = \prod_{i=1}^r B_i^{\alpha_i}$ avec $\alpha_i \ge 1$ et $(B_i)_{1\le i \le r}$ irréductibles deux à deux distincts.
    
    Alors, il existe un unique $E \in \Kset[X]$ et une unique famille $(A_{i,j})_{\underset{i \in [\![1,r]\!]}{j \in [\![1,\alpha_i]\!]}}$ d'éléments de $\Kset[X]$ tels que :

    $$F = E + \sum_{i=1}^r \sum_{j=1}^{\alpha_i} \frac{A_{i,j}}{B_i^j} \et \forall i \in [\![1,r]\!], \forall j \in [\![1,\alpha_i]\!], \deg(A_{i,j}) < \deg(B_i)$$
\end{theorem}

\begin{proof}
    Nous utilisons d'abord la division euclidienne de $A$ par $B$ pour trouver un quotient $E$ et un reste $R$ tels que $A = EB + R$ avec $\deg(R) < \deg(B)$. Par conséquent, $F = E + \frac RB$, qui reste une fraction de degré négatif. 
    
    Ensuite, on utilise la première propriété pour trouver des éléments $A_{i,1}$ tels que $\frac RB = \sum_{i=1}^r \frac{A_{i,1}}{B_i}$ avec $\deg(A_{i,1}) < \deg(B_i)$ pour tout $i \in [\![1,r]\!]$, sachant que $B$ peut être décomposé en un produit $\prod_{i=1}^r B_i^{\alpha_i}$ avec $\deg(\frac{A_k}{B_k^{\alpha_k}}) < 0$ pour tout $k \in [\![1,r]\!]$.
    $$\frac{R}{\prod_{i=1}^r B_i^{\alpha_i}} = \sum_{i=1}^r \frac{A_{i,1}}{B_i^{\alpha_i}}$$
    Ensuite, on utilise la seconde propriété pour trouver des éléments $A_{i,j}$ tels que $\frac{A_{i,1}}{B_i} = \sum_{j=1}^{\alpha_i} \frac{A_{i,j}}{B_i^j}$ avec $\deg(A_{i,j}) < \deg(B_i)$ pour tout $j \in [\![1,\alpha_i]\!]$ et tout $i \in [\![1,r]\!]$.
    
    En combinant toutes ces égalités, on trouve que $F = E + \sum_{i=1}^r \sum_{j=1}^{\alpha_i} \frac{A_{i,j}}{B_i^j}$ avec les conditions de degrés sur les $A_{i,j}$.
\end{proof}

\begin{property}
    Soit $F = \frac AB \in \Kset(X)$ avec $A\wedge B = 1$ et $B = \prod_{i=1}^n (X - a_i)$ où les $(a_i)_{1\le i \le n}$ sont deux à deux distincts.

    Alors, $F$ peut être décomposée de la manière suivante :

    $$F = E + \sum_{i=1}^{n}\frac{\lambda_i}{X - a_i}$$
    
    avec $E \in \Kset[X]$ et $\lambda_i \in \Kset$ pour tout $i \in [\![1,n]\!]$.

    Alors $\forall k \in \{1,\dots,n\}, \lambda_k = \frac{A(a_k)}{B'(a_k)}$.

    \textit{(On peut se permettre de ``diviser'' par $B'(a_k)$ (en travaillant dans $\R$ ou $\C$) car $B'(a_k) \neq 0$ pour tout $k \in [\![1,n]\!]$, étant donné que les $a_i$ sont deux à deux distincts, ce qui implique que $B$ n'admet pas de racine multiple (ne possède que des racines simples), et donc que $B'(a_k) \neq 0$ pour tout $k \in [\![1,n]\!]$).}
\end{property}

\begin{proof}
    recop
\end{proof}

\begin{property}
    Soit $F = \frac AB \avec A \wedge B = 1$ et $B = (X - a)^2 \cdot Q$ avec $Q(a) \neq 0$ et $a \in \Kset$.

    Alors, $F$ peut être décomposée de la manière suivante :

    $$F = E + \frac{\alpha}{X - a} + \frac{\beta}{(X - a)^2} + \frac{A_1}{Q}$$

    avec $\deg(A_1) < \deg(Q)$, $(\alpha,\beta) \in \Kset^2$.

    $$\text{Et on a } \begin{cases}
    \beta = \frac{A(a)}{Q(a)}\\
    \alpha = \left(\frac AQ\right)'(a)
    \end{cases}$$
\end{property}

\begin{proof}
    \begin{align*}
    F &= \frac{A}{(X - a)^2 \cdot Q}
    \\
      &= E + \frac{\alpha}{X - a} + \frac{\beta}{(X - a)^2} + \frac{A_1}{Q}
    \\
    (X - a)^2 \cdot F &= \frac AQ
    \\
      &= (X-a)^2(E+\frac{A_1}{Q}) + \alpha(X-a)+ \beta
    \\
    \implies [(X-a)^2 \cdot F]_{X=a} &= \frac{A(a)}{Q(a)} = \beta
    \\
    \implies \parens{\frac AQ}' &= (X-a)\big[2(E+\frac{A_1}{Q}) + (X-a)\cdot(E+\frac{A_1}{Q})'\big]+\alpha
    \\
    \donc \parens{\frac AQ}'(a) &= \alpha
    \\
    B = (X-a)^2 \cdot Q &\implies B'' = 2Q + 4(X-a)Q' + (X-a)^2 Q'' \implies B''(a) = 2Q(a) \neq 0
    \\
    \implies \beta = \frac{2A(a)}{B''(a)}
    \end{align*}
\end{proof}


\vspace{1em}


\section*{Extraits : 260304 CH14}


\begin{property}
    Soit $E$ un $\Kset$-ev. Alors :
    \begin{itemize}
        \item $\forall x \in E, \forall \alpha \in \Kset, \alpha \cdot x = 0_E \iff \alpha = 0_\Kset \ou x = 0_E$.
        \item $\forall x \in E, -1_\Kset \cdot x = -x$.
    \end{itemize}
\end{property}

\begin{proof}
    \begin{outline}
        \1 Démonstration de la première propriété :
            \2 Dans le sens indirect $\impliedby$, $0_\Kset \cdot x = (0_\Kset + 0_\Kset) \cdot x = 0_\Kset \cdot x + 0_\Kset \cdot x$, donc $(0_\Kset \cdot x) + (-0_\Kset \cdot x) = (0_\Kset \cdot x + 0_\Kset \cdot x) + (-0_\Kset \cdot x)$, donc $0_E = 0_\Kset \cdot x$.
            \2 On a aussi $\forall \alpha \in \Kset$, $0_E \cdot \alpha = 0_E \implies 0_E = 0_E \cdot \alpha = (0_E + 0_E) \cdot \alpha = 0_E \cdot \alpha + 0_E \cdot \alpha$, donc $0_E = 0_E \cdot \alpha$.
            \2 Dans le sens direct $\implies$, on suppose que $\alpha \cdot x = 0_E$. Si $\alpha \ne 0_\Kset$, alors $\alpha^{-1} \cdot (\alpha \cdot x) = \alpha^{-1} \cdot 0_E$, donc $1_\Kset \cdot x = 0_E$, donc $x = 0_E$.
        \1 Démonstration de la deuxième propriété :
            \2 On a $-1_\Kset \cdot x + x = (-1_\Kset + 1_\Kset) \cdot x = 0_\Kset \cdot x = 0_E$, donc $-1_\Kset \cdot x = -x$.
    \end{outline}
\end{proof}

\begin{property}[Caractérisation d'un sous-espace vectoriel]
    Soit $E$ un $\Kset$-ev, et soit $F$ une partie de $E$.

    Les propriétés suivantes sont équivalentes :
    \begin{enumerate}
        \item $F$ est un sous-ev. de $E$.
        \item $F \neq \emptyset$, et pour tous $x,y\in F$ et $\alpha,\beta \in\Kset$, on a $\alpha x + \beta y \in F$.
        \item $0_E \in F$, et pour tous $x,y\in F$ et $\lambda \in\Kset$, on a $\lambda x + y \in F$ (très pratique pour montrer que $F$ est un sous-ev. de $E$).
    \end{enumerate}
\end{property}

\begin{proof}
    Pour démontrer ces équivalences, on va suivre un raisonnement circulaire classique en montrant que $(1) \implies (2)$, $(2) \implies (3)$, et $(3) \implies (1)$.

    \begin{outline}
        \1 Montrons que $(1) \implies (2)$ :
            \2 On a $F$ est un $\Kset$-ev., alors $(F,+)$ est un groupe abélien, donc $F \neq \empty$.
            \2 Donc $\forall \alpha,\beta \in \Kset,\forall x,y \in F, \begin{cases}
            \alpha x \in F\\
            \beta y \in F
            \end{cases}$
            \2 Donc $\alpha x + \beta y \in F$.
        \1 Montrons que $(2) \implies (3)$ :
            \2 On sait que $F \neq \emptyset$.
            \2 Donc pour $\alpha = \beta = 0_\Kset$, et $y = x \in F$, on a $0_E = 0_\Kset \cdot x + 0_\Kset \cdot x \in F$.
            \2 Donc $\forall \alpha \beta \text{ qu'on fixe sur } (1_\Kset,\lambda) \in\Kset^2, \forall x,y \in F, \lambda x + y \in F$.
        \1 Montrons que $(3) \implies (1)$ :
            \2 $0_E \in F$, et pour tous $x,y \in F, x-y \in F$ (en prenant $\lambda = -1_\Kset$).
            \2 Donc $F$ est un sous-groupe de $(E,+)$.
            \2 Donc $(F,+)$ est un groupe abélien.
            \2 On a pour $x = 0_E$, et pour tous $\lambda, y \in \Kset\times F, 0_E + \lambda \cdot y \in F$, donc $\lambda \cdot y \in F$.
            \2 Donc $\cdot$ est une loi de composition externe de $F$.
            \2 Les propriétés de la loi de composition externe sont vérifiées dans $F$ car elles sont vérifiées dans $E$.
            \2 Donc $F$ est un $\Kset$-ev.
    \end{outline}
\end{proof}

\begin{property}
    Soit $E$ un $\Kset$-ev.
    \begin{enumerate}
        \item Si $(F_i)_{i\in I}$ est une famille de sous-ev. de $E$, alors $\bigcap_{i\in I} F_i$ est un sous-ev. de $E$.
        \item Soient $F,G$ deux sous-ev. de $E$. Alors $F \cup G$ est un sous-ev. de $E$ si et seulement si $F \subset G$ ou $G \subset F$.
    \end{enumerate}
\end{property}

\begin{proof}
    \begin{outline}
        \1 Démonstration de la première propriété...
            \2 \textbf{Inclusion : } $\bigcap_{i \in I} F_i \subset E$.
            \2 On sait aussi que $\forall i \in I, 0_E \in F_i \implies 0_E \in \bigcap_{i\in I} F_i$.
            \2 Soient $x,y \in \bigcap_{i\in I} F_i$, et $\lambda \in \Kset$. Alors $\forall i \in I, x,y \in F_i$, donc $\lambda x + y \in F_i$, donc $\lambda x + y \in \bigcap_{i\in I} F_i$, sachant que $F_i$ est un sous-ev. de $E$.
            \2 On en déduit que $\bigcap_{i\in I} F_i$ est un sous-ev. de $E$.
        \1 Démonstration de la deuxième propriété...
            \2 Dans le sens indirect $\impliedby$, si $F \subset G$, alors $F \cup G = G$ est un sous-ev. de $E$. De même, si $G \subset F$, alors $F \cup G = F$ est un sous-ev. de $E$.
            \2 Dans le sens direct $\implies$, on utilise la définition d'un sous-espace vectoriel.
            \2 En effet, comme $F \cup G$ est un sous-ev. de $E$, alors $F \cup G$ est un sous-groupe de $(E,+)$.
            \2 D'après les propriétés des groupes, on a bien $F \subset G$ ou $G \subset F$.
    \end{outline}
\end{proof}

\begin{property}
    Il s'agit du plus petit sous-ev. de $E$ contenant $A$.
\end{property}

\begin{property}
    Soient $E$ un $\Kset$-ev. et $A$,$B$ deux parties de $E$. Alors :
    \begin{enumerate}
        \item $\Vect(A) = A \iff A \text{ est un sous-ev. de } E$.
        \item $\Vect(\Vect(A)) = \Vect(A)$.
        \item $A \subset B \implies \Vect(A) \subset \Vect(B)$.
        \item $\Vect(A) = \Vect(B) \iff \begin{cases}
        A \subset \Vect(B)\\
        B \subset \Vect(A)
        \end{cases}$.
    \end{enumerate}
\end{property}

\begin{proof}
    \begin{outline}
        \1 Démonstration de la quatrième propriété.
            \2 Dans le sens direct $(\implies)$, $A \subset \Vect(A) = \Vect(B)$, et $B \subset \Vect(B) = \Vect(A)$.
            \2 Dans le sens indirect $(\impliedby)$
                \3 $A \subset \Vect(B) \overset{\text{(3)}}\implies \Vect(A) \subset \Vect(\Vect(B)) \overset{\text{(2)}}\implies \Vect(A) \subset \Vect(B)$,
                \3 et $B \subset \Vect(A) \overset{\text{(3),(2)}}\implies \Vect(B) \subset \Vect(A)$, donc $\Vect(A) = \Vect(B)$.
    \end{outline}
\end{proof}

\begin{property}
    Toute surfamille\index{surfamille} d'une famille génératrice est une famille génératrice.
    
    i.e. si $(e_i)_{i \in I}$ est une famille génératrice de $E$, et $J$ est un sous-ensemble de $I$ tel que $I \subset I'$, alors $(e_i)_{i \in I'}$ est une famille génératrice de $E$.
\end{property}

\begin{proof}
    On a $\forall x \in E, \exists J \subset I \text{ fini}, \exists (\lambda_j)_{j \in J} \in \Kset^\card{J} \tq x = \sum_{j \in J} \lambda_j e_j$, or $I \subset I'$, donc $J$ est aussi un sous-ensemble de $I'$, donc $(e_i)_{i \in I'}$ est une famille génératrice de $E$.
\end{proof}

\begin{property}
    Toute sous-famille d'une famille libre est une famille libre.

    i.e. si $(e_i)_{i \in I}$ est une famille libre de $E$, et $I' \subset I$, alors $(e_i)_{i \in I'}$ est une famille libre de $E$.
\end{property}

\begin{proof}
    Supposons que $(e_i)_{i \in I}$ est libre et soit $I' \subset I$.
    
    Soit $J \subset I'$ un sous-ensemble fini de $I'$, alors $J \subset I$ et comme $(e_i)_{i \in I}$ est libre, alors $(e_j)_{j \in J}$ est libre, donc $(e_i)_{i \in I'}$ est libre.
\end{proof}

\begin{theorem}
    Soient $e_1,\dots,e_{n+1}$ des éléments de $E$, un $\Kset$-ev.

    $(e_1,\dots,e_n,e_{n+1})$ est une famille libre de $E$ si et seulement si :
    
    $$(e_1,\dots,e_n) \text{ est une famille libre de } E \et e_{n+1} \notin \Vect(e_1,\dots,e_n)$$
    \label{thm:famille_libre_ajout_element}
\end{theorem}

\begin{proof}
    \begin{outline}
        \1 \textbf{Démonstration dans le sens direct :}
            \2 On a $(e_1,\dots,e_n)$ est une sous-famille de $(e_1,\dots,e_n,e_{n+1})$, donc $(e_1,\dots,e_n)$ est une famille libre de $E$.
            \2 Si $e_{n+1} \in \Vect(e_1,\dots,e_n)$, alors $\exists \lambda_1,\dots,\lambda_n \in \Kset \tq e_{n+1} = \sum_{k=1}^n \lambda_k e_k$, donc $\sum_{k=1}^n (-\lambda_k) e_k + 1_\Kset e_{n+1} = 0$, ce qui est absurde, donc $e_{n+1} \notin \Vect(e_1,\dots,e_n)$.
        \1 \textbf{Démonstration dans le sens indirect :}
            \2 Soient $\lambda_1,\dots,\lambda_{n+1} \in \Kset$ tels que $\sum_{k=1}^{n+1} \lambda_k e_k = 0$.

            \2 Si $\lambda_{n+1} \ne 0$, alors $e_{n+1} = -\lambda_{n+1}^{-1} \sum_{k=1}^n \lambda_k e_k \in \Vect(e_1,\dots,e_n)$, ce qui est absurde, donc $\lambda_{n+1} = 0$.

            \2 Alors $\lambda_{n+1} = 0$, donc $\sum_{k=1}^n \lambda_k e_k = 0$, et comme $(e_1,\dots,e_n)$ est une famille libre de $E$, alors $\forall k \in \{1,\dots,n\}, \lambda_k = 0$.

            \2 Donc $\forall k \in \{1,\dots,n+1\}, \lambda_k = 0$, donc $(e_1,\dots,e_n,e_{n+1})$ est une famille libre de $E$.
    \end{outline}
\end{proof}

\begin{property}
    Soit $E$ un $\Kset$-ev.
    \begin{enumerate}
        \item $B = (e_1,\dots,e_n)$ est une base de $E$ si et seulement si : 
        $$\forall x \in E, \exists ! (\lambda_1,\dots,\lambda_n) \in \Kset^n \tq x = \sum_{k=1}^n \lambda_k e_k$$
        \item $B$ est une base de $E$ si et seulement si : 
        $$\forall x \in E, \exists n\in \N^*, \exists ! (\lambda_1,\dots,\lambda_n) \in \Kset^n \tq x= \sum_{i=1}^{n}\lambda_i e_i$$
    \end{enumerate}
\end{property}

\begin{proof}
    \begin{outline}
        \1 \textbf{Démonstration de la première propriété :}
            \2 $\implies$ : Soit $B$ une famille génératrice de $E$. On a donc $\forall x \in E, \exists (\lambda_1,\dots,\lambda_n) \in \Kset^n \tq x = \sum_{k=1}^n \lambda_k e_k$. Si $\exists \beta_1,\dots,\beta_n \in \Kset$ tels que $x = \sum_{k=1}^n \beta_k e_k$, alors $\sum_{k=1}^n (\lambda_k - \beta_k) e_k = 0$, donc $\forall k \in \{1,\dots,n\}, \lambda_k - \beta_k = 0$, donc $\lambda_k = \beta_k$ (sachant que $B$ est une famille libre).
            \2 $\impliedby$ : On a $\forall x \in E$, $\exists (\lambda_1,\dots,\lambda_n) \in \Kset^n \tq x = \sum_{k=1}^n \lambda_k e_k$, donc $B$ est une famille génératrice de $E$. De plus, si $\sum_{k=1}^n \lambda_k e_k = 0$, alors $\sum_{k=1}^n \lambda_k e_k = \sum_{k=1}^n 0 \cdot e_k = 0$, (par unicité des scalaires) donc $\forall k \in \{1,\dots,n\}, \lambda_k = 0$, donc $B$ est une famille libre.
        \1 \textbf{Démonstration de la deuxième propriété :}
            
    \end{outline}
\end{proof}

\begin{theorem}
    Soient $n \in \N^*$, $E$ un $\Kset$-ev. et $e_1,\dots,e_n$ des éléments de $E$.

    Si $v_1,\dots,v_{n+1} \in \Vect(e_1,\dots,e_n)$ alors $(v_1,\dots,v_{n+1})$ est liée.
\end{theorem}

\begin{proof}
    La démonstration se fait par récurrence sur $n$.

    \textbf{Initialisation :} $n=1$

    Donc $\exists \lambda_1,\lambda_2 \in \Kset$ tels que : $\displaystyle\begin{cases}
    v_1 = \lambda_1 e_1\\
    v_2 = \lambda_2 e_1
    \end{cases}$

    Si $\lambda_1 = 0$, alors $v_1 = 0$, donc $(v_1,v_2)$ est liée.

    Si $\lambda_1 \ne 0$, alors $v_1 = \lambda_1 e_1$, donc $e_1 = \lambda_1^{-1} v_1$, donc $v_2 = \lambda_2 e_1 = \lambda_2 \lambda_1^{-1} v_1$, donc $(v_1,v_2)$ est liée.

    \textbf{Hérédité :} Supposons que la propriété est vérifiée pour $n-1 \in \N^*$, et montrons qu'elle est vérifiée pour $n$.

    \begin{outline}
        \1 Soient $v_1,\dots,v_{n+1} \in \Vect(e_1,\dots,e_n)$.
    
        \1 Donc $\forall j \in \{1,\dots,n+1\}, v_j \in \Vect(e_1,\dots,e_n)$.

        \1 Donc $\forall j \in \{1,\dots,n+1\}, \exists (\lambda_{i,j})_{1 \le i \le n} \in \Kset^n \tqp v_j = \sum_{i=1}^n \lambda_{i,j} e_i$, i.e. :

        $$\begin{cases}
v_1 = \lambda_{1,1} e_1 + \dots + \lambda_{n,1} e_n\\
v_2 = \lambda_{1,2} e_1 + \dots + \lambda_{n,2} e_n\\
\vdots\\
v_{n+1} = \lambda_{1,n+1} e_1 + \dots + \lambda_{n,n+1} e_n (*)
        \end{cases}$$

        \1 Si $\forall k \in \{1,2,\dots,n+1\}, \lambda_{n,k} = 0$, alors $v_1,\dots,v_{n+1} \in \Vect(e_1,\dots,e_{n-1})$, donc d'après l'hypothèse de récurrence, $(v_1,\dots,v_{n+1})$ est liée.

        \1 Sinon, il existe $j_0 \in \{1,2,\dots,n+1\}$ tel que $\lambda_{n,j_0} \ne 0$.

        \1 Comme les $v_j$ jouent des rôles symétriques, on peut prendre $j_0 = n+1$ sans perte de généralité, donc $\lambda_{n,n+1} \ne 0$.

        \1 Alors, d'après l'équation $(*)$, on a $e_n = \lambda_{n,n+1}^{-1} \left(v_{n+1} - \sum_{i=1}^{n-1} \lambda_{i,n+1} e_i\right)$.

        \1 Donc, pour tout $j \in \{1,2,\dots,n\}$, on a :
        $$\begin{aligned}
v_j &= \sum_{i=1}^n \lambda_{i,j} e_i\\
&= \sum_{i=1}^{n-1} \lambda_{i,j} e_i + \lambda_{n,j} e_n\\
&= \left(\sum_{i=1}^{n-1} \lambda_{i,j} e_i\right) + \lambda_{n,j} \lambda_{n,n+1}^{-1} \left(v_{n+1} - \sum_{i=1}^{n-1} \lambda_{i,n+1} e_i\right)\\
\end{aligned}$$

    \1 On pose $w_j = v_j - \lambda_{n,j} \lambda_{n,n+1}^{-1} v_{n+1}$.

    \1 On remarque que $(w_1,\dots,w_n) \in \Vect(e_1,\dots,e_{n-1})$.

    \1 Donc, d'après l'hypothèse de récurrence, la famille $(w_1,\dots,w_n)$ est liée.

    \1 Donc $\exists (\alpha_1,\dots,\alpha_n) \in \Kset^n\setminus\{0_\Kset^n\} \tq \sum_{j=1}^n \alpha_j w_j = 0$.

    \1 Donc $\sum_{j=1}^n \alpha_j v_j + \left(-\sum_{j=1}^n \alpha_j\lambda_{n,j} \lambda_{n,n+1}^{-1}\right) = \alpha_{n+1} v_{n+1}$, avec $\alpha_{n+1} = -\sum_{j=1}^n \alpha_j\lambda_{n,j} \lambda_{n,n+1}^{-1}$.

    \1 Donc $\sum_{j=1}^{n+1} \alpha_j v_j = 0$, avec $\alpha_{n+1} = -\sum_{j=1}^n \alpha_j\lambda_{n,j} \lambda_{n,n+1}^{-1}$.

    \1 Alors $(\alpha_1,\dots,\alpha_{n+1}) \in \Kset^{n+1}\setminus\{0_\Kset^{n+1}\}$.

    \1 Donc $(v_1,\dots,v_{n+1})$ est liée. La propriété est vérifiée pour $n$, d'où la propriété est vérifiée pour tout $n \in \N^*$.
    \end{outline}
\end{proof}

\begin{proof}
    Supposons que $n > m$.

    Donc $e_1,\dots,e_n \in \Vect(g_1,\dots,g_m)$, donc $(e_1,\dots,e_n)$ est liée, ce qui est absurde, donc $n \le m$.
\end{proof}

\begin{property}
    Soient $E$ un $\Kset$-ev. de dimension finie et $G = \{g_1,\dots,g_m\}$ une génératrice de $E$.

    Si $L_0$ est une famille libre telle que $L_0 \subset G$, alors il existe $B$ une base de $E$ telle que $L_0 \subset B \subset G$.
    \label{prop:base_incomplete_famille_libre_subset_generatrice}
\end{property}

\begin{proof}
    On pose $I = \{\card{L} \tq L \text{ famille libre de } E, L_0 \subset L \subset G\}$.

    \begin{outline}
        \1 On a $\card{L_0} \in I$, et $\forall L$ libre dans $E$, $\card{L} \le \card G = m$.

        \1 Donc $I$ est un ensemble non-vide de $\N$, et majoré par $m$, donc $I$ admet une borne supérieure, que nous allons noter $\max(I) = n$.

        \1 D'où $\exists B$ libre dans $E$ et $L_0 \subset B \subset G$ tel que $\card{B} = n$.

        \1 Montrons que $B$ est génératrice de $E$, i.e. $\Vect(B) = E = \Vect(g_1,\dots,g_m)$.
            \2 Il faut donc montrer que
            $$\begin{cases}
            B \subset \Vect(g_1,\dots,g_m)\\
            g_1,\dots,g_m \in \Vect(B)
            \end{cases}$$
            \2 On sait déjà que $B \subset \Vect(g_1,\dots,g_m)$, car $B \subset G$ et $G$ est une famille génératrice de $E$.
            \2 Il faut maintenant montrer que $\{g_1,\dots,g_m\} \in \Vect(B)$.
            \3 Nous allons procéder par absurde en supposant qu'il existe un $k$ dans $\{1,\dots,m\}$ tel que $g_k \notin \Vect(B)$.
            \3 Alors, comme $B$ est libre, alors d'après le théorème précédent (Th. \ref{thm:famille_libre_ajout_element}), la famille $B' = B\cup \{g_k\}$ est libre, avec $L_0 \subset B \subset B\cup\{g_k\} \subset G$.
            \3 Donc $\underbrace{\card{B\cup\{g_k\}}}_{=n+1} \in I$, ce qui est absurde.
        \2 D'où $\forall i \in \{1,\dots,m\}, g_i \in \Vect(B)$, donc $\Vect(B) = E \implies B$ est génératrice de $E$.
        \2 On en conclut que $B$ est une base de $E$ avec $L_0 \subset B \subset G$.
    \end{outline}
\end{proof}

\begin{proof}
    Soit $G = \{g_1,\dots,g_m\}$ une famille génératrice de $E$.

    La famille $G$ peut contenir des $g_k$ nuls. Mais, comme $E$ est non-nul (et donc non réduit au singleton $\{0_E\}$), alors il existe au moins un $g_k$ non nul, que nous allons noter $g_{k_0}$.

    Alors $\{g_k\}$ est libre avec $\{g_k\} \subset G$, donc d'après la propriété précédente, il existe une base $B$ de $E$ telle que $\{g_{k_0}\} \subset B \subset G$.
\end{proof}

\begin{property}
    Soient $E$ un $\Kset$-ev. de dimension finie et $G$ une famille génératrice de $E$.

    Si $L_0$ est une famille libre de $E$ alors il existe forcément $B$ une base de $E$ telle que $L_0 \subset B \subset L_0\cup G$.
    \label{prop:base_incomplete_famille_libre}
\end{property}

\begin{proof}
    $G$ est une famille génératrice de $E$, alors $L_0 \cup G$ est aussi une famille génératrice de $E$, avec $L_0 \subset L_0 \cup G$. On réutilise alors la propriété précédente pour conclure qu'il existe une base $B$ de $E$ telle que $L_0 \subset B \subset L_0 \cup G$.
\end{proof}

\begin{theorem}[Théorème de la base incomplète/extraite]
    Soit $E$ un $\Kset$-ev. de dimension finie non réduit au singleton $\{0_E\}$.
    \begin{enumerate}
        \item Si $L$ est une famille libre de $E$, alors on peut la ``compléter'' en une base de $E$.\\Autrement dit, il existe $B$ une base de $E$ telle que $L \subset B$.
        \item Si $G$ est une famille génératrice de $E$, alors on peut en extraire une base de $E$.\\Autrement dit, il existe $B$ une base de $E$ telle que $B \subset G$.
    \end{enumerate}
    \label{thm:base_incomplete}
\end{theorem}

\begin{proof}
    \begin{outline}
        \1 \textbf{Démonstration de la première partie :}
            \2 $G$ est une famille génératrice de $E$, alors d'après la propriété \ref{prop:base_incomplete_famille_libre}, il existe une base $B$ de $E$ telle que $L \subset B \subset L \cup G$.
        \1 \textbf{Démonstration de la deuxième partie :}
            \2 $L = \emptyset$ est une famille libre de $E$, alors d'après la propriété \ref{prop:base_incomplete_famille_libre_subset_generatrice} et la remarque \ref{rem:ev_dim_finie_admet_base}, il existe une base $B$ de $E$ telle que $\emptyset \subset B \subset G$, donc $B \subset G$.
    \end{outline}
\end{proof}

\begin{theorem}
    Toutes les bases de $E$ ont le même cardinal, donc la dimension de $E$ est bien définie.
\end{theorem}

\begin{proof}
    On pose arbitrairement $B_1$ et $B_2$ deux bases de $E$.
    \begin{outline}
        \1 On a $B_1$ est libre et génératrice de $E$, et $B_2$ est une base de $E$, alors :
        $$\card{B_1} \le \card{B_2} \et \card{B_2} \le \card{B_1}$$
        \1 Donc $\card{B_1} = \card{B_2}$.
    \end{outline}
    On en déduit que toutes les bases de $E$ ont le même cardinal, donc la dimension de $E$ est bien définie.
\end{proof}

\begin{property}
    Soit $E$ un $\Kset$-ev. de dimension finie avec $\underset\Kset\dim E = n \in \N^*$.

    \begin{outline}[enumerate]
        \1 $\begin{cases}
        L \text{ est une famille libre de } E \\
        \card L = n
        \end{cases} \iff L \text{ est une base de } E$
        \1 $\begin{cases}
        G \text{ est une famille génératrice de } E \\
        \card G = n
        \end{cases} \iff G \text{ est une base de } E$
    \end{outline}

    \label{prop:base_iff_libre_card_dim}
\end{property}

\begin{proof}
    \begin{itemize}
        \item ($\impliedby$) Dans le sens indirect, le résultat est clair (si une famille est une base, alors elle est génératrice et dans un espace vectoriel de dimension $n$, elle est de cardinal $n$).
        \item ($\implies$) Dans le sens direct, on applique le théorème de la base incomplète (Th. \ref{thm:base_incomplete}) à la famille libre $L$ et à la famille génératrice $G$ pour construire une base $B$ de $E$ telle que $L \subset B \subset G$. Comme $\card L = n$ et $\card B = n$, alors $L = B$, donc $L$ est une base de $E$. De même, comme $\card G = n$ et $\card B = n$, alors $G = B$, donc $G$ est une base de $E$.
    \end{itemize}
\end{proof}

\begin{property}
    Soient $E$ un $\Kset$-ev. de dimension finie et non réduit au singleton $\{0_E\}$, et $F$ un sous-ev. de $E$.

    Alors $F$ est un $\Kset$-ev. de dimension finie, et $\underset\Kset\dim F \le \underset\Kset\dim E$.
\end{property}

\begin{proof}
    On pose $n = \underset\Kset\dim E$ et $B = \{\eps_1,\dots,\eps_n\}$ une base de $E$.

    \begin{outline}
        \1 Si $F = \{0_E\}$, alors $\underset\Kset\dim F = 0 \le n$.
        \1 Si $F \ne \{0_E\}$, alors on considère :
        $$I = \{\card L \tq L \text{ famille libre de } F\}$$
        \2 On a $F \ne \{0_E\}$, alors il existe $x \in F\setminus\{0_E\}$, alors $\{x\}$ est une famille libre de $F$, donc $\card{\{x\}} = 1 \in I$.
        \2 De plus, $\forall L$ famille libre de $F$, $L$ est aussi une famille libre de $E$, alors $\card L \le n$, donc $I$ est majoré par $n$.
        \2 Et comme $I \subset N$, alors $I$ admet une borne supérieure, que nous allons noter $\max(I) = d$.
        \2 Soit $B_F$ une famille libre de $F$ telle que $\card{B_F} = d$. Montrons que $B_F$ est une base de $F$, i.e. $F = \Vect(B_F)$.
            \3 On a $\Vect(B_F) \subset F$, car $B_F \subset F$ et $F$ est un sev. de $E$.
            \3 Supposons qu'il existe $y \in F$ tel que $y \notin \Vect(B_F)$.
            \3 Cela impliquerait que $B_F \cup \{y\}$ est une famille libre de $F$, alors $\card{B_F \cup \{y\}} = d+1 \in I$, ce qui est absurde, donc $\Vect(B_F) = F$.
            \3 Donc $\forall y \in F, y \in \Vect(B_F)$, donc $F = \Vect(B_F)$, et $B_F$ est une base de $F$.
            \2 Donc $F$ est un $\Kset$-ev. de dimension finie.
            \2 Et comme $B_F$ est une base de $F$ alors $\underset\Kset\dim F = \card{B_F} = d \le n = \underset\Kset\dim E$.
            \2 D'où $F$ est un $\Kset$-ev. de dimension finie, et $\underset\Kset\dim F \le \underset\Kset\dim E$.
    \end{outline}
\end{proof}

\begin{property}
    Soit $E$ un $\Kset$-ev. et $F,G$ deux sous-ev. de $E$ tels que $G$ est de dimension finie.

    Si $F \subset G$, alors $F$ est de dimension finie et $\underset\Kset\dim F \le \underset\Kset\dim G$...
    
    \textit{avec égalité si et seulement si $F = G$.}
\end{property}

\begin{proof}
    $F$ et $G$ sont deux $\Kset$-ev, avec $F \subset G$ et $G$ de dimension finie ($\underset\Kset\dim G < \infty$).

    Donc $F$ est de dimension finie, avec $\underset\Kset\dim F \le \underset\Kset\dim G$.

    \textbf{Démonstration du cas de l'égalité :}
    \begin{outline}
        \1 Dans le sens indirect ($\impliedby$), si $F = G$, alors $\underset\Kset\dim F = \underset\Kset\dim G$.
        \1 Dans le sens direct ($\implies$), on pose $d = \underset\Kset\dim F = \underset\Kset\dim G$.
        \2 Soit $B = (e_1,\dots,e_d)$ une base de $F$.
        \2 Comme $F \subset G$, alors $B$ est une famille libre de $G$.
        \2 Or $\card B = d = \underset\Kset\dim G$, alors $B$ est une base de $G$.
        \2 Donc $F = \Vect(B) = G$, donc $F = G$.
    \end{outline}
\end{proof}

\begin{property}
    Soit $E$ un $\Kset$-ev.

    $E$ est de dimension infinie si et seulement si $E$ admet une famille libre infinie.
\end{property}

\begin{proof}
    \begin{outline}
        \1 $\impliedby$ : Si $\underset\Kset\dim E < \infty$ ($E$ est de dimension finie), alors $E$ admet une famille génératrice finie de cardinal $d$. Donc $\forall L$ famille libre de $E$, $\card L \le d$, donc $E$ n'admet pas de famille libre infinie, ce qui est absurde, donc $E$ est de dimension infinie.
        \1 $\implies$ : Supposons que $E$ est de dimension infinie.
            \2 Comme $E \ne \{0_E\}$, alors il existe $e_1 \in E\setminus\{0_E\}$, alors $\{e_1\}$ est une famille libre de $E$.
            \2 On a aussi $E \ne \Vect(e_1)$ car $E$ est de dimension infinie et il ne peut pas avoir de famille génératrice finie, alors il existe $e_2 \in E\setminus\Vect(e_1)$, alors $\{e_1,e_2\}$ est une famille libre de $E$.
            \2 On a aussi $E \ne \Vect(e_1,e_2)$, et on continue ainsi de suite pour construire une famille libre infinie de $E$.
            \2 Pour démontrer cela rigoureusement, nous allons procéder par récurrence pour construire une famille libre de cardinal $n$ pour tout $n \in \N^*$, et ainsi construire une famille libre infinie de $E$.
                \3 \textbf{Initialisation :} $n=1$, on a déjà construit une famille libre de cardinal $1$ de $E$, à savoir $\{e_1\}$.
                \3 \textbf{Hérédité :} Supposons que nous avons construit une famille libre de cardinal $n$, à savoir $\{e_1,\dots,e_n\}$, alors comme $E$ est de dimension infinie, alors $E \ne \Vect(e_1,\dots,e_n)$, alors il existe $e_{n+1} \in E\setminus\Vect(e_1,\dots,e_n)$, alors $\{e_1,\dots,e_n,e_{n+1}\}$ est une famille libre de cardinal $n+1$ de $E$.
            \2 Donc $\forall n \in \N^*$, il existe une famille libre de cardinal $n$ de $E$, alors il existe une famille libre infinie de $E$.
        \1 Donc $E$ est de dimension infinie si et seulement si $E$ admet une famille libre infinie.
    \end{outline}
\end{proof}

\begin{property}
    Si $E$ et $F$ sont deux $\Kset$-ev. de dimension finie, alors $E\times F$ est aussi de dimension finie, et $\underset\Kset\dim (E\times F) = \underset\Kset\dim E + \underset\Kset\dim F$.
\end{property}

\begin{proof}
    Soient $n = \underset\Kset\dim E$ et $m = \underset\Kset\dim F$, et $B_E = (e_1,\dots,e_n)$ une base de $E$, et $B_F = (f_1,\dots,f_m)$ une base de $F$.

    Donc $\forall (x,y) \in E\times F, \exists ! \lambda_1,\dots,\lambda_n \in \Kset, \exists ! \mu_1,\dots,\mu_m \in \Kset$ tels que :
    
    \begin{align*}
        (x,y) = (x,0_F) + (0_E,y) = (\sum_{i=1}^n \lambda_i e_i,0_F) + (0_E, \sum_{i=1}^m \mu_i f_i) = (\sum_{i=1}^n \lambda_i e_i, \sum_{i=1}^m \mu_i f_i)
    \end{align*}
    
    Donc $B_{E\times F} = (e_i,0_F)_{1\le i\le n} \cup (0_E,f_j)_{1\le j\le m}$ est une base de $E\times F$.

    On en déduit que $E\times F$ est de dimension finie,
    
    et $\underset\Kset\dim (E\times F) = \card{B_{E\times F}} = n + m = \underset\Kset\dim E + \underset\Kset\dim F$.
\end{proof}

\begin{property}
    Soient $V_1,\dots,V_p \in E$.
    
    $$\rang(V_1,\dots,V_p) \le p$$
    
    avec égalité si et seulement si $(V_1,\dots,V_p)$ est une famille libre de $E$,
    
    i.e. $\rang(V_1,\dots,V_p) = p \iff (V_1,\dots,V_p) \text{ est une famille libre de } E$.

    \label{prop:rang_famille_libre_iff_rang_egal_card}
\end{property}

\begin{proof}
    On pose $F = \Vect(V_1,\dots,V_p)$, alors $\underset\Kset\dim F = \rang(V_1,\dots,V_p) = r$.

    Donc $r = \underset\Kset\dim F \le \underset\Kset\dim \Vect(V_1,\dots,V_p) \le \card{V_1,\dots,V_p} = p$.

    \textbf{Démonstration dans le sens indirect : }

    Si $(V_1,\dots,V_p)$ est une famille libre, alors $\Vect(V_1,\dots,V_p)$ est une base de $F$ et donc $\underset\Kset\dim F = \card{V_1,\dots,V_p} = p$, alors $\rang(V_1,\dots,V_p) = p$.

    \textbf{Démonstration dans le sens direct : }

    Si $(V_1,\dots,V_p)$ est une famille génératrice de $F$, avec $\card{V_1,\dots,V_p} = p = \rang(V_1,\dots,V_p)$, alors $(V_1,\dots,V_p)$ est une base de $F$, donc forcément $(V_1,\dots,V_p)$ est une famille libre de $E$.

    Ainsi, $\rang(V_1,\dots,V_p) = p \iff (V_1,\dots,V_p) \text{ est une famille libre de } E$.
\end{proof}

\begin{property}
    Soient $V_1,\dots,V_p \in E$.

    \begin{enumerate}
        \item $\rang(0_E,V_1,\dots,V_p) = \rang(V_1,\dots,V_p)$
        \item $\forall \sigma \in S_p, \rang(V_{\sigma(1)},\dots,V_{\sigma(p)}) = \rang(V_1,\dots,V_p)$ où $S_p$ est l'ensemble des permutations de $\{1,\dots,p\}$.
        \item $\forall \lambda \in \Kset^*, \rang(\lambda V_1,V_2,\dots,V_p) = \rang(V_1,V_2,\dots,V_p)$
        \item $\forall \lambda_2,\dots,\lambda_p \in \Kset, \rang(V_1 + \sum_{k=2}^p \lambda_k V_k, V_2,\dots,V_p) = \rang(V_1,V_2,\dots,V_p)$
    \end{enumerate}
\end{property}

\begin{proof}
    \textbf{Démonstration de la quatrième propriété :}

    On pose $V_1' = V_1 + \sum_{k=2}^p \lambda_k V_k$, alors $\Vect(V_1',V_2,\dots,V_p) = \Vect(V_1,V_2,\dots,V_p)$, donc $\underset\Kset\dim \Vect(V_1',V_2,\dots,V_p) = \underset\Kset\dim \Vect(V_1,V_2,\dots,V_p)$, donc $\rang(V_1',V_2,\dots,V_p) = \rang(V_1,V_2,\dots,V_p)$.
\end{proof}

\begin{property}
    Toute famille échelonnée de $E$ est une famille libre de $E$.
\end{property}

\begin{proof}
    Soit $(v_1,\dots,v_p)$ une famille échelonnée de $E$.

    Supposons que $(v_1,\dots,v_p)$ n'est pas une famille libre de $E$.
    
    Alors il existe $\lambda_1,\dots,\lambda_p \in \Kset$, pas tous nuls, tels que $\sum_{i=1}^p \lambda_i v_i = 0_E$.

    Soit $i_0 = \min\{i \in [\![1,p]\!] \tq \lambda_i \ne 0\}$, alors $\alpha_{i_0} \ne 0$ et $\forall i < i_0, \lambda_i = 0$.

    On isole la somme à partir de $i_0$ : $\sum_{i=i_0}^p \lambda_i v_i = 0_E$.

    Donc $\alpha_{i_0} v_{i_0} + \sum_{i=i_0+1}^p \lambda_i v_i = 0_E$.

    Donc $\alpha_{i_0} \lambda_{i_0,\varphi(i_0)} e_{\varphi(i_0)} + \underbrace{\alpha_{i_0} \sum_{j=\varphi(i_0)+1}^n \lambda_{i_0,j} e_j + \sum_{i=i_0+1}^p \lambda_i v_i}_{\in \Vect(e_{\varphi(i_0)+1},\dots,e_n)} = 0_E$ (sachant que $\varphi(i_0) + 1 > \varphi(i_0)$.

    Or $(e_{\varphi(i_0)}, e_{\varphi(i_0)+1},\dots,e_n)$ est une famille libre de $E$, alors $\alpha_{i_0} \lambda_{i_0,\varphi(i_0)} = 0$, ce qui est absurde, car $\alpha_{i_0} \ne 0$ et $\lambda_{i_0,\varphi(i_0)} \ne 0$.

    D'où $(v_1,\dots,v_p)$ est une famille libre de $E$.
\end{proof}

\begin{property}
    \begin{outline}
    \1 $F+G$ est un sev. de $E$.
    \1 $F+G = \Vect(F \cup G)$.
    \end{outline}
\end{property}

\begin{proof}
    \textbf{Démonstration de la première propriété :}

    \begin{outline}
        \1 $F+G \subset E$, car $F \subset E$ et $G \subset E$.
        \1 $0_E = 0_E + 0_E \in F+G$, alors $F+G \ne \emptyset$.
        \1 Soient $u,v \in F+G$, alors $\exists x_1,x_2 \in F, y_1,y_2 \in G$ tels que $u = x_1 + y_1$ et $v = x_2 + y_2$, alors $u+v = (x_1+x_2) + (y_1+y_2) \in F+G$.
        \1 On a aussi $\forall \lambda \in \Kset, \forall u,v \in F+G, u + \lambda v = (x_1 + y_1) + \lambda (x_2 + y_2) = (x_1 + \lambda x_2) + (y_1 + \lambda y_2) \in F+G$, sachant que $F$ et $G$ sont des sev. de $E$.
        \1 Soient $\lambda \in \Kset$ et $u \in F+G$, alors $\exists x \in F, y \in G$ tels que $u = x+y$, alors $\lambda u = (\lambda x) + (\lambda y) \in F+G$.
        \1 Donc $F+G$ est un sev. de $E$.
    \end{outline}

    \textbf{Démonstration de la deuxième propriété :}

    \begin{outline}
        \1 $F+G \subset \Vect(F \cup G)$, car, par définition $\Vect(F \cup G)$ est le plus petit sous-espace vectoriel de $E$ contenant $F \cup G$.
        \1 On sait que $F \subset F+G$ et $G \subset F+G$, alors $F \cup G \subset F+G$, alors $\Vect(F \cup G) \subset F+G$.
        \1 \textit{Méthode alternative : On peut considérer les combinaisons linéaires de $F \cup G$ pour construire $\Vect(F \cup G)$, et on peut montrer que toutes ces combinaisons linéaires sont aussi des éléments de $F+G$, alors $\Vect(F \cup G) \subset F+G$.}
        \1 Donc $F+G = \Vect(F \cup G)$.
    \end{outline}
\end{proof}

\begin{property}
    Soient $E$ un $\Kset$-ev. et $F,G$ deux sous-ev. de $E$.

    Les propriétés suivantes sont équivalentes :

    \begin{outline}
        \1 $E = F \oplus G$.
        \1 $E = F + G$ et $F \cap G = \{0_E\}$.
        \1 $\forall x \in E, \exists ! (y,z) \in F\times G, x = y+z$.
    \end{outline}
\end{property}

\begin{proof}
    \textbf{Démonstration de la première équivalence :} [(1) $\iff$ (2)]

    $E = F \oplus G$ signifie que $E = F + G$ et que $F \cap G = \{0_E\}$, donc $E = F + G$ et $F \cap G = \{0_E\}$.

    Réciproquement, si $E = F + G$ et $F \cap G = \{0_E\}$, alors $E = F \oplus G$.

    \textit{Finalement, on a simplement utilisé la définition de la somme directe de $F$ et $G$ pour démontrer la première équivalence.}
    
    \textbf{Démonstration de la deuxième équivalence :} [(2) $\iff$ (3)]

    Supposons que $E = F + G$ et que $F \cap G = \{0_E\}$.

    \begin{outline}
    \1 \textbf{Démonstration de l'existence de la décomposition :} Soit $x \in E$, alors $\exists y \in F, z \in G$ tels que $x = y+z$, car $E = F + G$.

    \1 \textbf{Démonstration de l'unicité de la décomposition :}
    
    \2 \textbf{Dans le sens direct $\implies$ :} Supposons qu'il existe aussi $y' \in F, z' \in G$ tels que $x = y' + z'$, alors on a $(y-y') + (z-z') = 0_E$, alors $(y-y') = -(z-z')$, alors $(y-y') \in F\cap G$, alors $(y-y') = 0_E$, alors $y=y'$ et aussi $z=z'$, donc $\exists ! (y,z) \in F\times G, x=y+z$.

    \2 \textbf{Réciproquement}, supposons que $\forall x \in E, \exists ! (y,z) \in F\times G, x=y+z$, alors $\forall x \in E, \exists (y,z) \in F\times G, x=y+z$, alors $E=F+G$. De plus, supposons qu'il existe un élément non nul de $F\cap G$, alors il existe un élément non nul de $F$ qui est aussi dans $G$, ce qui contredit l'unicité de la décomposition de cet élément en somme d'un élément de $F$ et d'un élément de $G$. Donc $F\cap G=\{0_E\}$.
    \end{outline}

    Ainsi, les propriétés sont équivalentes.
\end{proof}

\begin{property}
    Soient $E$ un $\Kset$-ev. et $F_1,\dots,F_n$ des sev. de $E$.
    \begin{enumerate}
        \item $\sum_{k=1}^{n}F_k$ est un sev. de $E$.
        \item $\sum_{k=1}^n F_k = \bigoplus_{k=1}^n F_k \iff \forall x \in \sum_{k=1}^n F_k, \exists ! (x_1,\dots,x_n) \in \prod_{k=1}^n F_k \tq x = \sum_{k=1}^n x_k$.
    \end{enumerate}
\end{property}

\begin{proof}
    \begin{outline}[enumerate]
        \1 \textbf{Démonstration de la première propriété :} La première démonstration est facile. On utilise simplement la caractérisation d'un sous-espace vectoriel, sans oublier de mentionner que la somme est incluse dans $E$.
        \1 \textbf{Démonstration de la deuxième propriété :}
            \2 Dans le sens direct ($\implies$) :
                \3 Pour démontrer l'existence, on prend $x \in \sum_{k=1}^n F_k$, alors $\exists (x_1,\dots,x_n) \in \prod_{k=1}^n F_k$ tels que $x = \sum_{k=1}^n x_k$.
                
                \3 Pour démontrer l'unicité, supposons qu'il existe aussi $(y_1,\dots,y_n) \in \prod_{k=1}^n F_k$ tels que $x = \sum_{k=1}^n y_k$, alors $\sum_{k=1}^n (x_k - y_k) = 0_E$, alors $\forall k, x_k - y_k = 0_E$, alors $\forall k, x_k = y_k$ (parce que la somme est directe), donc $\exists ! (x_1,\dots,x_n) \in \prod_{k=1}^n F_k$ tels que $x = \sum_{k=1}^n x_k$.

            \2 Dans le sens indirect ($\impliedby$) :
                \3 Soient $(x_1,\dots,x_n) \in \prod_{k=1}^n F_k$ tels que $\sum_{k=1}^n x_k = 0_E$.
                \3 Alors, par unicité de la décomposition, on a $\forall k, x_k = 0_E$, donc $\sum_{k=1}^n F_k$ est directe.
    \end{outline}
\end{proof}

\begin{property}
    Soient $E$ un $\Kset$-ev. et $F_1,\dots,F_n$ des sev. de $E$.

    Les propriétés suivantes sont équivalentes :

    \begin{enumerate}
        \item $\sum_{k=1}^n F_k = \bigoplus_{k=1}^n F_k$.
        \item $\forall i \in [\![1,n]\!], F_i \cap \sum_{\substack{k=1 \\ k \ne i}}^n F_k = \{0_E\}$.
        \item $\forall i \in [\![2,n]\!], F_i \cap \sum_{k=1}^{i-1} F_k = \{0_E\}$.
    \end{enumerate}
\end{property}

\begin{proof}
    La démonstration se fait par raisonnement circulaire : (1) $\implies$ (2) $\implies$ (3) $\implies$ (2).

    \begin{outline}

        \1 \textbf{Démonstration de la troisième implication :} [(3) $\implies$ (2)]

        On va raisonner par absurde.

        Supposons que $\exists i \in [\![1,n]\!]$ tel que $x_i \ne 0$.


        On pose $S = \{ k \in [\![1,n]\!] \tq x_k \ne 0\}$.


        Alors $\begin{cases}
            S \text{ est majorée.}\\
            S \text{ est non vide.}\\
            \implies S \text{ admet une valeur maximale.}
        \end{cases}$

        Posons $s = \max S$, alors $s \in S$ et $x_s \ne 0$.
        
        On a $\sum_{k=1}^n x_k = 0_E$.

        On décompose la somme $\sum_{k=1}^n x_k$ en trois parties :
        
        $$\sum_{k=1}^n x_k = \underbrace{\sum_{k=1}^{s-1} x_k}_{\in \sum_{k=1}^{s-1} F_k} + \underbrace{x_s}_{\in F_s} + \underbrace{\sum_{k=s+1}^n x_k}_{= 0_E}$$

        Alors : $\sum_{k=1}^{s-1} x_k = -x_s$.

        Donc $x_s \in F_s \cap \sum_{k=1}^{s-1} F_k$, et comme l'intersection est réduite au vecteur nul $\{0_E\}$, alors $x_s = 0_E$, ce qui est absurde.
    \end{outline}
\end{proof}

\begin{theorem}
    Soient $E$ un $\Kset$-ev. et $F_1,\dots,F_n$ des sev. de $E$ de dimension finie.

    On pose $\forall 1 \le i \le n, B_i$ est une base de $F_i$.

    Alors $B = \bigcup_{i=1}^n B_i$ est une base de $\bigoplus_{k=1}^n F_k$, et $\underset\Kset\dim(\bigoplus_{k=1}^n F_k) = \sum_{i=1}^n \underset\Kset\dim(F_i)$.
\end{theorem}

\begin{proof}
    Pour tout $1 \le i \le n$, on pose $B_i = (e_{i,j})_{1 \le j \le n_i}$, avec $n_i = \card B_i = \underset\Kset\dim(F_i)$.

    Soit $x \in \bigoplus_{k=1}^n F_k$, alors $\exists ! (x_1,\dots,x_n) \in \prod_{k=1}^n F_k$ tels que $x = \sum_{k=1}^n x_k$.

    Et on a $\forall 1 \le i \le n, x_i \in F_i = \Vect(B_i)$, alors $\exists ! (\lambda_{i,j})_{1 \le j \le n_i} \in \Kset^{n_i}$ tels que $x_i = \sum_{j=1}^{n_i} \lambda_{i,j} e_{i,j}$.

    Donc $x = \sum_{i=1}^n\sum_{j=1}^{n_i} \lambda_{i,j} e_{i,j}$.

    Et donc $B = (e_i,j)_{\substack{1 \le i \le n \\ 1 \le j \le n_i}}$ est une famille génératrice de $\bigoplus_{k=1}^n F_k$.
\end{proof}

\begin{theorem}[Lemme]
    Tout sous-espace vectoriel $F$ d'un $\Kset$-ev. de dimension finie $E$ admet un supplémentaire dans $E$, c'est à dire que :

    $$\forall F \text{ sev. de } E, \exists G \text{ sev. de } E, E = F \oplus G$$
    \label{lem:sous_espace_vectoriel_admet_suppl}
\end{theorem}

\begin{proof}
    Soit $F$ un sous-espace vectoriel de $E$. On admettra que celui-ci est différent de $\{0_E\}$, sinon $E = \{0_E\} \oplus E$ et la propriété est vérifiée.

    \textbf{Montrons que $\exists H$ sev. de $E$ tel que $F \oplus H = E$.}

    \begin{outline}
        \1 On sait que $F$ est de dimension finie (parce que $E$ est de dimension finie). On pose $p = \underset\Kset\dim(F)$ et on considère une base $B_1 = (e_1,\dots,e_p)$ de $F$.
        \1 $B_1$ est une famille libre de $E$, alors on peut compléter $B_1$ en une base $B = (e_1,\dots,e_p,e_{p+1},\dots,e_n)$ de $E$, en posant $n = \underset\Kset\dim(E) \ge p$, en utilisant le théorème de la base incomplète (voir théorème \ref{thm:base_incomplete}).
        \1 On pose $H = \Vect(e_{p+1},\dots,e_n)$, alors $H$ est un sev. de $E$.
        \1 On a donc $\forall x \in E, \exists ! (\lambda_1,\dots,\lambda_n) \in \Kset^n$ tels que $x = \sum_{i=1}^n \lambda_i e_i$.
        \1 Alors $x = \underbrace{\sum_{i=1}^p \lambda_i e_i}_{\in F} + \underbrace{\sum_{i=p+1}^n \lambda_i e_i}_{\in H}$.
        \1 Donc $\forall x \in E, \exists ! (y,z) \in F\times H, x = y+z$, alors $E = F \oplus H$.
    \end{outline}
\end{proof}

\begin{theorem}[Formule de Grassmann]
    Soient $E$ un $\Kset$-ev. et $F,G$ des sev. de $E$ de dimension finie.

    Alors $\underset\Kset\dim(F+G) = \underset\Kset\dim(F) + \underset\Kset\dim(G) - \underset\Kset\dim(F \cap G)$.
\end{theorem}

\begin{proof}
    \begin{outline}
        \1 On a $F= \Vect(B_1)$ et $G = \Vect(B_2)$, avec $B_1$ et $B_2$ des bases de $F$ et $G$ respectivement.
        \1 Donc $F+G = \Vect(B_1) + \Vect(B_2) = \Vect(B_1 \cup B_2)$.
        \1 Donc $F+G$ est de dimension finie.
        \1 On a $F \cap G$ est un sev. de $F$, avec $F$ de dimension finie.
        \1 Alors, d'après le lemme \ref{lem:sous_espace_vectoriel_admet_suppl}, $F \cap G$ admet un supplémentaire dans $F$, c'est à dire que $\exists H$ sev. de $F$ tel que $F = (F \cap G) \oplus H$.
        \1 Montrons alors que $F+G = H \oplus G$.
        \2 On a $G \cap H = G \cap (F \cap H) = (F \cap G) \cap H = \{0_E\}$, sachant que $H \subset F$. (1)
        \2 Soit $x \in F+G$, alors $\exists y \in F, z \in G$ tels que $x = y+z$.
        \2 Or $y \in F \implies \exists ! (z_1,h) \in (F\cap G)\times H$ tels que $y = z_1 + h$, alors $x = h + (z_1 + z)$, avec $h \in H$ et $z_1 + z \in G$, alors $x \in H + G$. (2)
        \2 De (1) et (2), on a $F+G = H \oplus G$.
        \1 Donc $\underset\Kset\dim(F+G) = \underset\Kset\dim(H) + \underset\Kset\dim(G)$, et $\underset\Kset\dim(F) = \underset\Kset\dim(F \cap G) + \underset\Kset\dim(H)$, alors $\underset\Kset\dim(F+G) = \underset\Kset\dim(F) + \underset\Kset\dim(G) - \underset\Kset\dim(F \cap G)$.
    \end{outline}
\end{proof}

\begin{proof}
    La démonstration se fait par raisonnement circulaire : (1) $\implies$ (2) $\implies$ (3) $\implies$ (1).

    \begin{outline}
        \1 La première implication [(1) $\implies$ (2)] est facile à démontrer. On utilise simplement la définition de la somme directe et la formule de Grassmann pour montrer que $F\cap G = \{0_E\}$ et que $\underset\Kset\dim(E) = \underset\Kset\dim(F) + \underset\Kset\dim(G)$.
        \1 Démonstration de la deuxième implication [(2) $\implies$ (3)] :
        
            \2 On sait que la dimension de $F+G$ est égale à $\underset\Kset\dim(F) + \underset\Kset\dim(G) - \underset\Kset\dim(F \cap G)$, or $F \cap G$ est réduit au singleton vecteur nul $\{0_E\}$.
            \2 Alors $\underset\Kset\dim(F \cap G) = 0$, alors $\underset\Kset\dim(F+G) = \underset\Kset\dim(F) + \underset\Kset\dim(G) = \underset\Kset\dim(E)$.
            \2 Et on sait que $F + G \subset E$, alors $E = F + G$.

        \1 Démonstration de la troisième implication [(3) $\implies$ (1)] :
            \2 On sait que $E = F+G$, donc $\underset\Kset\dim(E) = \underset\Kset\dim(F+G)$.
            \2 C'est-à-dire que $\underset\Kset\dim(E) = \underset\Kset\dim(F) + \underset\Kset\dim(G) - \underset\Kset\dim(F \cap G)$, alors $\underset\Kset\dim(F \cap G) = 0$, alors $F \cap G = \{0_E\}$ (c'est le seul sous-espace vectoriel de dimension nulle), alors $E = F \oplus G$.
    \end{outline}
\end{proof}

\begin{property}
    Soient $E$ un $\Kset$-ev. et $H$ un sev. de $E$ tel que $H \ne E$. Alors :

    $H$ est un hyperplan de $E$ $\iff$ $\forall b \in E\setminus H, E = H \oplus \Kset \cdot b$.
\end{property}

\begin{proof}
    \begin{outline}
        \1 Démonstration dans le sens indirect $(\impliedby)$ : En utilisant la définition.
        \1 Démonstration de le sens direct $(\implies)$ :
        \2 $H$ est un hyperplan de $E$, alors $\exists a \in E\setminus\{0_E\}$ tel que $E = H \oplus \Kset \cdot a$.
        \2 Prenons $b \in E\setminus H$.
        \2 Comme $b \in E = H \oplus \Kset \cdot a$.
        \2 Alors $\exists ! (h_1,\lambda_1) \in H\times\Kset$ tels que : 
        $$b = h_1 + \lambda_1 a$$
        \2 Si $\lambda_1 = 0$, alors $b = h_1 \in H$, ce qui est absurde, alors $\lambda_1 \ne 0$.
        \2 Et on a, $\forall x \in E, \exists ! (h_2,\lambda_2) \in H\times\Kset$ tels que :
        \begin{align*}
            x &= h_2 + \lambda_2 a \\
            &= h_2 + \lambda_2 \cdot \lambda_1^{-1} (b - h_1) \\
            &= (h_2 - \lambda_2 \cdot \lambda_1^{-1} h_1) + \lambda_2 \cdot \lambda_1^{-1} b
        \end{align*}
        \2 Donc $E = H \oplus \Kset \cdot b$.
        \2 Si $x \in H_1 \cap \Kset \cdot b$, alors $\exists \lambda \in \Kset$ tel que $x = \lambda b$ et $x\in H$.
        \2 Mais comme $b \notin H$ alors $\lambda = 0$ alors $x = 0_E$.
        \2 Alors $H \cap \Kset \cdot b = \{0_E\}$, alors $E = H \oplus \Kset \cdot b$. CQFD.
    \end{outline}
\end{proof}

\begin{property}
    Soient $E$ un $\Kset$-ev. et $\underset\Kset\dim(E) = n \in \N^*$, et $H$ un sous-espace vectoriel de $E$.

    $H$ est un hyperplan de $E$ $\iff$ $\underset\Kset\dim(H) = n-1$.
\end{property}

\begin{proof}
    \begin{outline}
        \1 Dans le sens direct $\implies$ la démonstration est facile et se fait en utilisant la définition de l'hyperplan et la dimension d'une somme directe.
        \1 Dans le sens indirect $\impliedby$ :
        \2 On pose $H$ un sous-espace vectoriel de $E$ de dimension $n-1$.
        \2 On sait que $E = H \oplus F$, sachant que tout sous-espace vectoriel de $E$ admet un supplémentaire dans $E$ (voir le lemme \ref{lem:sous_espace_vectoriel_admet_suppl}).
        \2 Donc $\underbrace{\underset\Kset\dim(E)}_{=n} = \underbrace{\underset\Kset\dim(H)}_{=n-1} + \underset\Kset\dim(F)$.
        \2 Alors $\underset\Kset\dim(F) = 1$.
        \2 Donc $\exists a \in E\setminus\{0_E\}$ tel que $F = \Kset \cdot a$.
        \2 Donc $E = H \oplus \Kset \cdot a$, alors $H$ est un hyperplan de $E$.
    \end{outline}
\end{proof}


\vspace{1em}


\section*{Extraits : 260410 CH14-B}


\begin{property}
    Soient $E,F$ deux $\Kset$-ev. et $f \in \mathcal L(E,F)$.
    \begin{enumerate}
        \item $f(0_E) = 0_F$.
        \item $\forall x_1,\dots,x_n\in E, \forall \lambda_1,\dots,\lambda_n\in\Kset$, $f(\lambda_1 x_1 + \dots + \lambda_n x_n) = \lambda_1 f(x_1) + \dots + \lambda_n f(x_n)$.
    \end{enumerate}
\end{property}

\begin{proof}
    \begin{outline}[enumerate]
        \1 $f(0_E) = f(0_E + 0_E) = f(0_E) + f(0_E)$, donc $f(0_E) = 0_F$.
        \1 Par récurrence sur $n$. Le cas $n=1$ est trivial. Supposons que la propriété soit vérifiée pour $n$, montrons-la pour $n+1$. Soient $x_1,\dots,x_{n+1}\in E$ et $\lambda_1,\dots,\lambda_{n+1}\in\Kset$, alors
        \begin{align*}
            f(\lambda_1 x_1 + \dots + \lambda_{n+1} x_{n+1}) &= f((\lambda_1 x_1 + \dots + \lambda_n x_n) + \lambda_{n+1} x_{n+1})\\
            &= f(\lambda_1 x_1 + \dots + \lambda_n x_n) + f(\lambda_{n+1} x_{n+1})\\
            &= (\lambda_1 f(x_1) + \dots + \lambda_n f(x_n)) + \lambda_{n+1} f(x_{n+1})\\
            &= \lambda_1 f(x_1) + \dots + \lambda_{n+1} f(x_{n+1}) \text{ par H.R.}
        \end{align*}
    \end{outline}
\end{proof}

\begin{property}
    Soient $E,F$ deux $\Kset$-ev.
    $$\parens{\mathcal L(E,F), +, \cdot} \text{ est un } \Kset\text{-ev.}$$
\end{property}

\begin{proof}
    \begin{outline}
        \1 On a $\mathcal L(E,F) \subset F^E$.
        \1 L'application nulle (notée $0_{F^E}$) est bien linéaire. Donc $0_F \in \mathcal L(E,F)$.
        \1 Soient $f,g \in \mathcal L(E,F)$ et $\alpha \in \Kset$. Montrons que $f + \alpha g \in \mathcal L(E,F)$. Soient $x,y \in E$ et $\lambda \in \Kset$, alors :
        \begin{align*}
            (f + \alpha g)(x + \lambda y) &= f(x + \lambda y) + \alpha g(x + \lambda y)\\
            &= (f(x) + \lambda f(y)) + \alpha(g(x) + \lambda g(y))\\
            &= (f(x) + \alpha g(x)) + (\lambda f(y) + \alpha \lambda g(y))\\
            &= (f + \alpha g)(x) + \lambda (f + \alpha g)(y)
        \end{align*}
        \1 Ainsi, $f + \alpha g \in \mathcal L(E,F)$.
        \1 Par conséquent, $\mathcal L(E,F)$ est stable par addition et multiplication par un scalaire, et contient l'élément neutre pour l'addition. Donc $\parens{\mathcal L(E,F), +, \cdot}$ est un sous-espace vectoriel de $F^E$. Or $F^E$ est un $\Kset$-ev. (chapitre précédent), donc $\mathcal L(E,F)$ est un $\Kset$-ev.
    \end{outline}
\end{proof}

\begin{property}
    Soient $E,F,G$ trois $\Kset$-ev. et $f \in \mathcal L(E,F)$, $g \in \mathcal L(F,G)$, alors $g \circ f \in \mathcal L(E,G)$.
\end{property}

\begin{proof}
    Soient $x,y \in E$ et $\lambda \in \Kset$, alors :
    \begin{align*}
        (g \circ f)(x + \lambda y) &= g(f(x + \lambda y))\\
        &= g(f(x) + \lambda f(y))\\
        &= g(f(x)) + \lambda g(f(y))\\
        &= (g \circ f)(x) + \lambda (g \circ f)(y)
    \end{align*}
    Ainsi, $g \circ f$ est une application linéaire de $E$ vers $G$, c'est-à-dire $g \circ f \in \mathcal L(E,G)$.
\end{proof}

\begin{property}
    $$(\mathcal L(E), +, \circ) \text{ est un anneau non-commutatif.}$$
\end{property}

\begin{proof}
    On doit montrer que :
    \begin{outline}
        \1 La loi $\circ$ est une loi de composition interne sur $\mathcal L(E)$.
        \1 $(\mathcal L(E), +)$ est un groupe abélien.
        \1 Il existe un élément neutre pour $\circ$ dans $\mathcal L(E)$.
        \1 $\circ$ est associative.
        \1 $\circ$ est distributive à gauche et à droite par rapport à $+$.
    \end{outline}
    On procède alors à la vérification de chacune de ces propriétés :
    \begin{outline}
        \1 On a $(\mathcal L(E), +, \cdot)$ est un $\Kset$-ev. Donc $(\mathcal L(E), +)$ est un groupe abélien.
        \1 D'après la propriété précédente, on a $\circ$ est une loi de composition interne dans $\mathcal L(E)$.
        \1 On a $\Id_E \in \mathcal L(E)$ et pour tout $f \in \mathcal L(E)$, on a $\Id_E \circ f = f \circ \Id_E = f$. Donc $\Id_E$ est l'élément neutre pour la loi $\circ$.
        \1 On a $\circ$ est associative, c'est-à-dire que pour tous $f,g,h \in \mathcal L(E)$, on a $(f \circ g) \circ h = f \circ (g \circ h)$.
        \1 Aussi, $\circ$ est distributive par rapport à $+$, c'est-à-dire que pour tous $f,g,h \in \mathcal L(E)$, on a $f \circ (g + h) = f \circ g + f \circ h$ et $(f + g) \circ h = f \circ h + g \circ h$.
        \1 Cependant, $\circ$ n'est pas commutative en général, c'est-à-dire qu'il existe des $f,g \in \mathcal L(E)$ tels que $f \circ g \neq g \circ f$ (on démontre ça par contre-exemple).
            \2 On prend $E = \R^n$ avec $n \ge 2$
            \2 On pose $f : \begin{matrix}
            R^n &\to &R^n\\
            (x_1,\dots,x_n) &\mapsto &(x_1,0,\dots,0)
            \end{matrix}$ et $g : \begin{matrix}
            R^n &\to &R^n\\
            (x_1,\dots,x_n) &\mapsto &(x_2,0,\dots,0)
            \end{matrix}$
            \2 Alors $\forall (x_1,\dots,x_n) \in \R^n$, on a $f(x) = x_1\cdot e_1$ et $g(x) = x_2\cdot e_1$.
            \2 $(f \circ g)(e_2) = f(g(e_2)) = f(e_1) = e_1$ alors que $(g \circ f)(e_2) = g(f(e_2)) = g(0) = 0_{\R^n}$. Donc $f \circ g \neq g \circ f$.
            \2 Ainsi, $\circ$ n'est pas commutative en général.
            \3 On pourra aussi noter que cet anneau n'est ni intègre (car il n'est pas commutatif et un anneau intègre est défini comme étant un anneau commutatif sans diviseur de zéro) ni semi-intègre\index{Semi-intègre} (car, dans ce dernier contre-exemple $f\circ g(x) = 0_{R^n}$ mais $f \neq 0$ et $g \neq 0$).
    \end{outline}
\end{proof}

\begin{property}
    Soient $E,F$ deux $\Kset$-ev. et $f \in \mathcal L(E,F)$.
    \begin{enumerate}
        \item Si $H$ est un sev de $E$ alors $f(H)$ est un sev de $F$.
        \item Si $H$ est un sev de $F$ alors $f^{-1}(H)$ est un sev de $E$.
    \end{enumerate}
\end{property}

\begin{proof}
    recop <З
\end{proof}

\begin{property}
    Soit $f \in \mathcal L(E,F)$.
    \begin{enumerate}
        \item $f$ est injective $\iff$ $\Ker(f) = \{0_E\}$.
        \item $f$ est surjective $\iff$ $f(E) = F$.
    \end{enumerate}
\end{property}

\begin{proof}
    La preuve est très similaire à celle faite dans le chapitre 10, propriété 17, pages 29-30. On utilise juste les lois $+$ qu'on vient de définir maintenant au lieu des lois $*$ et $T$ définies dans le chapitre 10.
\end{proof}

\begin{property}
    Soient $f \in \mathcal L(E,F)$ et $B$ une base de $E$ (pas forcément finie).

    \begin{enumerate}
        \item $f$ est injective $\iff$ $f(B)$ est libre.
        \item $f$ est surjective $\iff$ $f(B)$ est génératrice de $F$.
        \item $f \in \Isom(E,F) \iff f(B)$ est une base de $F$.
    \end{enumerate}
\end{property}

\begin{proof}
    \begin{outline}
        \1 Démonstration de la première équivalence :
        \2 ($\implies$) Dans le sens direct :
        \3 On a $f(B) = \{f(e) \tq e \in B\}$.
        \3 Soient $u_1,\dots,u_n \in f(B)$ avec $n \in \N^*$. Alors $\exists e_1,\dots,e_n \in B$ tels que $u_i = f(e_i)$ pour $i=1,\dots,n$.
        \3 Soient $\lambda_1,\dots,\lambda_n \in \Kset$ tels que $\sum_{k=1}^{n}\lambda_k u_k = 0_F$.
        \begin{align*}
            \implies \sum_{k=1}^{n}\lambda_k u_k = 0_F &\implies \sum_{k=1}^{n}\lambda_k f(e_k) = 0_F\\
            &\implies f\left(\sum_{k=1}^{n}\lambda_k e_k\right) = 0_F\\
            &\implies \sum_{k=1}^{n}\lambda_k e_k \in \Ker(f)\\
            &\implies \sum_{k=1}^{n}\lambda_k e_k = 0_E \text{ (car $f$ est injective)}\\
            &\implies \lambda_1 = \dots = \lambda_n = 0
        \end{align*}
        \3 Donc $(u_1,\dots,u_n)$ est libre, et par conséquent, $f(B)$ est libre.
        \2 ($\impliedby$) Dans le sens inverse :
        \3 Soit $x \in \Ker(f)$, alors $x\in E$ et $f(x) = 0_F$.
        \3 Comme $x \in E$ et $B$ est une base de $E$, alors $\exists e_1,\dots,e_n \in B$ et $\lambda_1,\dots,\lambda_n \in \Kset$ tels que $x = \sum_{k=1}^{n}\lambda_k e_k$.
    \end{outline}
\end{proof}

\begin{proof}[2e remarque précédente]
    \begin{outline}
        \1 On pose $B = (e_1,\dots,e_n)$ une base de $E$ avec $n = \dim(E)$.
        \1 Donc $\Imf(f) = f(E) = f(\Vect(e_1,\dots,e_n))$ et $\Imf(f) = \Vect(f(e_1),\dots,f(e_n))$.
        \1 Donc $\Imf(f)$ est de dimension finie ($< \infty$), et par conséquent, le rang de $f$ existe.
    \end{outline}
\end{proof}

\begin{property}
    Soient $E,F$ deux $\Kset$-ev. et $f \in \mathcal L(E,F)$.

    Si $H$ est un supplémentaire de $\Ker(f)$ dans $E$, alors $H$ est isomorphe à $\Imf(f)$, donc :
    $$H \sim \Imf(f)$$
\end{property}

\begin{proof}
    Soit $H$ un supplémentaire de $\Ker(f)$ dans $E$. Alors $E = \Ker(f) \oplus H$.

    On cherche à déterminer une application linéaire bijective de $H$ vers $\Imf(f)$.

    On pose $g : \begin{matrix}
    H &\to &\Imf(f)\\
    x &\mapsto &f(x)
    \end{matrix}$

    \textit{Au final, cela revient à considérer la restriction de $f$ à $H$.}

    \begin{outline}
        \1 $\forall x \in H$, $g(x) = f(x) \in \Imf(f)$, donc $g$ est une application de $H$ vers $\Imf(f)$.
        \1 Comme $f$ est linéaire, alors $g$ est linéaire $\in \mathcal L(H, \Imf(f))$.
        \1 Soit $x \in \Ker(g)$. Cela équivaut à dire que $x\in H et g(x) = f(x) = 0_F$, c'est-à-dire que $x \in \Ker(f)$, ce qui équivaut à dire que $H \cap \Ker(f) = \{0_E\}$, alors $x = 0_E$. Donc $g$ est injective.
        \1 On a $E = H \oplus \Ker(f)$, donc $f(E) = f(H) + f(\Ker(f)) = f(H) + \{0_F\} = f(H)$, donc $g$ est surjective, car $g(H) = f(H) = f(E) = \Imf(f)$.
        \1 Donc $g$ est bijective, et par conséquent, $H$ est isomorphe à $\Imf(f)$.
    \end{outline}
\end{proof}

\begin{theorem}[Théorème du rang]
    Soient $E,F$ deux $\Kset$-ev. et $f \in \mathcal L(E,F)$, et $\underset\Kset\dim(E) = n \in \N^* < \infty$.
    
    Alors $\underset\Kset\dim(\Ker(f)) < \infty$, et $\rang(f) = n - \underset\Kset\dim(\Ker(f))$.
\end{theorem}

\begin{proof}
    \begin{outline}
        \1 $\Ker(f)$ est un sev de $E$, donc $\Ker(f)$ est de dimension finie, et $H \subset E$ est un supplémentaire de $\Ker(f)$ dans $E$, alors $E = \Ker(f) \oplus H$.
        \1 D'après la propriété précédente, $H$ est isomorphe à $\Imf(f)$, donc $\underset\Kset\dim(H) = \underset\Kset\dim(\Imf(f)) = \rang(f)$.
        \1 Or on sait que $E = \Ker(f) \oplus H$, donc $\underset\Kset\dim(E) = \underset\Kset\dim(\Ker(f)) + \underset\Kset\dim(H)$, et par conséquent, $\rang(f) = \underset\Kset\dim(H) = n - \underset\Kset\dim(\Ker(f))$.
    \end{outline}
\end{proof}

\begin{theorem}
    Si $f \in \mathcal L(E,F)$ avec $\underset\Kset\dim(E) = \underset\Kset\dim(F) = n < \infty$, alors :
    $$f \text{ est injective} \iff f \textrm{ est surjective} \iff f \in \Isom(E,F)$$
\end{theorem}

\begin{proof}
    \begin{align*}
    f \text{ est injective} &\iff \Ker(f) = \{0_E\}\\
    &\iff \rang(f) = n - \underset\Kset\dim(\Ker(f)) = n - 0 = n\\
    &\iff \dim(\Imf(f)) = n = \dim(F)\\
    &\iff \Imf(f) = F\\
    &\iff f \text{ est surjective}
    \end{align*}
\end{proof}

\begin{property}[Rang d'une composée d'applications linéaires]
    Soient $E, F, G$ trois $\Kset$-ev. et $f \in \mathcal L(E,F)$, $g \in \mathcal L(F,G)$, tels que $\underset\Kset\dim(E) < \infty$ et $\underset\Kset\dim(F) < \infty$.

    \begin{enumerate}
        \item $\rang(g\circ f) = \rang(f) - \underset\Kset\dim(\underbrace{\Ker(g) \cap \Imf(f)}_{=\Ker(g\circ f)})$.
        \item Si $g$ est injective, alors $\rang(g\circ f) = \rang(f)$.
        \item Si $f$ est surjective, alors $\rang(g\circ f) = \rang(g)$.
        \item \textbf{Les applications linéaires bijectives ne changent pas le rang :} Si $f \in \Isom(E,F)$ et $g \in \Isom(F,G)$, alors $\rang(g\circ f) = \rang(f) = \rang(g)$.
    \end{enumerate}
\end{property}

\begin{proof}
    \begin{outline}[enumerate]
    \1 On peut considérer $g\circ f$ comme l'application linéaire partant de $\Imf(f)$ vers $G$, et $\Imf(f)$ est de dimension finie, donc le rang de $g\circ f$ existe, et on applique donc le théorème du rang à $g\circ f$ pour retrouver le résultat.
    \1 Si $g$ est injective, alors $\Ker(g) = \{0_F\}$, donc $\Ker(g) \cap \Imf(f) = \{0_F\}$, et par conséquent, $\rang(g\circ f) = \rang(f)$.
    \1 Si $f$ est surjective, alors $\Imf(f) = F$, donc $\Ker(g) \cap \Imf(f) = \Ker(g)$, et par conséquent, $\rang(g\circ f) = \rang(f) - \underset\Kset\dim(\Ker(g)) = \underset\Kset\dim(F) - \underset\Kset\dim(\Ker(g)) = \rang(g)$.
    \end{outline}
\end{proof}

\begin{property}
    Soient $E$ un $\Kset$-ev. et $(E_i)_{1\le i \le n}$ une famille de sev de $E$ tels que $E = \bigoplus_{i=1}^n E_i$.

    Soient $u_i \in \mathcal L(E_i,F)$ pour $i=1,\dots,n$.

    Alors $\exists! u \in \mathcal L(E,F)$ tel que $U_{|E_i} = u_i$ pour $i=1,\dots,n$ (on note $u_{|E_i}$ la restriction de $u$ à $E_i$).

    $u$ est donnée par $u : \begin{matrix}
    E &\to &F\\
    \sum_{i=1}^n x_i &\mapsto &\sum_{i=1}^n u_i(x_i)
    \end{matrix}$ avec $x_i \in E_i$ pour $i=1,\dots,n$.
\end{property}

\begin{proof}
    À faire en tant qu'exercice

    (il faut montrer que $u$ est bien définie, linéaire, que sa restriction à $E_i$ est bien égale à $u_i$, et que $u$ est la seule application linéaire vérifiant ces propriétés).
\end{proof}

\begin{proof}
    Comme on sait que $p(x) = x_1$ pour $x = x_1 + x_2$ avec $x_1 \in F$ et $x_2 \in G$, alors $p^2(x) = p(p(x)) = p(x_1) = x_1 = p(x)$, et par conséquent, $p^2 = p$, donc $p$ est un projecteur.
\end{proof}

\begin{property}
    Soit $f \in \mathcal L(E)$.
    
    Alors $f$ est un projecteur $\implies$ $E = \Imf(f) \oplus \Ker(f)$.
\end{property}

\begin{proof}
    Réalisée en tant qu'exercice.

    \begin{outline}
        \1 On montre que l'intersection de $\Imf(f)$ et $\Ker(f)$ est réduite au vecteur nul.
        \2 On prend $x \in \Imf(f) \cap \Ker(f)$, alors $\exists y \in E$ tel que $x = f(y)$, et $f(x) = 0_E$, donc $f(f(y)) = 0_E$, donc $f(y) = 0_E$, donc $x = f(y) = 0_E$.
        \1 On montre que $E = \Imf(f) + \Ker(f)$.
        \2 On remarque directement que $\forall x \in E, f^2(x) = f(x) \implies f(f(x)) = f(x) \implies f(f(x)-x)=0_E \implies f(x)-x \in \Ker(f)$. Donc $x = f(x) + (x - f(x))$ avec $f(x) \in \Imf(f)$ et $x - f(x) \in \Ker(f)$, et par conséquent, $E = \Imf(f) + \Ker(f)$.
        \1 Donc $E = \Imf(f) \oplus \Ker(f)$.
    \end{outline}
\end{proof}

\begin{proof}
    On a $\forall x \in \Ker(f), f(x) = 0_E$ (par définition) et $\forall x \in \Imf(f), \exists y \in E$ tel que $x = f(y)$, donc $f(x) = f(f(y)) = f(y) = x$ (car $f$ est un projecteur), et par conséquent, sachant que $E = \Imf(f) \oplus \Ker(f)$, $f$ est une projection sur $\Imf(f)$ parallèlement à $\Ker(f)$.
\end{proof}

\begin{proof}
    La démonstration se fait en utilisant le lemme des noyaux, qui sera vu prochainement. On peut passer de $f^2 = f$ à un polynôme annulateur de $f$ qui est $X^2 - X$, et alors on peut factoriser ce polynôme annulateur en $(X)(X-1)$, et alors on peut appliquer le lemme des noyaux pour trouver que $E = \Ker(f) \oplus \Ker(f - \Id_E)$, et par conséquent, $\Imf(f) = \Ker(f - \Id_E)$.
\end{proof}

\begin{property}
    Soit $f \in \mathcal L(E)$. Alors $f$ est une involution $\implies$ $E = \Ker(f - \Id_E) \oplus \Ker(f + \Id_E)$.
\end{property}

\begin{proof}
    \begin{outline}
        \1 Montrons que $\Ker(f - \Id_E) \cap \Ker(f + \Id_E) = \{0_E\}$.
        \2 Soit $x \in \Ker(f - \Id_E) \cap \Ker(f + \Id_E)$, alors $f(x) = x$ et $f(x) = -x$, donc $x = 0_E$, et par conséquent, $\Ker(f - \Id_E) \cap \Ker(f + \Id_E) = \{0_E\}$.
        \1 Montrons que $E = \Ker(f - \Id_E) + \Ker(f + \Id_E)$.
        \2 Soit $x \in E$. On va procéder par analyse-synthèse.
        \2 \textit{Analyse :}
        \3 On suppose que $x = x_1 + x_2$ avec $x_1 \in \Ker(f - \Id_E)$ et $x_2 \in \Ker(f + \Id_E)$.
        \3 Alors $f(x) = f(x_1) + f(x_2) = x_1 - x_2$.
        \3 Donc $\begin{cases}x_1 + x_2 = x\\ x_1 - x_2 = f(x)\end{cases}$, et par conséquent, $\begin{cases}x_1 = \frac{x + f(x)}{2}\\ x_2 = \frac{x - f(x)}{2}\end{cases}$.
        \2 \textit{Synthèse :}
        \3 On pose $x_1 = \frac{x + f(x)}{2}$ et $x_2 = \frac{x - f(x)}{2}$.
        \3 Alors $f(x_1) = \frac{f(x) + f^2(x)}{2} = \frac{f(x) + x}{2} = x_1$.
        \3 Et $f(x_2) = \frac{f(x) - f^2(x)}{2} = \frac{f(x) - x}{2} = -x_2$.
        \3 Donc $x_1 \in \Ker(f - \Id_E)$ et $x_2 \in \Ker(f + \Id_E)$.
        \3 Alors $x = x_1 + x_2$.
        \1 Donc $E = \Ker(f - \Id_E) \oplus \Ker(f + \Id_E)$.
    \end{outline}
\end{proof}

\begin{theorem}[Lemme]
    Soient $E$ un $\Kset$-ev. et $f,g \in \mathcal L(E)$.
    $$f \circ g = 0 \iff \Imf(g) \subset \Ker(f)$$
\end{theorem}

\begin{proof}
    \begin{outline}
        \1 ($\implies$) Soit $y \in \Imf(g)$, alors $\exists x \in E$ tel que $y = g(x)$, et alors $f(y) = f(g(x)) = 0_E$, et par conséquent, $y \in \Ker(f)$, et par conséquent, $\Imf(g) \subset \Ker(f)$.
        \1 ($\impliedby$) Soit $x \in E$, alors $g(x) \in \Imf(g)$, et par conséquent, $g(x) \in \Ker(f)$, donc $f(g(x)) = 0_E$, et par conséquent, $f\circ g = 0$.
    \end{outline}
\end{proof}

\begin{theorem}[Lemme de la décomposition des noyaux]
    Soient $E$ un $\Kset$-ev. et $f \in \mathcal L(E)$ et $P \in \Kset[X]$ tel que $P = P_1\cdot P_2$ avec $P$ un polynôme annulateur de $f$ et $\pgcd(P_1,P_2) = 1$.
    $$E = \Ker(P_1(f)) \oplus \Ker(P_2(f))$$
\end{theorem}

\begin{proof}
    \begin{outline}
        \1 En utilisant le théorème de Bézout, on trouve $U,V \in \Kset[X]$ tels que $UP_1 + VP_2 = 1$.
        \1 On passe au polynôme d'endomorphisme pour trouver $(UP_1)(f) + (VP_2)(f) = \Id_E$.
        \1 Cela nous donne :
        $$\begin{cases}
        U(f)\circ P_1(f) + V(f)\circ P_2(f) = \Id_E\\
        P_1(f)\circ U(f) + P_2(f)\circ V(f) = \Id_E
        \end{cases}$$
        \1 Soit $x \in \Ker(P_1(f)) \cap \Ker(P_2(f))$.
        \1 Donc $x \in E$ et $\begin{cases} P_1(f)(x) = 0_E\\ P_2(f)(x) = 0_E\end{cases}$.
        \1 Donc 
    \end{outline}
\end{proof}

\begin{property}
    Soient $E$ un $\Kset$-ev. et $H$ un sev. de $E$.

    $H$ est un hyperplan de $E$ $\iff$ $\exists \varphi \in E^* \setminus \{0_{E^*}\}$ telle que $H = \Ker(\varphi)$.
    
    i.e. : Tout hyperplan $H$ de $E$ est le noyau d'une forme linéaire non nulle de $E$.
\end{property}

\begin{proof}
    recop

    % si x in e tq x in ker phi, alors phi x = 0
    % donc ** implique x = h_x + 0_E in H
    % c/c ker phi = H
\end{proof}


\vspace{1em}


\section*{Extraits : 260422 CH14-C}


\begin{proof}
    Pour montrer que $\mathcal M_{n,p}(\Kset)$ est un $\Kset$-ev., il faut montrer que :
    \begin{outline}
        \1 $(\mathcal M_{n,p}(\Kset),+)$ est un groupe abélien.
        \2 ... et donc que $+$ est bien une LCI\dots
        \2 qu'elle est associative et commutative...
        \2 que $0_{n,p} = (0)_{\substack{1 \le i \le n \\ 1 \le j \le p}}$ est élément neutre pour la loi + dans $\mathcal M_{n,p}(\Kset)$.
        \2 Enfin, que les éléments de $\mathcal M_{n,p}(\Kset)$ sont symétrisables par loi +.
        \1 Ensuite, $\cdot$ est bien une LCE et qu'elles vérifient les propriétés d'une LCE pour un espace vectoriel comme énoncées dans le chapitre 14-A en page 3.
    \end{outline}
\end{proof}

\begin{theorem}[Isomorphisme entre $\mathcal L(E,F)$ et $\mathcal M_{n,p}(\Kset)$]
    Soient $E,F$ deux $\Kset$-ev. tels que $\begin{cases}
    \underset\Kset\dim E = p\\ \underset\Kset\dim F = n
    \end{cases}$

    Soient $\begin{cases}
    B = (e_1,\dots,e_p) \text{ base de } E\\
    C = (c_1,\dots,c_n) \text{ base de } F
    \end{cases}$

    Alors $\psi : \begin{matrix}
    \mathcal{L}(E,F) &\to &\mathcal M_{n,p}(\Kset)\\
    f &\mapsto & \Mat(f,B,C)
    \end{matrix}$ est un isomorphisme de $\Kset$-ev.
\end{theorem}

\begin{proof}
    Pour montrer que $\psi$ est un isomorphisme de $E$ vers $F$, il faut vérifier que :
    \begin{itemize}
        \item $\psi$ est bien définie.
        \item $\psi$ est une application linéaire.
        \item $\psi$ est injective.
        \item $\psi$ est surjective.
    \end{itemize}

    \textbf{$\psi$ est bien définie :}

    \begin{itemize}
        \item On a $\forall f \in \mathcal L(E,F)$.
        \item On a $\Mat(f,B,C) = \psi(f) \in \mathcal M_{n,p}(\Kset)$.
        \item Donc $\psi$ est définie.
    \end{itemize}

    \textbf{$\psi$ est une application linéaire :}

    \begin{outline}
    \1 On pose $N = \Mat(f,B,C) = (a_{i,j})_{\substack{1 \le i \le n\\ 1 \le j \le p}}$.
    \2 On a donc $N = \left(\begin{matrix} f(e_1) & \cdots & f(e_p)\\
    a_{1,1} & \cdots & a_{1,p} \\
    \vdots & & \vdots \\
    a_{n,1} & \cdots & a_{n,p}
    \end{matrix}\right)$.
    \1 On pose $M = \Mat(g,B,C) = (b_{i,j})_{\substack{1 \le i \le n\\ 1 \le j \le p}}$.
    \2 Donc $\forall j \in [\![1,p]\!], g(e_j) = \sum_{i=1}^{n}b_{i,j}c_i$.
    \1 Donc $N + \lambda M = (a_{i,j} + b_{i,j})_{\substack{1 \le i \le n\\ 1 \le j \le p}}$.
    \1 Or $\forall j \in [\![1,p]\!], (f+\lambda g)(e_j)$.
    \begin{align*}
        (f+ \lambda g)(e_j) &= f(e_j) + \lambda g(e_j)\\
        &= \sum_{i=1}^n a_{i,j} c_i + \lambda \sum_{i=1}^{n}b_{i,j}c_i\\
        &= \sum_{i=1}^{n}(a_{i,j}+\lambda b_{i,j})\cdot c_i
    \end{align*}
    \1 Cela implique que $\underbrace{\Mat(f+\lambda g,B,C)}_{=\psi(f+\lambda g)}= \underbrace{N}_{\psi(g)} + \lambda \underbrace{M}_{\psi(g)}$.
    \1 Donc $\psi \in \mathcal L(\mathcal L(E,F), \mathcal M_{n,p}(\Kset))$.
    \end{outline}

    \textbf{$\psi$ est injective :}

    \begin{outline}
    \1 Soit $f \in \Ker(\psi)$. Donc $\psi(f) = 0_{n,p}$.
    \2 i.e. : (matrices en latex ehhh)
    \1 Donc $\forall 1 \le j \le p, f(e_j) = 0_F$.
    \1 Or $f(E) = \Vect(f(B)) = \{0_F\}$.
    \1 Donc $f = 0_{\mathcal L(E,F)}$.
    \1 i.e. $\Ker\psi = 0_{\mathcal L(E,F)}$.
    \1 Et donc $\psi$ est injective !
    \end{outline}

    \textbf{$\psi$ est surjective :}

    \begin{outline}
        \1 Soit $A = (a_{i,j})_{\substack{1 \le i \le n\\ 1 \le j \le p}} \in \mathcal M_{n,p}(\Kset)$.
        \1 On pose $f : \begin{matrix} E &\to &F \\ e_j &\mapsto &\sum_{i=1}^n a_{i,j} c_i \end{matrix}$ pour tout $1 \le j \le p$.
        \1 $f$ est bien définie, car $B$ est une base de $E$.
        \1 $f$ est linéaire, car $C$ est une base de $F$.
        \1 Donc $f \in \mathcal L(E,F)$.
        \1 Or $\psi(f) = \Mat(f,B,C) = A$.
        \1 Donc $\psi$ est surjective !
    \end{outline}

    On a montré que $\psi$ est un isomorphisme de $\Kset$-ev. entre $\mathcal L(E,F)$ et $\mathcal M_{n,p}(\Kset)$.
\end{proof}

\begin{proof}
    En effet, si $\underset\Kset\dim(E) = p$ et $\underset\Kset\dim(F) = n$, alors $\mathcal L(E,F)$ est isomorphe à $\mathcal M_{n,p}(\Kset)$, qui est de dimension $n \cdot p$.

    $$\mathcal L(E,F) \sim \mathcal M_{n,p}(\Kset) \implies \underset\Kset\dim(\mathcal L(E,F)) = \underset\Kset\dim(\mathcal M_{n,p}(\Kset)) = n \cdot p$$
\end{proof}

\begin{property}
    Soit $E$ un $\Kset$-ev. de dimension $n \in \N^*$.
    $$E \sim \Kset^n$$
\end{property}

\begin{proof}
    Soit $C = (c_1,\dots,c_n)$ une base de $E$.

    On pose $B_C = (e_1,\dots,e_n)$ la base canonique de $\Kset^n$.

    On définit $\varphi : \begin{matrix} K^n &\to &E \\ (\lambda_1,\dots,\lambda_n) &\mapsto &\sum_{i=1}^n \lambda_i c_i \end{matrix}$.

    \begin{outline}
        \1 $\varphi$ est bien définie, car $C$ est une base de $E$.
        \1 $\varphi$ est linéaire, car $C$ est une base de $E$.
        \1 Donc $\varphi \in \mathcal L(\Kset^n,E)$.
        \1 On a $\varphi(B_C) = C$, qui est une base de $E$.
        \1 Donc $\varphi$ est un isomorphisme de $\Kset$-ev. entre $\Kset^n$ et $E$.
    \end{outline}
\end{proof}

\begin{theorem}[Lemme]
    Pour tout $A,B \in \mathcal M_{n,p}(\Kset)$ et pour tout $\lambda \in \Kset$, on a :
    $$^t(A+\lambda B) = \,{^t}A + \lambda \cdot ^tB$$
    \label{th:transposée}
\end{theorem}

\begin{proof}
    On pose $A = (a_{i,j})_{\substack{1 \le i \le n \\ 1 \le j \le p}}$ et $B = (b_{i,j})_{\substack{1 \le i \le n \\ 1 \le j \le p}}$.
    
    On pose aussi $^tA = (c_{i,j})_{\substack{1 \le i \le p \\ 1 \le j \le n}}$ et $^tB = (d_{i,j})_{\substack{1 \le i \le p \\ 1 \le j \le n}}$.
    
    On pose aussi $(A+\lambda B) = (e_{i,j})_{\substack{1 \le i \le n \\ 1 \le j \le p}}$ et $^t(A+\lambda B) = (f_{i,j})_{\substack{1 \le i \le p \\ 1 \le j \le n}}$.

    On a, $\forall 1 \le i \le p, \forall 1 \le j \le n$ :
    \begin{align*}
        f_{i,j} &= e_{j,i}\\
        &= a_{j,i} + \lambda b_{j,i}\\
        &= c_{i,j} + \lambda d_{i,j}
    \end{align*}

    Donc $^t(A+\lambda B) = \,{^t}A + \lambda \cdot ^tB$.
\end{proof}

\begin{property}
    \begin{enumerate}
        \item $T_{n,i}(\Kset)$ et $T_{n,s}(\Kset)$ sont des sous-espaces vectoriels isomorphes de $\mathcal M_n(\Kset)$.
        \item $S_n(\Kset)$ et $A_n(\Kset)$ sont des sous-espaces vectoriels supplémentaires de $\mathcal M_n(\Kset)$ avec $\underset\Kset\dim(S_n(\Kset)) = \frac{n(n+1)}{2}$ et $\underset\Kset\dim(A_n(\Kset)) = \frac{n(n-1)}{2}$.
        \item $D_n(\Kset)$ est un sous-espace vectoriel de $\mathcal M_n(\Kset)$ avec $\underset\Kset\dim(D_n(\Kset)) = n$.
    \end{enumerate}
\end{property}

\begin{proof}
    \begin{outline}
        \1 Première propriété :
            \2 $T_{n,s}(\Kset)$ et $T_{n,i}(\Kset)$ sont des sous-espaces vectoriels de $\mathcal M_n(\Kset)$ (trivial).
            \2 Pour montrer qu'ils sont isomorphes, on pose $T: \begin{matrix} T_{n,s}(\Kset) &\to &T_{n,i}(\Kset) \\ A &\mapsto &^tA \end{matrix}$.
            \3 $T$ est bien définie (trivial).
            \3 D'après le lemme \ref{th:transposée}, $T$ est linéaire.
            \3 Et $\forall B \in T_{n,i}(\Kset), \exists ! A \in T_{n,s}(\Kset), B = \,^tA$. Donc $T$ est bijective.
            \2 Donc $T$ est un isomorphisme de $\Kset$-ev. entre $T_{n,s}(\Kset)$ et $T_{n,i}(\Kset)$.
            \2 Ils ont donc la même dimension.
            \2 Pour déterminer leur dimension, on pose $A = (a_{i,j})_{1 \le i,j \le n} \in T_{n,s}(\Kset)$.
            \2 On sait que $\forall i > j, a_{i,j} = 0$.
            \2 Donc $A = \sum_{1 \le i,j \le n} a_{i,j} E_{i,j}(n) = \sum_{1 \le i \le j \le n} a_{i,j} E_{i,j}(n)$.
            \2 Donc $B = (E_{i,j}(n))_{1 \le i \le j \le n}$ est une base de $T_{n,s}(\Kset)$.
            \2 Donc $\underset\Kset\dim(T_{n,s}(\Kset)) = \card{B} = \frac{n(n+1)}{2}$ (il s'agit d'une somme de $1 + 2 + \cdots + n$).
            \2 Donc $\underset\Kset\dim(T_{n,i}(\Kset)) = \frac{n(n+1)}{2}$.
            \2 \textit{Note additionnelle :} Une base de $T_{n,i}(\Kset)$ est $B' = (E_{i,j}(n))_{1 \le j \le i \le n}$ (on permute les indices $i$ et $j$).
        \1 Deuxième propriété :
            \2 $S_n(\Kset)$ et $A_n(\Kset)$ sont des sous-espaces vectoriels de $\mathcal M_n(\Kset)$ (trivial).
            \2 Pour montrer que $S_n(\Kset)$ et $A_n(\Kset)$ sont supplémentaires dans $\mathcal M_n(\Kset)$, on vérifie que $S_n(\Kset) \cap A_n(\Kset) = \{0_{\mathcal M_n(\Kset)}\}$ et que $S_n(\Kset) + A_n(\Kset) = \mathcal M_n(\Kset)$.
            \3 Soit $A \in S_n(\Kset) \cap A_n(\Kset)$. Donc $^tA = A$ et $^tA = -A$.
            \3 Donc $A = -A$, donc $2A = 0_{\mathcal M_n(\Kset)}$.
            \3 Donc $A = 0_{\mathcal M_n(\Kset)}$
            \3 Donc $S_n(\Kset) \cap A_n(\Kset) = \{0_{\mathcal M_n(\Kset)}\}$.
            \3 Soit $A = (a_{i,j})_{1 \le i,j \le n} \in \mathcal M_n(\Kset)$.
            \3 On pose $B = \frac{1}{2}(A + \,^tA)$ et $C = \frac{1}{2}(A - \,^tA)$.
            \3 Donc $^tB = \frac{1}{2}(^tA + A) = B$ et $^tC = \frac{1}{2}(^tA - A) = -C$.
            \3 Donc $B \in S_n(\Kset)$ et $C \in A_n(\Kset)$.
            \3 Donc $A = B + C \in S_n(\Kset) + A_n(\Kset)$.
            \3 Donc $S_n(\Kset) + A_n(\Kset) = \mathcal M_n(\Kset)$. (Dans le cours, on a procédé par analyse-synthèse).
            \2 Donc $S_n(\Kset)$ et $A_n(\Kset)$ sont des sous-espaces vectoriels supplémentaires de $\mathcal M_n(\Kset)$.
            \2 Pour déterminer leur dimension, on pose $A = (a_{i,j})_{1 \le i,j \le n} \in S_n(\Kset)$.
            \2 On va construire une base de $S_n(\Kset)$ à partir de la base canonique de $\mathcal M_n(\Kset)$.
            \3 On sait que $\forall 1 \le i,j \le n, a_{i,j} = a_{j,i}$.
            \3 Donc $A = \sum_{1 \le i,j \le n} a_{i,j} E_{i,j} = \sum_{1 \le i \le j \le n} a_{i,j} E_{i,j} + \sum_{1 \le j < i \le n} a_{i,j} E_{i,j}$.
            \3 Or $\forall 1 \le j < i \le n, a_{i,j} = a_{j,i}$.
            \3 Donc $A = \sum_{i=1}^n a_{i,i} E_{i,i} + \sum_{1 \le i < j \le n} a_{i,j} (E_{i,j} + E_{j,i})$.
            \3 Donc $B = (E_{i,i}(n))_{1 \le i \le n} \cup (E_{i,j}(n) + E_{j,i}(n))_{1 \le i < j \le n}$ est une base de $S_n(\Kset)$.
            \2 Donc $\underset\Kset\dim(S_n(\Kset)) = \card{B} = n + \frac{n(n-1)}{2} = \frac{n(n+1)}{2}$.
            \2 Donc $\underset\Kset\dim(A_n(\Kset)) = \frac{n(n-1)}{2}$.
            \3 On peut aussi déterminer une base de $A_n(\Kset)$ à partir de la base de $S_n(\Kset)$.
            \3 En effet, on sait que $B = (E_{i,i}(n))_{1 \le i \le n} \cup (E_{i,j}(n) + E_{j,i}(n))_{1 \le i < j \le n}$ est une base de $S_n(\Kset)$.
            \3 Donc $B' = (E_{i,j}(n) - E_{j,i}(n))_{1 \le i < j \le n}$ est une base de $A_n(\Kset)$ (on permute les signes dans $\sum_{1 \le i < j \le n} a_{i,j} (E_{i,j} + E_{j,i})$ parce que $a_{i,j} = -a_{j,i}$ et on enlève les éléments de la forme $a_{i,i} E_{i,i}$ parce que $a_{i,i} = 0$).
        \1 Troisième propriété :
            \2 Méthode 1 : Réutiliser la première propriété et le fait que $D_n(\Kset) = T_{n,s}(\Kset) \cap T_{n,i}(\Kset)$.
            \2 Méthode 2 : $D_n(\Kset) = \Vect(E_{i,i}(n))_{1 \le i \le n}$, qui est de dimension $n$.
    \end{outline}
\end{proof}

\begin{property}
    Soient $E,F,G$ trois $\Kset$-ev. tels que :
    $$\underset\Kset\dim(E) = q, B = (e_1,\dots,e_q) \text{ base de } E$$
    $$\underset\Kset\dim(F) = p, C = (c_1,\dots,c_p) \text{ base de } F$$
    $$\underset\Kset\dim(G) = n, D = (d_1,\dots,d_n) \text{ base de } G$$

    Soient $f \in \mathcal L(E,F)$ et $g \in \mathcal L(F,G)$.

    On pose $\displaystyle\begin{cases} A_1 = \Mat(f,B,C) \in \mathcal M_{p,q}(\Kset) = (b_{i,j})_{\substack{1 \le i \le p \\ 1 \le j \le q}}\\
    A_2 = \Mat(g,C,D) \in \mathcal M_{n,p} = (a_{i,j})_{\substack{1 \le i \le n \\ 1 \le j \le p}}
    \end{cases}$

    Alors $\Mat(g \circ f, B, D) = (\alpha_{i,j})_{\substack{1 \le i \le n \\ 1 \le j \le q}}$ tels que $\alpha_{i,j} = \sum_{k=1}^p a_{i,k} b_{k,j}$.
\end{property}

\begin{proof}
\begin{itemize}
    \item On pose $A = \Mat(g \circ f, B, D) = (\alpha_{i,j})_{\substack{1 \le i \le n \\ 1 \le j \le q}}$.

    \item On a $\forall j \in [\![1,q]\!], (g \circ f)(e_j) = \sum_{i=1}^n \alpha_{i,j} d_i$.

    \item Or $\forall j \in [\![1,q]\!], f(e_j) = \sum_{k=1}^p b_{k,j} c_k$.

    \item Donc $\forall j \in [\![1,q]\!], (g \circ f)(e_j) = g(f(e_j)) = g\left(\sum_{k=1}^p b_{k,j} c_k\right) = \sum_{k=1}^p b_{k,j} g(c_k)$.

    \item Or $\forall k \in [\![1,p]\!], g(c_k) = \sum_{i=1}^n a_{i,k} d_i$.

    \item Donc $\forall j \in [\![1,q]\!], (g \circ f)(e_j) = \sum_{k=1}^p b_{k,j} g(c_k) = \sum_{k=1}^p b_{k,j} \left(\sum_{i=1}^n a_{i,k} d_i\right) = \sum_{i=1}^n \left(\sum_{k=1}^p a_{i,k} b_{k,j}\right) d_i$.

    \item Donc $\forall j \in [\![1,q]\!], (g \circ f)(e_j) = \sum_{i=1}^n \left(\sum_{k=1}^p a_{i,k} b_{k,j}\right) d_i$.

    \item Donc $A = (\alpha_{i,j})_{\substack{1 \le i \le n \\ 1 \le j \le q}}$ tels que $\alpha_{i,j} = \sum_{k=1}^p a_{i,k} b_{k,j}$.
\end{itemize}
\end{proof}

\begin{property}
    Soit $n \ge 2$. $\mathcal M_n(\Kset)$ muni de $+$ et $\times$ (produit matriciel) est un anneau non commutatif. L'élément neutre de $\times$ est $I_n$.
\end{property}

\begin{proof}
    \begin{outline}
        \1 On sait que $(\mathcal M_n(\Kset),+,\cdot)$ est un $\Kset$-ev. 
        \1 Donc $(\mathcal M_n(\Kset),+)$ est un groupe abélien.
        \1 Soient $A,B,C \in \mathcal M_n(\Kset)$ et $a,b,c$ les endomorphismes associés canoniquement à $A,B,C$ respectivement.
        $$[a]_{B_C} = A,~~[b]_{B_C} = B,~~[c]_{B_C} = C$$
        \1 On a $A \times (B \times C) = A \times [b\circ c]_{B_C} = [a \circ (b \circ c)]_{B_C}$.
        \1 Comme $\circ$ est associative, on a $[a \circ (b \circ c)]_{B_C} = [(a \circ b) \circ c]_{B_C} = [(a \circ b)]_{B_C} \times C = (A \times B) \times C$.
        \1 Donc $\times$ est associative.
        \1 On a $A\times(B+C) = [a]_{B_C} \times [b+c]_{B_C} = [a \circ (b+c)]_{B_C} = [a\circ b + a \circ c]_{B_C} = [a\circ b]_{B_C} + [a \circ c]_{B_C} = A\times B + A \times C$.
        \2 On peut séparer $[a \circ b + a \circ c]_{B_C}$ en $[a\circ b]_{B_C} + [a \circ c]_{B_C}$ en se référant à $\psi$, l'application considérée précédemment dans la section sur les applications linéaires entre espaces de matrices, qui est un isomorphisme de $\Kset$-ev. entre $\mathcal L(E,F)$ et $\mathcal M_{n,p}(\Kset)$.
        \1 Donc $\times$ est distributive par rapport à $+$.
        \1 Soit $A = (a_{i,j})_{1 \le i,j \le n} \in \mathcal M_n(\Kset)$. On pose :
        $$\left|~\begin{matrix}
        A\times I_n = (\alpha_{i,j})_{1 \le i,j \le n} \\
        I_n \times A = (\beta_{i,j})_{1 \le i,j \le n}
        \end{matrix}\right.$$
        \1 On a $\forall 1 \le i,j \le n, \alpha_{i,j} = \sum_{k=1}^n a_{i,k} \delta_{k,j} = a_{i,j}$.
        \1 Et $\forall 1 \le i,j \le n, \beta_{i,j} = \sum_{k=1}^n \delta_{i,k} a_{k,j} = a_{i,j}$.
        \1 Donc $A\times I_n = A$ et $I_n \times A = A$.
        \1 Donc $I_n$ est l'élément neutre de $\times$.
    \end{outline}
    À ce stade-là, on a montré que $(\mathcal M_n(\Kset),+,\times)$ est un anneau. Il reste à montrer que $\times$ n'est pas commutatif. On va procéder par contre-exemple, qui seront plus du domaine des propriétés que de la preuve.
\end{proof}

\begin{property}
    Soient $p,q,r,s \in [\![1,n]\!]$ des indices.

    On pose $A = E_{p,q}$ et $B = E_{r,s}$.\index{Élémentaire (Matrice)}.

    Pour expliciter, on aura :
    $$\begin{cases}
    A = E_{p,q} = (\underbrace{\delta_{i,p}\delta_{j,q}}_{a_{i,j}})_{1 \le i,j \le n}\\
    B = E_{r,s} = (\underbrace{\delta_{i,r}\delta_{j,s}}_{b_{i,j}})_{1 \le i,j \le n}
    \end{cases}$$

    On pose :
    
    $$\begin{cases}A\times B = (c_{i,j})_{1 \le i,j \le n}
    \\
    B\times A = (d_{i,j})_{1 \le i,j \le n}\end{cases}$$

    On a $\forall 1 \le i,j \le n$ :
    
    $$c_{i,j} = \sum_{k=1}^n a_{i,k} b_{k,j} = \sum_{k=1}^n \delta_{i,p}\delta_{k,q} \delta_{k,r}\delta_{j,s}$$
    
    Si $q \ne r$, alors $\forall 1 \le i,j \le n, c_{i,j} = 0$.

    Si $q = r$, alors $\forall 1 \le i,j \le n, c_{i,j} = \delta_{i,p}\delta_{j,s}$.

    De même, $\forall 1 \le i,j \le n$,
    
    $$d_{i,j} = \sum_{k=1}^n b_{i,k} a_{k,j} = \sum_{k=1}^n \delta_{i,r}\delta_{k,s} \delta_{k,p}\delta_{j,q}$$
    
    $$\implies \boxed{d_{i,j} = \delta_{i,r}\delta_{j,q}}$$

    Donc $A\times B = E_{p,s}$ si $q=r$ et $A\times B = 0_{\mathcal M_n(\Kset)}$ sinon.

    $$A\times B = \delta_{q,r} E_{p,s}$$

    De même, $B\times A = E_{r,q}$ si $s=p$ et $B\times A = 0_{\mathcal M_n(\Kset)}$ sinon.

    $$B\times A = \delta_{s,p} E_{r,q}$$

    Alors, si par exemple, $r \ne q$ et $p=s$, on a $A\times B = 0_{\mathcal M_n(\Kset)}$ et $B\times A = E_{r,q} \ne 0_{\mathcal M_n(\Kset)}$.

    Donc $A\times B \ne B\times A$.

    Cela montre que $\times$ n'est pas commutatif.
\end{property}

\begin{property}
    Soient $E$ un $\Kset$-ev. de dimension $n$ et $B$ une base de $E$. Alors :
    $$\psi : \begin{matrix} (\mathcal L(E),+,\circ,\cdot) &\to &(\mathcal M_n(\Kset),+, \times, \cdot) \\ f &\mapsto &[f]_B \end{matrix}$$
    est un isomorphisme de $\Kset$-algèbres.
\end{property}

\begin{property}
    Soient $E,F$ deux $\Kset$-ev. de dimensions respectives $p$ et $n$. Soient $B=(e_1,\ldots,e_p)$ une base de $E$ et $C=(c_1,\ldots,c_n)$ une base de $F$. Soit $f \in \mathcal L(E,F)$.

    Soit $x \in E$. On pose :
    $$\begin{cases}
    X = \Mat_B(x)\\
    A = \Mat(f,B,C)\\
    Y = \Mat_C(f(x))
    \end{cases}$$

    $$\text{Alors } Y = A \times X$$
\end{property}

\begin{proof}
    \begin{itemize}
        \item On a $\exists ! \lambda_1, \dots, \lambda_p \in \Kset$ tels que $x = \sum_{i=1}^p \lambda_i e_i$.
        \item Donc $X = \Mat_B(x) = \left(\begin{matrix} \lambda_1 \\ \vdots \\ \lambda_p \end{matrix}\right)$.
        \item On pose $A = \Mat(f,B,C) = (a_{i,j})_{\substack{1 \le i \le n \\ 1 \le j \le p}}$.
        \item Donc $\forall i \in [\![1,p]\!], f(e_i) = \sum_{j=1}^n a_{i,j} c_j$.
        \item Donc $f(x) = f\left(\sum_{i=1}^p \lambda_i e_i\right) = \sum_{i=1}^p \lambda_i f(e_i) = \sum_{i=1}^p \lambda_i \left(\sum_{j=1}^n a_{i,j} c_j\right) = \sum_{j=1}^n \left(\sum_{i=1}^p a_{i,j} \lambda_i\right) c_j$.
        \item Donc $Y = \Mat_C(f(x)) = \left(\begin{matrix} \sum_{i=1}^p a_{1,j} \lambda_i \\ \vdots \\ \sum_{i=1}^p a_{n,j} \lambda_i \end{matrix}\right)$.
        \item Or $A \times X = \left(\begin{matrix} a_{1,1} & \cdots & a_{1,p} \\ \vdots & & \vdots \\ a_{n,1} & \cdots & a_{n,p} \end{matrix}\right) \times \left(\begin{matrix} \lambda_1 \\ \vdots \\ \lambda_p \end{matrix}\right) = \left(\begin{matrix} \sum_{i=1}^p a_{1,i} \lambda_i \\ \vdots \\ \sum_{i=1}^p a_{n,i} \lambda_i \end{matrix}\right)$.
        \item Donc $Y = A \times X$.
    \end{itemize}
\end{proof}

\begin{property}
    Soit $A \in \mathcal M_n(\Kset)$. Les propriétés suivantes sont équivalentes :
    \begin{enumerate}
        \item $A \in \GL_n(\Kset)$ est inversible.
        \item L'application linéaire $f \in \mathcal L(\Kset^n)$ associée canoniquement à $A$ est un isomorphisme (application bijective).
        \item $\forall X \in \mathcal M_{n,1}(\Kset), AX = 0_{\mathcal M_{n,1}(\Kset)} \implies X = 0_{\mathcal M_{n,1}(\Kset)}$.
        \item $\forall Y \in \mathcal M_{n,1}(\Kset), \exists ! X \in \mathcal M_{n,1}(\Kset), AX = Y$.
        \item $\exists ! B \in \mathcal M_n(\Kset) \tqf AB = I_n$.
        \item $\exists ! B \in \mathcal M_n(\Kset) \tqf BA = I_n$.
    \end{enumerate}
\end{property}

\begin{proof}
    \textit{Seule l'équivalence entre les propriétés 1 et 3 sera démontrée. Les autres équivalences sont plus faciles (il suffit d'utiliser les propriétés de l'application linéaire associée canoniquement à $A$).}

    \begin{outline}
        \1 Montrons que $1 \implies 3$.
        \2 Soit $X \in \mathcal M_{n,1}(\Kset)$ tel que $AX = 0_{\mathcal M_{n,1}(\Kset)}$.
        \2 Cela implique que $A^{-1} \times (AX) = A^{-1} \times 0_{\mathcal M_{n,1}(\Kset)}$ (sachant que $A^{-1}$ existe parce que $A$ est inversible).
        \2 Donc $(A^{-1} \times A) \times X = 0_{\mathcal M_{n,1}(\Kset)}$.
        \2 Donc $I_n \times X = 0_{\mathcal M_{n,1}(\Kset)}$.
        \2 Donc $X = 0_{\mathcal M_{n,1}(\Kset)}$.
        \1 Montrons que $3 \implies 1$.
        \2 Soit $f \in \mathcal L(\Kset^n)$ l'application linéaire associée canoniquement à $A$.
        \2 On a $\forall X \in \mathcal M_{n,1}(\Kset), AX = 0_{\mathcal M_{n,1}(\Kset)} \implies X = 0_{\mathcal M_{n,1}(\Kset)}$.
        \2 Alors $\forall x \in \Kset^n, f(x) = 0_{\Kset^n} \implies x = 0_{\Kset^n}$.
        \2 Donc le noyau de $f$ est réduit à $\{0_{\Kset^n}\}$.
        \2 Donc $f$ est injective.
        \2 Or $\Kset^n$ est de dimension finie, donc $f$ est bijective.
        \2 Donc $f$ est un isomorphisme.
        \2 Donc $A$ est inversible, $\in \GL_n(\Kset)$.
    \end{outline}
\end{proof}

\begin{property}
    Soit $A = (a_{i,j})_{1 \le i,j \le n} \in T_{n,s}(\Kset)$ (matrice triangulaire supérieure). Alors $A$ est inversible si et seulement si $\forall i \in [\![1,n]\!], a_{i,i} \ne 0$.

    Dans ce cas, $\forall i \in [\![1,n]\!], [A^{-1}]_{i,i} = \frac{1}{a_{i,i}}$.
\end{property}

\begin{proof}
    \begin{outline}
        \1[$\impliedby$] Dans le sens indirect, on suppose que $\forall i \in [\![1,n]\!], a_{i,i} \ne 0$.
        \2 Soit $X = \left(\begin{matrix} x_1 \\ \vdots \\ x_n \end{matrix}\right) \in \mathcal M_{n,1}(\Kset)$ tel que $AX = 0_{\mathcal M_{n,1}(\Kset)}$, i.e. :
        $$\left(\begin{matrix} a_{1,1} & \cdots & \cdots & a_{1,n} \\ 0 & a_{2,2} & \cdots & a_{2,n} \\ \vdots & \ddots & \ddots & \vdots \\ 0 & \cdots & 0 & a_{n,n} \end{matrix}\right) \times \left(\begin{matrix} x_1 \\ x_2 \\ \vdots \\ x_n \end{matrix}\right) = 0_{\mathcal M_{n,1}(\Kset)}$$
        \2 Donc $\begin{cases} a_{1,1} x_1 + a_{1,2} x_2 + \cdots + a_{1,n} x_n = 0 \\ a_{2,2} x_2 + \cdots + a_{2,n} x_n = 0 \\ \vdots \\ a_{n,n} x_n = 0 \end{cases}$.
        \2 Cela implique directement que $x_n = x_{n-1} = \cdots = x_1 = 0$.
        \2 Donc $X = 0_{\mathcal M_{n,1}(\Kset)}$.
        \2 Donc $A$ est inversible, $\in \GL_n(\Kset)$.
        \1[$\implies$] Dans le sens direct, on suppose que $A$ est inversible, $\in \GL_n(\Kset)$.
        \2 On suppose que $\exists i_0 \in [\![1,n]\!], a_{i_0,i_0} = 0$.
        \2 Soit $f \in \mathcal L(\Kset^n)$ associée canoniquement à $A$. i.e. :
        % faire apparaitre i_0 dans la matrice
        $$[f]_{B_C} = A = \left(\begin{matrix} a_{1,1} & a_{1,2} & \cdots & \cdots & a_{1,n} \\ 0 & a_{2,2} & \cdots & a_{2,n} \\ \vdots & \ddots & a_{i_0,i_0} & \cdots & a_{i_0,n} \\ \vdots & & \ddots & \ddots & \vdots \\ 0 & \cdots & 0 & 0 & a_{n,n} \end{matrix}\right)$$
        \2 On peut donc écrire $(f(e_1),\dots,f(e_{i_0})) \in \Vect(e_1,\dots,e_{i_0-1})$.
        \2 Or le cardinal de $(f(e_1),\dots,f(e_{i_0}))$ est égal à $i_0$ et le cardinal de $\Vect(e_1,\dots,e_{i_0-1})$ est égal à $i_0-1$.
        \2 Donc $(f(e_1),\dots,f(e_{i_0}))$ est une famille liée.
        \2 Donc toute surfamille de $(f(e_1),\dots,f(e_{i_0}))$ est liée, donc la base $f(e_1),\dots,f(e_n)$ de $\Kset^n$ est liée, ce qui est absurde.
        \2 Donc $\forall i \in [\![1,n]\!], a_{i,i} \ne 0$.
    \end{outline}
\end{proof}

\begin{property}
    Soient $A \in \mathcal M_{n,p}(\Kset)$ et $B \in \mathcal M_{p,q}(\Kset)$. Alors :
    $$\underbrace{^t(AB)}_{\in \mathcal M_{q,n}(\Kset)} = \underbrace{\underbrace{(^tB)}_{\in \mathcal M_{q,p}(\Kset)} \cdot \underbrace{(^tA)}_{\in \mathcal M_{p,n}(\Kset)}}_{\in \mathcal M_{q,n}(\Kset)}$$
\end{property}

\begin{proof}
.
\end{proof}

\begin{property}
    \begin{enumerate}
        \item $T_{n,s}(\Kset)$ et $T_{n,i}(\Kset)$ sont des sous-algèbres de $\mathcal M_n(\Kset)$. En particulier, $T_{n,s}(\Kset)$ et $T_{n,i}(\Kset)$ sont stables par produit matriciel.
        \item $\forall A \in T_{n,s}(\Kset)\cap \GL_n(\Kset), A^{-1} \in T_{n,s}(\Kset)$.
    \end{enumerate}
\end{property}

\begin{proof}
    \begin{enumerate}
        \item On considère le produit matriciel suivant dans $T_{n,s}(\Kset)$ :
        $$\left(\begin{matrix}\alpha_1 & * & \cdots & \cdots & * \\ 0 & \alpha_2 & * & \cdots & * \\ \vdots & \ddots & \ddots & \ddots & \vdots \\ 0 & \cdots & 0 & \alpha_{n-1} & * \\ 0 & \cdots & \cdots & 0 & \alpha_n\end{matrix}\right) \times \left(\begin{matrix}\beta_1 & * & \cdots & \cdots & * \\ 0 & \beta_2 & * & \cdots & * \\ \vdots & \ddots & \ddots & \ddots & \vdots \\ 0 & \cdots & 0 & \beta_{n-1} & * \\ 0 & \cdots & \cdots & 0 & \beta_n\end{matrix}\right) = \left(\begin{matrix}\alpha_1 \beta_1 & * & \cdots & \cdots & * \\ 0 & \alpha_2 \beta_2 & * & \cdots & * \\ \vdots & \ddots & \ddots & \ddots & \vdots \\ 0 & \cdots & 0 & \alpha_{n-1} \beta_{n-1} & * \\ 0 & \cdots & \cdots & 0 & \alpha_n \beta_n\end{matrix}\right)$$
        Donc $T_{n,s}(\Kset)$ est stable par produit matriciel.
        \item Soit $A \in T_{n,s}(\Kset)\cap \GL_n(\Kset)$.
        On pose $f : \begin{matrix} T_{n,s}(\Kset) &\to &T_{n,s}(\Kset) \\ M &\mapsto &AM \end{matrix}$.
        
        L'application est bien définie parce que $T_{n,s}(\Kset)$ est stable par produit matriciel.
        
        $f$ est linéaire parce que $f(M+N) = A(M+N) = AM + AN = f(M) + f(N)$ et $f(\lambda M) = A(\lambda M) = \lambda (AM) = \lambda f(M)$.
        
        $f$ est injective parce que $f(M) = 0_{T_{n,s}(\Kset)} \implies AM = 0_{T_{n,s}(\Kset)} \implies M = 0_{T_{n,s}(\Kset)}$ (sachant que $A$ est inversible). Alors $\Ker(f) = \{0_{T_{n,s}(\Kset)}\}$.
        
        Or $T_{n,s}(\Kset)$ est de dimension finie, donc $f$ est bijective.

        Donc $\forall Y \in T_{n,s}(\Kset), \exists ! M \in T_{n,s}(\Kset), f(M) = AM = Y$.
        
        Donc $\exists ! M \in T_{n,s}(\Kset), f(M) = AM = I_n$.
        
        Donc $A^{-1} = M \in T_{n,s}(\Kset)$.
    \end{enumerate}
\end{proof}

\begin{property}
    Soient $E$ un $\Kset$-ev. de dimension finie $n \in \N^*$ et $B_1, B_2$ deux bases de $E$.
    \begin{enumerate}
        \item $P_{B_1}^{B_2} \in \GL_n(\Kset)$ avec $(P_{B_1}^{B_2})^{-1} = P_{B_2}^{B_1}$.
        \item Soit $x \in E$. On pose $x_1 = \Mat_{B_1}(x)$ et $x_2 = \Mat_{B_2}(x)$. Alors $x_1 = P_{B_1}^{B_2} \times x_2$, i.e. $\Mat_{B_1}(x) = \mathrm{Pass}(B_1,B_2) \times \Mat_{B_2}(x)$.
    \end{enumerate}
\end{property}

\begin{proof}
    \begin{outline}
        \1 \textbf{Démonstration de la première propriété.}
        \2 On a $P_{B_1}^{B_2} = \Mat(\Id_E, B_1, B_2)$.
        \2 Or $\Id_E \in \GL(E)$, donc $P_{B_1}^{B_2} \in \GL_n(\Kset)$.
        \2 Et on a $(P_{B_1}^{B_2})^{-1} = \Mat(\Id_E, B_2, B_1) = P_{B_2}^{B_1}$.
        \1 \textbf{Démonstration de la deuxième propriété.}
        \2 On prend $\Id_E : \underset{\text{Base }B_1}{E} \to \underset{\text{Base }B_2}{E}$.
        \2 Alors $\forall x \in E, \Id_E(x) = x$. On a donc $\Mat(\Id_E, B_1, B_2) \times \Mat_{B_1}(x) = \Mat_{B_2}(\Id_E(x))$.
        \2 Donc $P_{B_1}^{B_2} \times x_1 = x_2$.
    \end{outline}
\end{proof}

\begin{property}
    Soient $E,F$ deux $\Kset$-ev. de dimensions respectives $p$ et $n$. Soient $B_1, B_2$ deux bases de $E$ et $C_1, C_2$ deux bases de $F$. Soit $f \in \mathcal L(E,F)$.

    Soient $B_1, B_2$ deux bases de $E$. Soient $C_1, C_2$ deux bases de $F$.

    On pose $\begin{cases} A_1 = \Mat(f,B_1,C_1) \\ A_2 = \Mat(f,B_2,C_2) \end{cases}$.

    Alors $A_2 = P_{C_2}^{C_1} \times A_1 \times P_{B_1}^{B_2}$.

    i.e. $\Mat(f,B_2,C_2) = \mathrm{Pass}(C_2,C_1) \times \Mat(f,B_1,C_1) \times \mathrm{Pass}(B_1,B_2)$.
\end{property}

\begin{proof}
    On considère l'application $g : \underset{B_2}{E} \overset{\Id_E}{\to} \underset{B_1}{E} \overset{f}{\to} \underset{C_1}{F} \overset{\Id_F}{\to} \underset{C_2}{F}$.

    Donc $f = g = \Id_F \circ f \circ \Id_E$.

    Donc $\Mat(g, B_2, C_2) = \Mat(\Id_F, C_1, C_2) \times \Mat(f, B_1, C_1) \times \Mat(\Id_E, B_2, B_1)$.

    On en déduit que $\Mat(f, B_2, C_2) = P_{C_2}^{C_1} \times A_1 \times P_{B_1}^{B_2}$.

    D'où $A_2 = P_{C_2}^{C_1} \times A_1 \times P_{B_1}^{B_2}$.
\end{proof}

\begin{proof}
    En effet, $[f]_{B_1} = \Mat(f, B_1, B_1) = P_{B_1}^{B_2} \times \Mat(f, B_2, B_2) \times P_{B_2}^{B_1} = P_{B_1}^{B_2} \times [f]_{B_2} \times P_{B_2}^{B_1}$.
\end{proof}

\begin{proof}
    \textbf{Démonstration de la deuxième remarque.}

    \begin{outline}
    \1 On a : $f^p = 0$ et $f^{p-1} \ne 0$.
    \1 Donc $\exists x_0 \in E, f^{p-1}(x_0) \ne 0$.
    \1 On montre que la famille $(x_0, f(x_0), \dots, f^{p-1}(x_0))$ est libre.
    \2 Soit $\lambda_0,\dots,\lambda_{p-1} \in \Kset$ tels que $\lambda_0 x_0 + \lambda_1 f(x_0) + \cdots + \lambda_{p-1} f^{p-1}(x_0) = 0_E$.
    \2 Alors $f^{p-1}(\lambda_0 x_0 + \lambda_1 f(x_0) + \cdots + \lambda_{p-1} f^{p-1}(x_0)) = 0_E$.
    \2 Donc $\lambda_0 f^{p-1}(x_0) = 0_E$.
    \2 Donc $\lambda_0 = 0$.
    \2 Donc $\lambda_1 f^{p-1}(x_0) = 0_E$.
    \2 Donc $\lambda_1 = 0$.
    \2 En continuant ainsi, on trouve que $\lambda_0 = \lambda_1 = \cdots = \lambda_{p-1} = 0$.
    \2 Donc la famille $(x_0, f(x_0), \dots, f^{p-1}(x_0))$ est libre. \textit{On peut aussi procéder par absurde en supposant que la famille est liée, et en montrant que cela contredit le fait que $f^{p-1}(x_0) \ne 0$.}
    \2 Donc $p \le n$, car le cardinal de la famille $(x_0, f(x_0), \dots, f^{p-1}(x_0))$ est égal à $p$ et le cardinal de toute famille libre de $E$ est inférieur ou égal à $\dim E = n$.
    \end{outline}
\end{proof}

\begin{property}[Transitivité des matrices de passage]
    $$\mathrm{Pass}(B_1,B_3) = \mathrm{Pass}(B_1,B_2) \times \mathrm{Pass}(B_2,B_3)$$
\end{property}

\begin{proof}
    $\Id_E : E \to E \to E$ avec $B_1$ comme base de départ, $B_2$ comme base d'arrivée pour la première application, et $B_3$ comme base d'arrivée pour la seconde application.

    Donc $\Id_E = \Id_E \circ \Id_E$.

    Donc $\Mat(\Id_E, B_1, B_3) = \Mat(\Id_E, B_1, B_2) \times \Mat(\Id_E, B_2, B_3)$.

    D'où $\mathrm{Pass}(B_1,B_3) = \mathrm{Pass}(B_1,B_2) \times \mathrm{Pass}(B_2,B_3)$.
\end{proof}

\begin{property}
    ``éq'' est une relation d'équivalence dans $\mathcal M_{n,p}(\Kset)$.
\end{property}

\begin{proof}
    \begin{outline}
        \1 \textbf{Réflexivité.} Soit $A \in \mathcal M_{n,p}(\Kset)$. On a $A = I_n A I_p$, donc $A \text{ éq } A$.
        \1 \textbf{Symétrie.} Soient $A,B \in \mathcal M_{n,p}(\Kset)$ tels que $A \text{ éq } B$. Alors $\exists P \in \GL_n(\Kset), Q \in \GL_p(\Kset)$ tels que $B = PAQ$. Donc $A = P^{-1} B Q^{-1}$, donc $B \text{ éq } A$.
        \1 \textbf{Transitivité.} Soient $A,B,C \in \mathcal M_{n,p}(\Kset)$ tels que $A \text{ éq } B$ et $B \text{ éq } C$. Alors $\exists P_1,P_2 \in \GL_n(\Kset), Q_1,Q_2 \in \GL_p(\Kset)$ tels que $B = P_1 A Q_1$ et $C = P_2 B Q_2$. Donc $C = P_2 P_1 A Q_1 Q_2$, donc $A \text{ éq } C$.
    \end{outline}
\end{proof}

\begin{property}
    Soient $A_1,A_2 \in \mathcal M_{n,p}(\Kset)$.

    $A_1$ éq $A_2 \iff A_1 \et A_2$ représentent la même application linéaire dans 2 couples de bases différents

    $\iff \exists f \in \mathcal L(E,F), \exists B_1,B_2 \text{ bases de } E, \exists C_1,C_2 \text{ bases de } F, A_2 = \Mat(f,B_2,C_2) \text{ et } A_1 = \Mat(f,B_1,C_1)$
\end{property}

\begin{proof}
    
\end{proof}

\begin{property}
    ``$\sim$'' est une relation d'équivalence dans $\mathcal M_n(\Kset)$.
\end{property}

\begin{proof}
    Preuve analogue à celle de la relation ``éq''.
\end{proof}

\begin{property}
    Soient $A_1,A_2 \in \mathcal M_n(\Kset)$.

    $A_1 \sim A_2 \iff A_1$ et $A_2$ représentent le même endomorphisme dans 2 bases différentes.

    $$A_1 \sim A_2 \iff \exists f \in \mathcal L(E), \exists B_1,B_2 \text{ bases de } E, A_2 = [f]_{B_2} \text{ et } A_1 = [f]_{B_1}$$
\end{property}

\begin{proof}
    \begin{outline}
        \1 \textbf{Démonstration dans le sens indirect $\impliedby$ :}
        \2 Soit $f \in \mathcal L(E)$ et $B_1,B_2$ deux bases de $E$ telles que $A_2 = [f]_{B_2}$ et $A_1 = [f]_{B_1}$.
        \2 On utilise la conséquence \ref{cons:ch_base_endo} pour trouver que $[f]_{B_1} = P_{B_1}^{B_2} \times [f]_{B_2} \times P_{B_2}^{B_1}$.
        \2 Donc $A_1 = P_{B_1}^{B_2} \times A_2 \times P_{B_2}^{B_1}$.
        \2 Comme $P_{B_1}^{B_2} = (P_{B_2}^{B_1})^{-1}$, on trouve que $A_1 = P_{B_1}^{B_2} \times A_2 \times (P_{B_1}^{B_2})^{-1}$.
        \2 Donc $A_1 \sim A_2$.
        \1 \textbf{Démonstration dans le sens direct $\implies$ :}
        \2 Soit $P \in \GL_n(\Kset)$ telle que $A_1 = P A_2 P^{-1}$.
        \2 On sait qu'il existe un endomorphisme associé canoniquement à chacune des deux matrices $A_1$ et $A_2$.
        \2 On pose arbitrairement $f \in \mathcal L(\Kset^n)$ tel que $[f]_{B_C^{(n)}} = A_1$.
        \2 On va maintenant chercher une deuxième base $B$ de $\Kset^n$ telle que $[f]_B = A_2$.
        \2 Soit $B$ la base formée par les colonnes de $P^{-1}$.
        \2 Alors $P^{-1} = P_{B_C^{(n)}}^B$.
        \2 D'où $A_2 = P^{-1} A_1 P = P_{B_C^{(n)}}^B \times A_1 \times P_{B}^{B_C^{(n)}}$.
        \2 Donc $A_2 = P_{B_C^{(n)}}^B \times [f]_{B_C^{(n)}} \times P_{B}^{B_C^{(n)}}$.
        \2 Donc $A_2 = [f]_B$.
    \end{outline}
\end{proof}

\begin{property}
    Soient $E,F$ deux $\Kset$-ev. tels que $\dim E = p$ et $\dim F = n$, et $f \in \mathcal L(E,F)$ telle que $\rang(f) = r \ne 0$.

    $$\exists B \text{ base de } E, \exists C \text{ base de } F, \Mat(f,B,C) = \left(\begin{matrix} J_r & 0 \\ 0 & 0 \end{matrix}\right)_{(n,p)}$$

    avec $J_r$ la matrice qui contient une matrice identité de taille $r$ dans le coin supérieur gauche, et des zéros partout ailleurs.

    $J_r = \left(\begin{matrix} 1 & 0 & \cdots & 0 \\ 0 & 1 & \cdots & 0 \\ \vdots & \vdots & \ddots & \vdots \\ 0 & 0 & \cdots & 0 \end{matrix}\right)_{(n,p)}$ (à la ligne et colonne $r$, on retrouve le dernier 1 diagonal).
\end{property}

\begin{proof}
    \begin{itemize}
        \item On a $\Ker(f)$ est un sev. de $E$ de dimension $p-r$ (par théorème du rang).

        \item Soit $(e_{r+1}, \dots, e_p)$ une base de $\Ker(f)$.

        \item On va compléter cette famille en une base de $E$ de la forme $B = (e_1, \dots, e_r, e_{r+1}, \dots, e_p)$.

        \item Par théorème de la base incomplète, comme $(e_{r+1}, \dots, e_p)$ est une famille libre de $E$, il existe $e_1, \dots, e_r \in E$ tels que $B = (e_1, \dots, e_r, e_{r+1}, \dots, e_p)$ soit une base de $E$.
        
        \item Or $\Imf(f) = \Vect(f(B))$, mais les images de $f$ par $(e_{r+1}, \dots, e_p)$ sont nulles, donc $\Imf(f) = \Vect(f(e_1), \dots, f(e_r))$.
        \item Donc $(f(e_1), \dots, f(e_r))$ est une famille génératrice de $\Imf(f)$. Son cardinal est $r$.
        \item Or $\dim \Imf(f) = r$, donc $(f(e_1), \dots, f(e_r))$ est une base de $\Imf(f)$.
        \item On va compléter cette famille en une base de $F$ de la forme $C = (f(e_1), \dots, f(e_r), \varepsilon_{r+1}, \dots, \varepsilon_n)$.
        \item Par théorème de la base incomplète, comme $(f(e_1), \dots, f(e_r))$ est une famille libre de $F$, il existe $\varepsilon_{r+1}, \dots, \varepsilon_n \in F$ tels que $C = (f(e_1), \dots, f(e_r), \varepsilon_{r+1}, \dots, \varepsilon_n)$ soit une base de $F$.
        \item Alors $\Mat(f,B,C) = \left(\begin{matrix} I_r & 0 \\ 0 & 0 \end{matrix}\right)_{(n,p)} = J_r(n,p)$.
        
        $$\begin{matrix}
            f(e_1) \\
            \vdots \\
            f(e_r) \\
            \varepsilon_{r+1} \\
            \vdots \\
            \varepsilon_n
        \end{matrix}
        \overset{\scriptsize{f(e_1) ~ f(e_2) ~ \cdots ~ f(e_r) ~ f(e_{r+1}) ~ \cdots ~ f(e_p)}}{\left(\begin{matrix}
            1 & 0 & \cdots & 0 & 0 & \cdots & 0 \\
            0 & 1 & \cdots & 0 & 0 & \cdots & 0 \\
            \vdots & \vdots & \ddots & \vdots & \vdots & \ddots & \vdots \\
            0 & 0 & \cdots & 1 & 0 & \cdots & 0 \\
            0 & 0 & \cdots & 0 & 0 & \cdots & 0 \\
            \vdots & \vdots & \ddots & \vdots & \vdots & \ddots & \vdots \\
            0 & 0 & \cdots & 0 & 0 & \cdots & 0
        \end{matrix}\right)_{(n,p)}}$$
    \end{itemize}
\end{proof}

\begin{property}
    Soit $A \in \mathcal M_{n,p}(\Kset)$.
    
    Soit $P \in \GL_p(\Kset)$ et $Q \in \GL_n(\Kset)$.

    Alors $\begin{cases} \rang(AP) = \rang(A) \\ \rang(QA) = \rang(A) \end{cases}$
\end{property}

\begin{proof}
    Soient $h \in \mathcal L(\Kset^p, \Kset^n)$, $f \in \mathcal L(\Kset^p)$, $g \in \mathcal L(\Kset^n)$ associées canoniquement à $A$, $P$ et $Q$ respectivement.

    Donc $\rang(AP) = \rang(h \circ f) = \rang h = \rang A$ car $P$ est inversible donc $f$ est une bijection (isomorphisme de $\Kset^p$ dans lui-même).

    De même pour $\rang(QA) = \rang(g \circ h) = \rang h = \rang A$ car $Q$ est inversible donc $g$ est une bijection (isomorphisme de $\Kset^n$ dans lui-même).
\end{proof}

\begin{proof}
    \begin{outline}
        \1 Dans le sens direct, on utilise la propriété précédente pour trouver que $A \text{ éq } J_r(n,p)$.
        \1 Dans le sens indirect, on a $A = Q\,J_r(n,p)\,P$ avec $P \in \GL_p(\Kset)$ et $Q \in \GL_n(\Kset)$. Donc $\rang(A) = \rang(Q\,J_r(n,p)\,P) = \rang(J_r(n,p)) = r$ car $P$ et $Q$ sont inversibles.
    \end{outline}
\end{proof}

\begin{proof}
    Dans le sens direct, on utilise simplement la transitivité de la relation équivalence.

    Dans le sens indirect, on a $A \text{ éq } B$, donc $\exists P \in \GL_n(\Kset), Q \in \GL_p(\Kset)$ tels que $B = PAQ$. Donc $\rang(B) = \rang(PAQ) = \rang(A)$ car $P$ et $Q$ sont inversibles.
\end{proof}

\begin{property}
    Soit $A \in \mathcal M_{n,p}(\Kset)$.

    Alors $\rang(\,^tA) = \rang(A)$.
\end{property}

\begin{proof}
    Soit $r = \rang(A)$.

    Alors $\exists P \in \GL_n(\Kset), Q \in \GL_p(\Kset)$ tels que $A = P\,J_r(n,p)\,Q$.

    Donc $^tA = \,{^tQ}\,{^tJ_r(n,p)}\,{^tP}$.

    Avec $^tQ \in \GL_p(\Kset)$ et $^tP \in \GL_n(\Kset)$ et $\,^tJ_r(n,p) = J_r(p,n)$.

    Donc $\rang(\,^tA) = \rang(J_r(p,n)) = r$.
\end{proof}

\begin{property}
    $T_{j,k}(\lambda) \in \GL_p(\Kset)$ et $T_{j,k}(\lambda)^{-1} = T_{j,k}(-\lambda)$.
\end{property}

\begin{proof}
    On a $T_{j,k}(\lambda) \times T_{j,k}(-\lambda) = (I_p + \lambda E_{j,k})(I_p - \lambda E_{j,k}) = I_p - \lambda^2 E_{j,k}^2 = I_p$ car $E_{j,k}^2 = \delta_{j,k}E_{j,k} = 0$ car $j \ne k$.
\end{proof}

\begin{property}
    Soient $A \in \mathcal M_{n,p}(\Kset)$, $j,k \in \{1,\dots,p\}$ tels que $j \ne k$, et $\lambda \in \Kset^*$.

    Si $A' = A \times T_{k,j}(\lambda)$, alors $A'$ est la matrice obtenue à partir de $A$ en effectuant l'opération élémentaire $c_j \leftarrow c_j + \lambda c_k$.
\end{property}

\begin{proof}
    On suppose que $A' = A \times T_{k,j}(\lambda)$.

    On doit montrer que :

    $$\begin{cases}
        C_j(A') = C_j(A) + \lambda C_k(A) \\
        C_i(A') = C_i(A) \text{ pour tout } i \in \{1,\dots,p\} \text{ tel que } i \ne j
    \end{cases}$$

    On a $C_j(A') = A' E_j = A T_{k,j}(\lambda) E_j$.

    $\implies C_j(A') = A(I_p + \lambda E_{k,j}) E_j = A E_j + \lambda A (E_{k,j}\times E_j)$.

    Or $E_{k,j}\times E_j = \delta_{j,j} E_k = E_k$.

    Donc $C_j(A') = A E_j + \lambda A E_k = C_j(A) + \lambda C_k(A)$.

    De même, pour $i \ne j$, on trouve que $C_i(A') = A' E_i = A T_{k,j}(\lambda) E_i = A (I_p + \lambda E_{k,j}) E_i = A E_i + \lambda A (E_{k,j}\times E_i)$.

    Or $E_{k,j}\times E_i = \delta_{j,i} E_k = 0$ car $i \ne j$.

    Donc $C_i(A') = A E_i + \lambda A (E_{k,j}\times E_i) = A E_i = C_i(A)$.

    % $C_j(A') = A' E_j = A T_{j,k}(\lambda) E_j = A (I_p + \lambda E_{j,k}) E_j = A E_j + \lambda A (E_{j,k}\times E_j)$.
\end{proof}

\begin{property}
    $D_j(\alpha) \in \GL_p(\Kset)$ et $D_j(\alpha)^{-1} = D_j(\alpha^{-1})$.
\end{property}

\begin{property}
    Soient $A \in \mathcal M_{n,p}(\Kset)$, $j \in \{1,\dots,p\}$, et $\alpha \in \Kset^*$.

    Si $A' = A \times D_j(\alpha)$, alors $A'$ est la matrice obtenue à partir de $A$ en effectuant l'opération élémentaire $c_j \leftarrow \alpha c_j$.
\end{property}

\begin{proof}
    La preuve est analogue à celle de la propriété précédente, en utilisant le fait que $D_j(\alpha) E_j = \alpha E_j$ et $D_j(\alpha) E_i = E_i$ pour $i \ne j$.
\end{proof}

\begin{property}
    $\forall \sigma,\tau \in S_n, P_\sigma \times P_\tau = P_{\sigma \circ \tau} = P_{\sigma \tau}$.

    En particulier, $P_\sigma \in \GL_n(\Kset)$ et $(P_\sigma)^{-1} = P_{\sigma^{-1}} = \,^tP_\sigma$, sachant que $\sigma^{-1}$ existe parce que $\sigma$ est une bijection de $\{1,\dots,n\}$ dans lui-même.
\end{property}

\begin{proof}
    \textit{Mentionné oralement : On peut voir $P_\sigma$ comme une sorte de matrice échelonnée obtenue à partir de $I_n$ en effectuant des opérations élémentaires de permutation de colonnes, qui sont inversibles.}

    Soit $1 \le k \le n$.

    \begin{itemize}
        \item On a $C_k(P_{\sigma \tau}) = E_{\sigma \tau(k)}$.
        \item On a $C_k(P_\sigma \times P_\tau) = P_\sigma \times C_k(P_\tau) = P_\sigma \times E_{\tau(k)} = E_{\sigma(\tau(k))} = E_{\sigma \tau(k)}$.
        \item D'où $C_k(P_{\sigma \tau}) = C_k(P_\sigma \times P_\tau)$ pour tout $k \in \{1,\dots,n\}$, donc $P_{\sigma \tau} = P_\sigma \times P_\tau$.
        \item Et on a $P_{\sigma^{-1}} \times P_\sigma = P_{\sigma^{-1} \sigma} = P_{\Id_{S_n}} = I_n$, donc $P_{\sigma^{-1}} = (P_\sigma)^{-1}$.
        \item On pose $P_{\sigma^{-1}} = (\alpha_{i,j})_{1 \le i,j \le n}$.
        \item On a $\forall i,j \in \{1,\dots,n\}, \alpha_{i,j} = \delta_{i,\sigma^{-1}(j)}$.
        \item Et on a $\delta_{i,\sigma^{-1}(j)} = 1 \iff \sigma^{-1}(j) = i \iff j = \sigma(i) \iff \delta_{j,\sigma(i)} = 1$.
        \item Si on pose $P_\sigma = (\beta_{i,j})_{1 \le i,j \le n}$, on trouve que $\beta_{i,j} = \delta_{i,\sigma(j)}$.
        \item Donc $^tP_\sigma = (\delta_{j,\sigma(i)})_{1 \le i,j \le n} = P_{\sigma^{-1}}$.
    \end{itemize}
\end{proof}

\begin{property}
    Soient $A \in \mathcal M_{n,p}(\Kset)$ et $1 \le i \ne j \le p$.

    Si $A' = A \times P_{\tau_{i,j}}$, alors $A'$ se déduit de $A$ par l'opération $c_i \leftrightarrow c_j$.
\end{property}

\begin{proof}
    \begin{itemize}
        \item $C_j(A') = A' E_j = A P_{\tau_{i,j}} E_j = A P_{\tau_{i,j}(j)} = A E_i = C_i(A)$.
        \item $C_i(A') = A' E_i = A P_{\tau_{i,j}} E_i = A P_{\tau_{i,j}(i)} = A E_j = C_j(A)$.
        \item Pour $k \ne i,j$, $C_k(A') = A' E_k = A P_{\tau_{i,j}} E_k = A E_k = C_k(A)$.
    \end{itemize}
\end{proof}

\begin{property}
    Soient $A \in \mathcal M_{n,p}(\Kset)$ et $L_1, \dots, L_n$ les lignes de $A$.
    
    \begin{itemize}
        \item $L_j \leftarrow L_j + \lambda L_k$ correspond à $A' = T_{j,k}(\lambda) \times A$.

        \item $L_j \leftarrow \alpha L_j$ correspond à $A' = D_j(\alpha) \times A$.

        \item $L_j \leftrightarrow L_k$ correspond à $A' = P_{\tau_{j,k}} \times A$.
    \end{itemize}
\end{property}

\begin{proof}
    Les lignes de $A$ sont les colonnes de $^tA$.

    \begin{enumerate}
        \item $L_i(A') = C_i(\,^tA') = C_i(\,^t(T_{j,k}(\lambda) \times A)) = C_i(\,^tA \times \,^tT_{j,k}(\lambda))$.\\Et on a $T_{i,j}(\lambda) = I_n + \lambda E_{i,j}$, donc $^tT_{i,j}(\lambda) = I_n + \lambda E_{j,i} = T_{j,i}(\lambda)$.\\Donc $L_i(A') = C_i(\,^tA \times T_{j,i}(\lambda)) = C_i(\,^tA) + \lambda C_j(\,^tA) = L_i(A) + \lambda L_j(A)$.
        \item Preuve similaire pour $L_j \leftarrow \alpha L_j$.
        \item Preuve similaire pour $L_j \leftrightarrow L_k$.
    \end{enumerate}
\end{proof}

\begin{property}
    Soient $A = (a_{i,j})_{\substack{1 \le i \le n \\ 1 \le j \le p}} \in \mathcal M_{n,p}(\Kset)$ et $B = (b_{i,j})_{\substack{1 \le i \le p \\ 1 \le j \le n}} \in \mathcal M_{p,n}(\Kset)$.
    $$\Tr(AB) = \Tr(BA)$$
\end{property}

\begin{proof}
    On pose :
    $$\begin{cases}
    AB = (c_{i,j})_{\substack{1 \le i \le n \\ 1 \le j \le n}} \in \mathcal M_n(\Kset) \\
    BA = (d_{i,j})_{\substack{1 \le i \le p \\ 1 \le j \le p}} \in \mathcal M_p(\Kset)
    \end{cases}$$

    On a $\forall i,j \in \{1,\dots,n\}, c_{i,j} = \sum_{k=1}^p a_{i,k} b_{k,j}$.

    Et on a $\forall i,j \in \{1,\dots,p\}, d_{i,j} = \sum_{k=1}^n b_{i,k} a_{k,j}$.
    
    Donc $\Tr(AB) = \sum_{i=1}^n c_{i,i} = \sum_{i=1}^n \sum_{k=1}^p a_{i,k} b_{k,i}$.

    Or $\sum_{k=1}^p \sum_{i=1}^n a_{i,k} b_{k,i} = \sum_{k=1}^p d_{k,k} = \Tr(BA)$.

    Donc $\Tr(AB) = \Tr(BA)$.
\end{proof}

\begin{proof}
    Pour toutes bases $B$ et $B'$ de $E$, on a :
    $$[f]_B \sim [f]_{B'}$$
\end{proof}

\begin{property}
    \begin{enumerate}
        \item $\Tr : \mathcal M_n(\Kset) \to \Kset$ est une forme linéaire, donc $\Tr \in (\mathcal M_n(\Kset))^*$.
        \item Si $\underset\Kset\dim E < +\infty$, alors $\Tr : \mathcal L(E) \to \Kset$ est une forme linéaire, donc $\Tr \in (\mathcal L(E))^*$.
    \end{enumerate}
\end{property}

\begin{property}
    Soient $E$ un $\Kset$-espace vectoriel de dimension finie $n$, et $f$ un projecteur de $E$.

    Alors $\Tr(f) = \rang(f)$.
\end{property}

\begin{proof}
    On pose $\rang(f) = r$.

    $f$ projecteur de $E$ $\implies$ $E = \Ker(f) \oplus \Imf(f)$.

    Soit $B' = (e_1,\dots,e_r)$ une base de $\Imf(f)$, et soit $B'' = (e_{r+1},\dots,e_n)$ une base de $\Ker(f)$.

    Une base adaptée à la somme directe $E = \Ker(f) \oplus \Imf(f)$ est $B = (e_1,\dots,e_r,e_{r+1},\dots,e_n)$.

    Or $f$ est une projection sur $\Imf(f)$ parallèlement à $\Ker(f)$, donc $\forall i \in [\![1,r]\!], f(e_i) = e_i$, et $\forall i \in [\! [r+1,n]\!], f(e_i) = 0$.

    $$\implies [f]_B = \left(\begin{matrix}
        1 & 0 & \cdots & 0 & 0 & \cdots & 0 \\
        0 & 1 & \cdots & 0 & 0 & \cdots & 0 \\
        \vdots & \vdots & \ddots & \vdots & \vdots & \ddots & \vdots \\
        0 & 0 & \cdots & 1 & 0 & \cdots & 0 \\
        0 & 0 & \cdots & 0 & 0 & \cdots & 0 \\
        \vdots & \vdots & \ddots & \vdots & \vdots & \ddots & \vdots \\
        0 & 0 & \cdots & 0 & 0 & \cdots & 0
    \end{matrix}\right)$$

    Donc $[f]_B = J_r(n)$.

    Donc $\Tr(f) = \Tr([f]_B) = r = \rang(f)$.
\end{proof}

\begin{property}
    Si $(L)$ est un système linéaire de Cramer, alors il possède une solution unique.

    $$S(L) = \{ X = A^{-1} Y \}$$

\end{property}

\begin{property}
    Si $A \notin \GL_n(\Kset)$ ou (plus généralement) $n\ne p$, alors $(L)$ admet une solution particulière $X_0$ ssi $Y \in \Imf(A)$ ssi $y \in \Imf(f)$.

    Dans ce cas, $S(L) = X_0 + S(H)$.

    C'est ce qu'on appelle un sous-espace affine de direction $S(H)$ et passant par $X_0$.
\end{property}

\begin{proof}
    Si $Y \notin \Imf(A)$, alors $AX=Y$ n'admet pas de solution, donc $S(L) = \emptyset$.

    Sinon, $\exists X_0 \in \mathcal M_{p,1}(\Kset)$ telle que $AX_0 = Y$.

    \begin{align*}
        x \in S(L) & \iff AX = Y = AX_0 \\
        & \iff A(X - X_0) = 0 \\
        & \iff X - X_0 \in \Ker(A) \\
        & \iff X - X_0 \in S(H) \scq \Ker(A) = S(H) \\
        & \iff X \in X_0 + S(H)
    \end{align*}
\end{proof}


\vspace{1em}

\end{document}
